% Produire des écrits utilitaires : le rapport — Français 1re Tech — Cours signé OnBuch+
% Programme officiel MINESEC. Compiler avec : tectonic N16-produire-ecrits-utilitaires-rapport.tex
\newcommand{\DOCMATIERE}{Français}
\newcommand{\DOCNIVEAU}{1re Tech}
\newcommand{\DOCMODULE}{Français – classe de Première technique}
\newcommand{\DOCLECON}{N16}
\newcommand{\DOCTITRE}{Produire des écrits utilitaires : le rapport}
% OnBuch+ — préambule « cours signés » ENRICHI (compilé avec tectonic / XeLaTeX)
% Système visuel « Pop » d'OnBuch+ : fond crème (jamais blanc), bordures encre
% épaisses, ombres « dures » décalées, titres Archivo Black en capitales.
% Version MATHÉMATIQUES : graphes (pgfplots/tikz), boîtes pédagogiques
% (définition, propriété, méthode, attention, le savais-tu, exercice résolu,
% exercice, TP, bilan), page de garde et signature OnBuch+ ancrée en bas.
%
% Macros attendues dans main.tex AVANT \input{preamble} :
%   \DOCMATIERE  \DOCNIVEAU  \DOCMODULE  \DOCLECON (numéro)  \DOCTITRE
\documentclass[11pt]{article}

\usepackage{fontspec}
\usepackage[french]{babel}
\usepackage[a4paper,margin=2cm,top=3.4cm,bottom=2.8cm,headheight=1.4cm,footskip=1.3cm]{geometry}
\usepackage{amsmath,amssymb,amsfonts}
\usepackage{mathtools} % \xmapsto, \coloneqq... souvent utilisés par les modèles
\usepackage{cancel} % simplifications barrées (bilans de forces)
% \abs{z} (module, valeur absolue) et \norm{u} (norme), utilisés par les modèles
\providecommand{\abs}{}\let\abs\relax\DeclarePairedDelimiter{\abs}{\lvert}{\rvert}
\providecommand{\norm}{}\let\norm\relax\DeclarePairedDelimiter{\norm}{\lVert}{\rVert}
\usepackage[table]{xcolor} % [table] : fournit aussi \rowcolors (alternance de lignes)
\usepackage{pagecolor}
\usepackage{tikz}
\usetikzlibrary{arrows.meta,positioning,calc,shapes.geometric,shapes.misc,shapes.arrows,decorations.pathmorphing,decorations.pathreplacing,decorations.markings,patterns,shadows,mindmap,trees,fit,backgrounds,matrix,intersections,angles,quotes}
\usepackage{pgfplots}
\usepackage[version=4]{mhchem} % équations nucléaires \ce{^{235}_{92}U}
\usepackage[european,siunitx]{circuitikz} % schémas électriques
\pgfplotsset{compat=1.18}
\usepgfplotslibrary{fillbetween,groupplots} % aires entre courbes (calcul intégral)
\usepackage{enumitem}
\usepackage{fancyhdr}
% [most] charge toutes les bibliothèques tcolorbox (listings, minted, raster,
% documentation...) et gonfle inutilement la mémoire TeX sur un document long
% (~300 boîtes) : on ne charge que ce qui est réellement utilisé.
\usepackage[skins,breakable]{tcolorbox}
\usepackage{graphicx}
\usepackage{colortbl}
\usepackage{booktabs}
\usepackage{tabularx}
\usepackage{multirow}
\usepackage{diagbox} % case d'en-tête diagonale des tableaux à double entrée
% Colonnes extensibles centrée / gauche / droite (C, L, R), souvent utilisées par les modèles
\newcolumntype{C}{>{\centering\arraybackslash}X}
\newcolumntype{L}{>{\raggedright\arraybackslash}X}
\newcolumntype{R}{>{\raggedleft\arraybackslash}X}
\usepackage{array}
\usepackage{multicol}
\usepackage{siunitx}
% parse-numbers=false : accepte aussi \SI{2,5 \times 10^{-3}}{...}
\sisetup{locale=FR,output-decimal-marker={,},group-minimum-digits=4,parse-numbers=false}
\DeclareSIUnit\ohm{\ensuremath{\symup{\Omega}}} % U+2126 absent de la police
\DeclareSIUnit\division{div} % réglages d'oscilloscope (ms/div, V/div)
\DeclareSIUnit\tr{tr} % tours (vitesse de rotation, ex. \SI{1500}{\tr\per\minute})
\DeclareSIUnit\molar{M} % concentration molaire (ex. \SI{3}{\molar}, \SI{1}{\micro\molar})
\DeclareSIUnit\an{an} % année (le modèle confond parfois avec \year, un compteur TeX) — ex. \SI{5}{\kilo\meter\per\an}
\DeclareSIUnit\FCFA{FCFA} % franc CFA (ex. \SI{500}{\FCFA\per\kilo\gram})
\DeclareSIUnit\degreeBrix{\textdegree Brix} % teneur en sucre d'un sirop/jus (réfractométrie)
\DeclareSIUnit\month{mois} % le modèle utilise parfois \month, un compteur TeX, comme unité de durée
\DeclareSIUnit\microliter{\micro\liter} % le modèle écrit parfois \microliter en un seul token
\DeclareSIUnit\million{millions} % dénombrement (ex. \SI{15}{\million\per\mL} spermatozoïdes)
\usepackage{qrcode}
% \degree : commande fréquemment utilisée par les modèles pour un angle ou une
% température en dehors de \SI{}{} (non fournie par les paquets chargés).
\providecommand{\degree}{^{\circ}}
% \qed (fin de démonstration), utilisé par les modèles sans amsthm
\providecommand{\qed}{\hfill$\square$}
% \barre : double barre des valeurs interdites dans un tableau de variations (tkz-tab)
\providecommand{\barre}{\ensuremath{\|}}
% environnement proof (amsthm non chargé) utilisé par les modèles
\ifdefined\proof\else\newenvironment{proof}[1][Démonstration]{\par\noindent\textit{#1.}\ }{\qed\par}\fi
\usepackage{lastpage}
\usepackage{hyperref}
\hypersetup{hidelinks}

% ============================================================
% Palette « Pop » OnBuch+ — mêmes hex que la classe P (plus_ui.dart)
% ============================================================
\definecolor{poporange}{HTML}{F59321}
\definecolor{popdark}{HTML}{E07A0C}
\definecolor{popcream}{HTML}{FFF6E8}
\definecolor{popink}{HTML}{1C1714}
\definecolor{poppurple}{HTML}{7A5AE0}
\definecolor{popgreen}{HTML}{2E9E6A}
\definecolor{poppink}{HTML}{E0457B}
\definecolor{popblue}{HTML}{2F7DE1}
\definecolor{popgold}{HTML}{C98A12}
\definecolor{popmuted}{HTML}{8A7F76}
\definecolor{popline}{HTML}{EADFD2}
\definecolor{popgray}{HTML}{8A7F76}
\colorlet{popgrey}{popgray}
% Alias tolérants (le modèle nomme parfois les couleurs différemment)
\colorlet{popwhite}{white}
\colorlet{popblack}{popink}
\colorlet{popred}{poppink}
\colorlet{popyellow}{popgold}
% Teintes claires (fonds de tableaux/encadrés) — le modèle nomme parfois
% les couleurs avec un suffixe L (ex. popblueL) sans les avoir définies.
\colorlet{poporangeL}{poporange!20}
\colorlet{popdarkL}{popdark!20}
\colorlet{poppurpleL}{poppurple!20}
\colorlet{popgreenL}{popgreen!20}
\colorlet{poppinkL}{poppink!20}
\colorlet{popblueL}{popblue!20}
\colorlet{popgoldL}{popgold!20}
\colorlet{popredL}{popred!20}
\colorlet{popgrayL}{popgray!20}
% Teintes claires pour fonds de tableaux / zones
\colorlet{popblueL}{popblue!10!white}
\colorlet{poporangeL}{poporange!14!white}
\colorlet{popgreenL}{popgreen!12!white}
\colorlet{poppurpleL}{poppurple!12!white}
\colorlet{poppinkL}{poppink!12!white}
\colorlet{popgoldL}{popgold!14!white}

\pagecolor{popcream}

% ============================================================
% Typographie
% ============================================================
\defaultfontfeatures{Ligatures=TeX}
\setmainfont{PlusJakartaSans-Medium.ttf}[Path=../../../fonts/,BoldFont=PlusJakartaSans-Bold.ttf]
% Maths en sans-serif (Fira Math) avec chiffres et lettres droites de Plus
% Jakarta, pour que les formules se fondent dans le texte.
\usepackage[math-style=ISO,bold-style=ISO]{unicode-math}
\setmathfont{FiraMath-Regular.otf}[Path=../../../fonts/]
\setmathfont{PlusJakartaSans-Medium.ttf}[Path=../../../fonts/,range={up/num,up/{Latin,latin},\mathup}]
\usepackage{icomma} % virgule décimale sans espace en mode math
% Caractères mathématiques Unicode absents des polices de texte (Plus Jakarta,
% Archivo Black) : rendus par leur équivalent mathématique, en texte comme en
% formule — sinon ils s'affichent en carré vide (ex. « Arithmétique dans ℤ »).
\usepackage{newunicodechar}
\newunicodechar{ℝ}{\ensuremath{\mathbb{R}}}
\newunicodechar{ℤ}{\ensuremath{\mathbb{Z}}}
\newunicodechar{ℂ}{\ensuremath{\mathbb{C}}}
\newunicodechar{ℕ}{\ensuremath{\mathbb{N}}}
\newunicodechar{ℚ}{\ensuremath{\mathbb{Q}}}
\newunicodechar{𝔻}{\ensuremath{\mathbb{D}}}
\newunicodechar{α}{\ensuremath{\alpha}}
\newunicodechar{β}{\ensuremath{\beta}}
\newunicodechar{γ}{\ensuremath{\gamma}}
\newunicodechar{δ}{\ensuremath{\delta}}
\newunicodechar{ε}{\ensuremath{\varepsilon}}
\newunicodechar{θ}{\ensuremath{\theta}}
\newunicodechar{λ}{\ensuremath{\lambda}}
\newunicodechar{μ}{\ensuremath{\mu}}
\newunicodechar{σ}{\ensuremath{\sigma}}
\newunicodechar{τ}{\ensuremath{\tau}}
\newunicodechar{φ}{\ensuremath{\varphi}}
\newunicodechar{ψ}{\ensuremath{\psi}}
\newunicodechar{ω}{\ensuremath{\omega}}
\newunicodechar{Σ}{\ensuremath{\Sigma}}
% majuscules produites par \MakeUppercase dans les titres (α-aminés -> Α)
\newunicodechar{Α}{\ensuremath{\alpha}}
\newunicodechar{Β}{\ensuremath{\beta}}
\newunicodechar{Γ}{\ensuremath{\Gamma}}
\newunicodechar{Δ}{\ensuremath{\Delta}}
\newunicodechar{Ε}{\ensuremath{\varepsilon}}
\newunicodechar{Θ}{\ensuremath{\Theta}}
\newunicodechar{Λ}{\ensuremath{\Lambda}}
\newunicodechar{Μ}{\ensuremath{\mu}}
\newunicodechar{Τ}{\ensuremath{\tau}}
\newunicodechar{Φ}{\ensuremath{\Phi}}
\newunicodechar{Ψ}{\ensuremath{\Psi}}
\newunicodechar{Ω}{\ensuremath{\Omega}}
\newunicodechar{↦}{\ensuremath{\mapsto}}
\newunicodechar{∈}{\ensuremath{\in}}
\newunicodechar{∉}{\ensuremath{\notin}}
\newunicodechar{∧}{\ensuremath{\wedge}}
\newunicodechar{⇔}{\ensuremath{\Leftrightarrow}}
\newunicodechar{⟺}{\ensuremath{\Longleftrightarrow}}
\newunicodechar{⇒}{\ensuremath{\Rightarrow}}
\newunicodechar{⟹}{\ensuremath{\Longrightarrow}}
\newunicodechar{∘}{\ensuremath{\circ}}
\newunicodechar{∪}{\ensuremath{\cup}}
\newunicodechar{∩}{\ensuremath{\cap}}
\newunicodechar{≡}{\ensuremath{\equiv}}
\newunicodechar{⊂}{\ensuremath{\subset}}
\newunicodechar{⊥}{\ensuremath{\perp}}
\newunicodechar{∃}{\ensuremath{\exists}}
\newunicodechar{∀}{\ensuremath{\forall}}
\newunicodechar{∛}{\ensuremath{\sqrt[3]{\,}}}
\newunicodechar{∜}{\ensuremath{\sqrt[4]{\,}}}
\newunicodechar{⃗}{\ensuremath{{}^{\to}}}
\newunicodechar{ⁿ}{\ifmmode^{n}\else\textsuperscript{\ensuremath{n}}\fi}
\newunicodechar{ˣ}{\ifmmode^{x}\else\textsuperscript{\ensuremath{x}}\fi}
\newunicodechar{ᵇ}{\ifmmode^{b}\else\textsuperscript{\ensuremath{b}}\fi}
\newunicodechar{ʳ}{\ifmmode^{r}\else\textsuperscript{\ensuremath{r}}\fi}
\newunicodechar{ᵘ}{\ifmmode^{u}\else\textsuperscript{\ensuremath{u}}\fi}
\newunicodechar{ʸ}{\ifmmode^{y}\else\textsuperscript{\ensuremath{y}}\fi}
\newunicodechar{ᵃ}{\ifmmode^{a}\else\textsuperscript{\ensuremath{a}}\fi}
\newunicodechar{ᵉ}{\ifmmode^{e}\else\textsuperscript{\ensuremath{e}}\fi}
\newunicodechar{ᵏ}{\ifmmode^{k}\else\textsuperscript{\ensuremath{k}}\fi}
\newunicodechar{ᵖ}{\ifmmode^{p}\else\textsuperscript{\ensuremath{p}}\fi}
\newunicodechar{ᴺ}{\ifmmode^{N}\else\textsuperscript{\ensuremath{N}}\fi}
\newunicodechar{ᶜ}{\ifmmode^{c}\else\textsuperscript{\ensuremath{c}}\fi}
\newunicodechar{ʰ}{\ifmmode^{h}\else\textsuperscript{\ensuremath{h}}\fi}
\newunicodechar{ᵠ}{\ifmmode^{q}\else\textsuperscript{\ensuremath{q}}\fi}
\newunicodechar{ₙ}{\ifmmode_{n}\else\textsubscript{\ensuremath{n}}\fi}
\newunicodechar{ₐ}{\ifmmode_{a}\else\textsubscript{\ensuremath{a}}\fi}
\newunicodechar{ᵢ}{\ifmmode_{i}\else\textsubscript{\ensuremath{i}}\fi}
\newunicodechar{ᵣ}{\ifmmode_{r}\else\textsubscript{\ensuremath{r}}\fi}
\newunicodechar{ₖ}{\ifmmode_{k}\else\textsubscript{\ensuremath{k}}\fi}
\newunicodechar{ᵦ}{\ifmmode_{\beta}\else\textsubscript{\ensuremath{\beta}}\fi}
\newunicodechar{ₓ}{\ifmmode_{x}\else\textsubscript{\ensuremath{x}}\fi}
\newfontfamily\popdisplay{ArchivoBlack-Regular.ttf}[Path=../../../fonts/]
\newfontfamily\poplogo{SpaceGrotesk-Bold.ttf}[Path=../../../fonts/]
\newfontfamily\popsemi{PlusJakartaSans-SemiBold.ttf}[Path=../../../fonts/]
\newfontfamily\popxbold{PlusJakartaSans-ExtraBold.ttf}[Path=../../../fonts/]
\newfontfamily\popxboldls{PlusJakartaSans-ExtraBold.ttf}[Path=../../../fonts/,LetterSpace=3.0]

\setlist[enumerate]{leftmargin=1.7em,itemsep=5pt,topsep=4pt}
\setlist[itemize]{leftmargin=1.7em,itemsep=4pt,topsep=4pt,label={\color{poporange}\textbullet}}
\setlength{\parindent}{0pt}
\setlength{\parskip}{5pt}
\renewcommand{\arraystretch}{1.25}

% pgfplots : style Pop par défaut
\pgfplotsset{
  every axis/.append style={axis line style={popink,line width=1pt},tick style={popink},
    grid style={popline,line width=0.5pt},label style={font=\small},tick label style={font=\footnotesize}},
  popcurve/.style={poporange,line width=1.8pt,no markers,smooth},
}
% Même style disponible dans un \draw TikZ simple (hors pgfplots)
\tikzset{popcurve/.style={poporange,line width=1.8pt}}
\tikzset{popgrid/.style={popline,very thin}}
\colorlet{popcurve}{poporange} % le nom du style est parfois utilisé comme couleur

% ============================================================
% Pill / squiggle
% ============================================================
\newcommand{\poppill}[3]{%
  \tikz[baseline=-0.55ex]\node[draw=popink,line width=1.2pt,rounded corners=2.6pt,%
    fill=#2,inner xsep=6.5pt,inner ysep=3pt]{%
    \popxboldls\fontsize{7}{7}\selectfont\color{#3}\MakeUppercase{#1}};%
}
\newcommand{\popsquiggle}[1]{%
  \tikz[baseline=-0.4ex]{
    \draw[#1,line width=2.6pt,line cap=round]
      plot [smooth] coordinates {(0,0.09) (0.34,0.01) (0.68,0.21) (1.02,0.01) (1.36,0.10)};
  }%
}
% Mot-clé surligné (marqueur orange sous le texte)
\newcommand{\cle}[1]{{\popxbold\color{popdark}#1}}

% ============================================================
% En-tête / pied
% ============================================================
\pagestyle{fancy}
\fancyhf{}
\renewcommand{\headrulewidth}{0pt}
\renewcommand{\footrulewidth}{0pt}
\lhead{\raisebox{-0.15ex}{\poplogo\fontsize{13}{13}\selectfont\color{popink}OnBuch\color{poporange}+}}
\rhead{\poppill{\DOCMATIERE\ \textbullet\ \DOCNIVEAU\ \textbullet\ Leçon \DOCLECON}{white}{popink}}
\lfoot{\popxbold\fontsize{7.6}{7.6}\selectfont\color{popmuted}\MakeUppercase{Cours signé OnBuch+}\ \textbullet\ {\popsemi\DOCTITRE}}
\rfoot{\tikz[baseline=-0.6ex]\node[rounded corners=8.5pt,fill=popink,minimum height=17pt,minimum width=17pt,inner xsep=5pt,inner sep=0]{\popxbold\fontsize{8}{8}\selectfont\color{popcream}\thepage\,/\,\pageref*{LastPage}};}

% ============================================================
% Titres
% ============================================================
\newcommand{\chaptitre}[2]{%
  \begin{tcolorbox}[enhanced,colback=poporange,colframe=popink,boxrule=2.6pt,arc=11pt,
      drop shadow=popink,left=15pt,right=15pt,top=12pt,bottom=11pt,before skip=3pt,after skip=18pt]
    \poppill{#2}{popink}{popcream}\par\vspace{9pt}
    {\raggedright\hyphenpenalty=10000\exhyphenpenalty=10000\popdisplay\fontsize{22}{24}\selectfont\color{popcream}\MakeUppercase{#1}\par}
  \end{tcolorbox}
}
% Section numérotée : pastille orange avec le numéro + titre Archivo
\newcounter{popsec}
\newcommand{\coursec}[1]{%
  \stepcounter{popsec}\setcounter{popsub}{0}%
  \par\vspace{16pt}\noindent
  \begin{minipage}{\linewidth}
  \tikz[baseline=-0.7ex]\node[draw=popink,line width=1.6pt,fill=poporange,rounded corners=4pt,minimum size=22pt,inner sep=0]{\popdisplay\fontsize{12}{12}\selectfont\color{popcream}\thepopsec};%
  \hspace{8pt}{\popdisplay\fontsize{15}{17}\selectfont\color{popink}\MakeUppercase{#1}}\par
  \vspace{4pt}\noindent{\color{poporange}\rule{36pt}{3.6pt}}
  \end{minipage}\par\nopagebreak\vspace{8pt}
}
\newcounter{popsub}
\newcommand{\courssub}[1]{%
  \stepcounter{popsub}%
  \par\vspace{9pt}\noindent{\popxbold\fontsize{11.5}{13}\selectfont\color{popink}{\color{poporange}\thepopsec.\thepopsub}\hspace{6pt}#1}\par\nopagebreak\vspace{4pt}
}

% ============================================================
% Boîtes pédagogiques — même recette : fond blanc, bordure épaisse colorée,
% ombre dure encre, badge plein. Titre optionnel [#1].
% ============================================================
\tcbset{popbase/.style={breakable,enhanced,colback=white,boxrule=2.3pt,arc=9pt,
  shadow={3pt}{-3pt}{0pt}{popink},left=14pt,right=14pt,top=11pt,bottom=11pt,before skip=11pt,after skip=15pt,
  fontupper=\popsemi\fontsize{10.6}{14.5}\selectfont\color{popink},
  fontlower=\popsemi\fontsize{10.6}{14.5}\selectfont\color{popink}}}
% Coche dessinée (le glyphe ✓ n'existe pas dans les polices du document)
\newcommand{\popcheck}[1]{\tikz[baseline=-0.5ex]\draw[#1,line width=1.6pt,line cap=round,line join=round](-0.13,0.0)--(-0.04,-0.09)--(0.13,0.1);}
\providecommand{\coche}{\popcheck{popgreen}}
\providecommand{\resultat}[1]{\par\smallskip\noindent\fcolorbox{popgreen}{popgreenL}{\parbox{\dimexpr\linewidth-2\fboxsep-2\fboxrule\relax}{\textbf{Résultat.} #1}}\par\smallskip}
\newcommand{\popboxhead}[3]{\poppill{#1}{#2}{white}\ifx\relax#3\relax\else\hspace{6pt}{\popxbold\fontsize{10.6}{12}\selectfont\color{popink}#3}\fi\par\vspace{7pt}}

\newtcolorbox{aretenir}[1][]{popbase,colframe=popgreen,before upper={\popboxhead{À retenir}{popgreen}{#1}}}
\newtcolorbox{exemplebox}[1][]{popbase,colframe=poporange,before upper={\popboxhead{Exemple}{poporange}{#1}}}
\newtcolorbox{definition}[1][]{popbase,colframe=popblue,before upper={\popboxhead{Définition}{popblue}{#1}}}
\newtcolorbox{propriete}[1][]{popbase,colframe=poppurple,before upper={\popboxhead{Propriété}{poppurple}{#1}}}
\newtcolorbox{methode}[1][]{popbase,colframe=popgold,before upper={\popboxhead{Méthode}{popgold}{#1}}}
\newtcolorbox{attention}[1][]{popbase,colframe=poppink,before upper={\popboxhead{Attention}{poppink}{#1}}}
\newtcolorbox{experience}[1][]{popbase,colframe=popblue,colback=popblueL,before upper={\popboxhead{Expérience}{popblue}{#1}}}
% « Le savais-tu ? » : carte violette pleine (contexte, culture, Cameroun)
\newtcolorbox{savaistu}[1][]{popbase,colback=poppurple,colframe=popink,
  fontupper=\popsemi\fontsize{10.4}{14.2}\selectfont\color{white},
  before upper={\poppill{Le savais-tu\,?}{white}{poppurple}\ifx\relax#1\relax\else\hspace{6pt}{\popxbold\color{white}#1}\fi\par\vspace{7pt}}}
% Exercice résolu : énoncé en haut, \tcblower puis la solution sur fond crème
\newcounter{popexo}\newcounter{popexr}
\newtcolorbox{exoresolu}[1][]{popbase,colframe=poporange,colbacklower=popcream,
  segmentation style={popink,line width=1.2pt,dashed},
  before upper={\stepcounter{popexr}\popboxhead{Exercice résolu \thepopexr}{poporange}{#1}},
  before lower={\poppill{Solution}{popink}{popcream}\par\vspace{6pt}}}
% Exercice (non résolu en place) — la correction est dans la partie Corrigés
\newtcolorbox{exercice}[1][]{popbase,colframe=popink,
  before upper={\stepcounter{popexo}\popboxhead{Exercice \thepopexo}{popink}{#1}}}
% Corrigé
\newtcolorbox{corrige}[1][]{popbase,colframe=popgreen,colback=popgreenL,
  before upper={\popboxhead{Corrigé}{popgreen}{#1}}}
% Objectifs / prérequis / bilan
\newtcolorbox{objectifs}{popbase,colframe=popink,colback=poporangeL,
  before upper={\popboxhead{Objectifs de la leçon}{popink}{}}}
\newtcolorbox{prerequis}{popbase,colframe=popink,colback=popblueL,
  before upper={\popboxhead{Prérequis}{popblue}{}}}
\newtcolorbox{bilan}{popbase,colframe=popink,colback=white,boxrule=2.8pt,
  before upper={\popboxhead{Fiche bilan}{poporange}{L'essentiel à connaître pour l'examen}}}
% Figure encadrée avec légende
\newtcolorbox{popfigure}[1][]{enhanced,breakable=false,colback=white,colframe=popline,boxrule=1.4pt,arc=8pt,
  left=10pt,right=10pt,top=10pt,bottom=8pt,before skip=10pt,after skip=12pt,halign=center,
  lower separated=false,halign lower=center,
  fontlower=\popsemi\fontsize{9}{11.5}\selectfont\color{popmuted},
  #1}
\newcommand{\legende}[1]{\par\vspace{6pt}{\popsemi\fontsize{9}{11.5}\selectfont\color{popmuted}#1}}

% ============================================================
% Signature OnBuch+
% ============================================================
\newcommand{\poplogoinline}[1]{{\poplogo\fontsize{#1}{#1}\selectfont\color{popink}OnBuch\color{poporange}+}}
\newcommand{\signatureonbuch}{%
  \begin{tcolorbox}[enhanced,colback=popink,colframe=popink,boxrule=2.2pt,arc=10pt,
      drop shadow=poporange,left=14pt,right=14pt,top=10pt,bottom=10pt,before skip=0pt,after skip=0pt]
    \tikz[baseline=-0.7ex]\node[circle,fill=popgreen,draw=popcream,line width=1.3pt,minimum size=22pt,inner sep=0]{\popcheck{white}};%
    \hspace{9pt}\begin{minipage}[c]{0.62\linewidth}
      {\popxbold\fontsize{12}{13}\selectfont\color{popcream}Signé\ \textbullet\ {\poplogo OnBuch\color{poporange}+}}\par\vspace{2pt}
      {\raggedright\popsemi\fontsize{8}{10.5}\selectfont\color{popline}\MakeUppercase{Équipe pédagogique OnBuch+}\\\MakeUppercase{Conforme au programme officiel MINESEC}\par}
    \end{minipage}\hfill
    {\poplogo\fontsize{22}{22}\selectfont\color{popcream}OnBuch\color{poporange}+}
  \end{tcolorbox}%
}

% ============================================================
% Page de garde
% #1 titre · #2 matière · #3 niveau · #4 module · #5 n° leçon · #6 durée ·
% #7 description / objectifs (texte ou itemize)
% ============================================================
\newcommand{\pagegarde}[7]{%
  \thispagestyle{empty}
  \newgeometry{a4paper,margin=1.7cm,top=1.6cm,bottom=1.4cm}%
  \begin{tikzpicture}[remember picture,overlay]
    \fill[poporange!18] ([xshift=-3.2cm,yshift=-3.6cm]current page.north east) circle (4.6cm);
    \fill[poppurple!14] ([xshift=2.4cm,yshift=7.6cm]current page.south west) circle (3.8cm);
  \end{tikzpicture}%
  {\poplogo\fontsize{30}{30}\selectfont\color{popink}OnBuch\color{poporange}+}\hfill
  \raisebox{0.6ex}{\poppill{Cours signé \textbullet\ Premium}{popink}{popcream}}\par
  \vspace{1pt}\popsquiggle{poppurple}\par
  \vspace{8pt}
  \begin{tcolorbox}[enhanced,colback=poporange,colframe=popink,boxrule=2.8pt,arc=13pt,
      drop shadow=popink,left=18pt,right=18pt,top=12pt,bottom=13pt,after skip=10pt]
    \poppill{#2 \textbullet\ #3}{popink}{popcream}\hspace{5pt}\poppill{#4}{white}{popink}\par\vspace{8pt}
    {\popdisplay\fontsize{14}{14}\selectfont\color{popink}LEÇON #5}\par\vspace{4pt}
    {\raggedright\hyphenpenalty=10000\exhyphenpenalty=10000\popdisplay\fontsize{25}{27}\selectfont\color{popcream}\MakeUppercase{#1}\par}
    \vspace{8pt}
    {\popxbold\fontsize{9.5}{9.5}\selectfont\color{popink}\MakeUppercase{Durée conseillée : #6}}
  \end{tcolorbox}
  \begin{tcolorbox}[enhanced,colback=white,colframe=popink,boxrule=2.2pt,arc=10pt,
      drop shadow=popink,left=16pt,right=16pt,top=10pt,bottom=10pt,after skip=8pt]
    \poppill{Dans cette leçon}{poppurple}{white}\par\vspace{5pt}
    \popsemi\fontsize{9.3}{12.6}\selectfont\color{popink}#7
  \end{tcolorbox}
  \noindent\begin{minipage}[t]{0.487\linewidth}
    \begin{tcolorbox}[enhanced,colback=popgreen,colframe=popink,boxrule=2.2pt,arc=10pt,
        drop shadow=popink,left=11pt,right=11pt,top=10pt,bottom=10pt,height=3.1cm,valign=center]
      \tcbox[colback=white,colframe=popink,boxrule=1.3pt,arc=5pt,boxsep=0pt,nobeforeafter,
        left=3pt,right=3pt,top=3pt,bottom=3pt,valign=center]{\qrcode[height=1.9cm]{https://play.google.com/store/apps/details?id=com.luvvix.onbuch&hl=fr}}%
      \hspace{8pt}\begin{minipage}[c]{0.5\linewidth}
        {\popdisplay\fontsize{11}{12.5}\selectfont\color{white}\MakeUppercase{Télécharge OnBuch}}\par\vspace{4pt}
        {\raggedright\popsemi\fontsize{8.4}{11}\selectfont\color{white}Gratuit \textbullet\ Google Play\\Tuteur IA, annales, concours, résultats\par}
      \end{minipage}
    \end{tcolorbox}
  \end{minipage}\hfill
  \begin{minipage}[t]{0.487\linewidth}
    \begin{tcolorbox}[enhanced,colback=poppurple,colframe=popink,boxrule=2.2pt,arc=10pt,
        drop shadow=popink,left=11pt,right=11pt,top=10pt,bottom=10pt,height=3.1cm,valign=center]
      \tikz[baseline=-0.7ex]\node[circle,fill=white,draw=popink,line width=1.3pt,minimum size=1.6cm,inner sep=0]{\popdisplay\fontsize{24}{24}\selectfont\color{poppurple}+};%
      \hspace{8pt}\begin{minipage}[c]{0.6\linewidth}
        {\popdisplay\fontsize{11}{12.5}\selectfont\color{white}\MakeUppercase{Passe à OnBuch+}}\par\vspace{4pt}
        {\raggedright\popsemi\fontsize{8.4}{11}\selectfont\color{white}Tous les cours signés, Tuteur IA illimité, fiches PDF — onglet OnBuch+ de l'app\par}
      \end{minipage}
    \end{tcolorbox}
  \end{minipage}
  \vspace{8pt}
  \popcontenu
  \vfill
  \signatureonbuch
  \restoregeometry
  \newpage
}

% Rangée « Ce cours contient » (page de garde)
\newcommand{\poptile}[3]{% #1 couleur #2 gros texte #3 légende
  \begin{tcolorbox}[enhanced,colback=#1,colframe=popink,boxrule=2pt,arc=9pt,shadow={2.5pt}{-2.5pt}{0pt}{popink},
      left=6pt,right=6pt,top=9pt,bottom=9pt,halign=center,valign=center,height=1.75cm,width=\linewidth,nobeforeafter]
    {\popdisplay\fontsize{12.5}{13}\selectfont\color{white}\MakeUppercase{#2}}\par\vspace{3pt}
    {\centering\popsemi\fontsize{7.8}{9.5}\selectfont\color{white}#3\par}
  \end{tcolorbox}}
\newcommand{\popcontenu}{%
  \par\noindent{\popxboldls\fontsize{7.5}{7.5}\selectfont\color{popmuted}CE COURS SIGNÉ CONTIENT}\par\vspace{7pt}
  \noindent\begin{minipage}[t]{0.235\linewidth}\poptile{popblue}{Cours}{complet, illustré, pas à pas}\end{minipage}\hfill
  \begin{minipage}[t]{0.235\linewidth}\poptile{poporange}{Méthodes}{et exercices résolus}\end{minipage}\hfill
  \begin{minipage}[t]{0.235\linewidth}\poptile{poppink}{Exercices}{type examen + corrigés}\end{minipage}\hfill
  \begin{minipage}[t]{0.235\linewidth}\poptile{popgreen}{Fiche bilan}{l'essentiel à retenir}\end{minipage}\par}

% Sommaire
\newtcolorbox{sommaire}{enhanced,colback=white,colframe=popink,boxrule=2.2pt,arc=10pt,
  drop shadow=popink,left=18pt,right=18pt,top=13pt,bottom=13pt,before skip=6pt,after skip=20pt,
  before upper={\poppill{Sommaire}{poppurple}{white}\par\vspace{9pt}\popsemi\fontsize{10.6}{14.5}\selectfont\color{popink}}}

% Signature de fin de cours — poussée tout en bas de la dernière page
\newcommand{\findecours}{%
  \par\vfill\nopagebreak
  \begin{center}\popsquiggle{poporange}\end{center}\vspace{4pt}
  \signatureonbuch
}
\AtBeginDocument{\def\llbracket{\mathopen{[\mkern-3.5mu[}}\def\rrbracket{\mathclose{]\mkern-3.5mu]}}} % intervalles d'entiers (glyphe ⟦ absent de la police)
% Glyphes absents de Fira Math : redéfinis à partir de glyphes disponibles
\AtBeginDocument{%
  \def\setminus{\mathbin{\backslash}}\let\smallsetminus\setminus
  \def\vdots{\mathord{\vbox{\baselineskip3.2pt\lineskiplimit0pt\kern4pt\hbox{.}\hbox{.}\hbox{.}}}}%
  \def\ddots{\mathinner{\mkern1mu\raise6.5pt\hbox{.}\mkern2mu\raise3.6pt\hbox{.}\mkern2mu\raise0.7pt\hbox{.}\mkern1mu}}%
  \def\triangle{\mathord{\scalebox{0.9}{$\symup{\Delta}$}}}%
  \def\nabla{\mathord{\raisebox{\depth}{\scalebox{1}[-1]{$\symup{\Delta}$}}}}%
  \def\checkmark{\popcheck{popdark}}%
  \def\star{\mathbin{\ast}}\def\bigstar{\mathord{\ast}}%
  \def\diamond{\mathbin{\scalebox{0.6}{$\Diamond$}}}%
  \def\bigcup{\mathop{\mathchoice{\raisebox{-0.3em}{\scalebox{1.7}{$\cup$}}}{\scalebox{1.25}{$\cup$}}{\cup}{\cup}}\displaylimits}%
  \def\bigcap{\mathop{\mathchoice{\raisebox{-0.3em}{\scalebox{1.7}{$\cap$}}}{\scalebox{1.25}{$\cap$}}{\cap}{\cap}}\displaylimits}%
  \def\lesseqqgtr{\mathrel{\lessgtr}}\def\bigoplus{\mathop{\oplus}\displaylimits}%
}
\providecommand{\ssi}{si et seulement si} % abréviation fréquente
\AtBeginDocument{\providecommand{\wideparen}{\overparen}} % arc de cercle

%% FIGFIX-BEGIN v5 — correctifs globaux des figures (docs/onbuch-generation/tools/figfix)
\usepackage{adjustbox}\usepackage{etoolbox}\tcbuselibrary{raster}
% toute figure TikZ/pgfplots plus large que la ligne est réduite à la largeur disponible
\BeforeBeginEnvironment{tikzpicture}{\begin{adjustbox}{max width=\linewidth}}
\AfterEndEnvironment{tikzpicture}{\end{adjustbox}}
% légendes pgfplots lisibles par-dessus les courbes
\pgfplotsset{every axis/.append style={legend style={fill=white,fill opacity=0.92,text opacity=1,draw=black!15,inner sep=3pt}}}
% étiquettes d'arêtes (syntaxe edge["..."]) avec fond blanc
\tikzset{every edge quotes/.append style={fill=white,inner sep=1.2pt}}
%% FIGFIX-END
\begin{document}

\pagegarde{Produire des écrits utilitaires : le rapport}{Français}{1re Tech}{Français – classe de Première technique}{N16}{9 séances (fiche de progression harmonisée)}{Cette leçon outille l'élève de Première technique pour concevoir et rédiger un rapport complet (introduction, corps, conclusion) à partir de situations professionnelles et scolaires concrètes. Elle mobilise le texte explicatif, la distinction entre lexique commun et lexique spécialisé, ainsi que la lecture de textes fonctionnels et d'affiches de sécurité.
\vspace{4pt}
\begin{multicols}{2}\small
\begin{itemize}
\item À la fin de cette leçon, je sais définir le rapport et identifier ses caractéristiques (objectivité, clarté, fidélité aux faits)
\item À la fin de cette leçon, je sais préparer un rapport : collecter, trier et organiser les informations sur le terrain
\item À la fin de cette leçon, je sais distinguer la structure externe (en-tête, titre, signature) et la structure interne (introduction, corps, conclusion) d'un rapport
\item À la fin de cette leçon, je sais reconnaître les procédés du texte explicatif (définition, exemple, comparaison, cause-conséquence)
\item À la fin de cette leçon, je sais différencier le lexique commun du lexique spécialisé et les employer à bon escient
\item À la fin de cette leçon, je sais lire et interpréter une affiche d'indications sécuritaires et un texte fonctionnel
\end{itemize}\end{multicols}}

\begin{sommaire}
\begin{multicols}{2}
\begin{enumerate}
\item Découverte du rapport : définition, types et préparation
\item La structure du rapport : externe et interne
\item Le texte explicatif (stylistique et rhétorique)
\item Lexique commun et lexique spécialisé (sémantique et lexicologie)
\item Lire les textes fonctionnels et les affiches de sécurité (réception)
\item Rédaction du rapport : introduction, corps, conclusion
\item Activité d'intégration
\item Méthodes et exercices résolus
\item Exercices
\item Corrigés des exercices
\item Fiche bilan
\end{enumerate}
\end{multicols}
\end{sommaire}

\begin{objectifs}
\begin{itemize}
\item À la fin de cette leçon, je sais définir le rapport et identifier ses caractéristiques (objectivité, clarté, fidélité aux faits)
\item À la fin de cette leçon, je sais préparer un rapport : collecter, trier et organiser les informations sur le terrain
\item À la fin de cette leçon, je sais distinguer la structure externe (en-tête, titre, signature) et la structure interne (introduction, corps, conclusion) d'un rapport
\item À la fin de cette leçon, je sais reconnaître les procédés du texte explicatif (définition, exemple, comparaison, cause-conséquence)
\item À la fin de cette leçon, je sais différencier le lexique commun du lexique spécialisé et les employer à bon escient
\item À la fin de cette leçon, je sais lire et interpréter une affiche d'indications sécuritaires et un texte fonctionnel
\item À la fin de cette leçon, je sais rédiger l'introduction et la conclusion d'un rapport
\item À la fin de cette leçon, je sais rédiger le corps d'un rapport en paragraphes cohérents et articulés
\item À la fin de cette leçon, je sais produire un rapport complet sur une situation réelle (stage, visite, incident) et justifier ma démarche
\end{itemize}
\end{objectifs}

\begin{prerequis}
\begin{itemize}
\item Le compte rendu (notion vue en classe de Seconde)
\item Les types de textes (narratif, descriptif, informatif)
\item Les connecteurs logiques et la cohérence textuelle
\item La construction du paragraphe (phrase d'idée, développement)
\item La conjugaison des temps du récit et du discours informatif (présent, passé composé, imparfait)
\end{itemize}
\end{prerequis}

\newpage

\coursec{Découverte du rapport : définition, types et préparation}

\courssub{Situation d'accroche et problématique}

Il est 14 h 30, un mardi après-midi, au lycée technique de Nkongsamba. Dans l'atelier de menuiserie, un dégagement de fumée s'échappe d'une raboteuse : en quelques minutes, le feu gagne un amas de copeaux de bois. Grâce au sang-froid du surveillant et à l'extincteur fixé près de la porte, l'incendie est maîtrisé. Bilan : deux établis endommagés, une machine hors d'usage, aucun blessé. Le lendemain matin, le proviseur convoque le délégué de classe et le responsable de l'atelier : « Je veux un rapport circonstancié sur mon bureau avant vendredi. »

Les deux élèves se regardent. Par où commencer ? Comment présenter les faits sans les déformer ? Dans quel ordre ? Avec quels mots ? Faut-il raconter leur peur, donner leur avis sur le responsable de l'accident, ou s'en tenir aux seuls faits observés ?

Cette situation n'a rien d'exceptionnel. Que ce soit pour un stage en entreprise à Douala, une visite d'étude au port autonome, un incident en atelier ou une mission confiée par une autorité, l'élève technicien — et demain le technicien en poste — sera constamment amené à \cle{rendre compte par écrit} de manière claire, objective et structurée. C'est toute l'utilité du \cle{rapport}, écrit utilitaire par excellence.

\textbf{Problématique :} qu'est-ce qu'un rapport, quels sont ses types et ses qualités, et comment le préparer méthodiquement avant même d'écrire la première phrase ?

\courssub{Définition du rapport}

Commençons par l'intuition. Quand quelqu'un te demande « raconte-moi ce qui s'est passé », tu peux répondre de mille façons : avec émotion, avec humour, en exagérant. Mais quand un proviseur, un chef d'atelier ou un maire te demande un \emph{rapport}, il n'attend ni émotion ni exagération : il attend des \textbf{faits}, présentés dans un \textbf{ordre logique}, pour pouvoir \textbf{décider} (réparer la machine, sanctionner, prévenir un nouvel accident). Le rapport est donc un écrit qui sert à quelque chose : c'est un \cle{écrit utilitaire}.

\begin{definition}[Le rapport]
Le \cle{rapport} est un écrit utilitaire qui rend compte de façon \textbf{objective}, \textbf{exacte} et \textbf{structurée} de faits, d'observations ou de l'exécution d'une mission, à un \textbf{destinataire déterminé} (chef de service, proviseur, maire, directeur d'entreprise), en vue de l'informer et souvent de l'aider à prendre une décision.
\end{definition}

Trois traits de cette définition doivent être retenus :

\begin{itemize}
\item Le rapport porte sur des \textbf{faits réels} : on n'invente rien, on n'embellit rien, on ne cache rien.
\item Il est adressé à un \textbf{destinataire précis}, généralement une autorité qui a \emph{commandé} le rapport : on ne rédige pas un rapport pour soi-même.
\item Il obéit à une \textbf{structure codifiée} (structure externe et interne), que nous étudierons dans la prochaine section.
\end{itemize}

\begin{savaistu}[Le rapport de mission dans la fonction publique camerounaise]
Dans la fonction publique camerounaise, tout déplacement en mission — par exemple un inspecteur pédagogique qui visite les lycées techniques de la région du Nord, ou un ingénieur des travaux publics envoyé sur un chantier à Maroua — se termine \textbf{obligatoirement} par un rapport de mission transmis à la hiérarchie. Sans ce document, la mission est considérée comme non accomplie et les indemnités peuvent ne pas être versées. Savoir rédiger un rapport est donc une compétence professionnelle concrète, pas seulement scolaire.
\end{savaistu}

\courssub{Rapport, compte rendu, procès-verbal : ne pas confondre}

Ces trois écrits utilitaires se ressemblent : tous trois relatent des faits réels de manière objective. Pourtant, ils ne se situent ni au même moment, ni dans le même cadre, ni avec le même ton. Le tableau ci-dessous permet de les distinguer.

\begin{popfigure}
\begin{center}
\begin{tikzpicture}[font=\small]
\fill[popblueL] (0,0) rectangle (4.4,1);
\fill[poporangeL] (4.6,0) rectangle (9,1);
\fill[popgreenL] (9.2,0) rectangle (13.6,1);
\node[font=\bfseries] at (2.2,0.5) {Le rapport};
\node[font=\bfseries] at (6.8,0.5) {Le compte rendu};
\node[font=\bfseries] at (11.4,0.5) {Le procès-verbal};
\node[align=left,anchor=north west,text width=4cm] at (0.2,-0.3)
  {\textbf{Destinataire :} une autorité (chef, proviseur, maire) qui l'a \emph{commandé}.\\[2pt]
   \textbf{Moment :} après une mission, un stage, un incident, une enquête.\\[2pt]
   \textbf{Ton et contenu :} constats, analyse, souvent des \emph{suggestions} ou des propositions de décision.};
\node[align=left,anchor=north west,text width=4cm] at (4.8,-0.3)
  {\textbf{Destinataire :} des lecteurs variés (journal, classe, administration).\\[2pt]
   \textbf{Moment :} après un événement (spectacle, match, visite, expérience).\\[2pt]
   \textbf{Ton et contenu :} restitution fidèle et neutre des faits, sans proposer de décision.};
\node[align=left,anchor=north west,text width=4cm] at (9.4,-0.3)
  {\textbf{Destinataire :} les membres d'une assemblée, l'administration, parfois la justice.\\[2pt]
   \textbf{Moment :} rédigé \emph{pendant} une réunion ou une séance, signé en fin de séance.\\[2pt]
   \textbf{Ton et contenu :} valeur officielle, voire juridique ; note les présences, les débats, les votes.};
\end{tikzpicture}
\end{center}
\legende{Tableau comparatif : le rapport, le compte rendu et le procès-verbal se distinguent par le destinataire, le moment de rédaction et le contenu attendu.}
\end{popfigure}

\begin{attention}[Erreur fréquente : confondre rapport et compte rendu]
Beaucoup d'élèves utilisent les deux mots comme des synonymes. Or le \textbf{rapport} est \emph{commandé} par une autorité et comporte souvent des \textbf{constats} et des \textbf{suggestions} (que faut-il faire ? quelle décision prendre ?). Le \textbf{compte rendu}, lui, se contente de restituer fidèlement un événement, sans proposer de décision. À l'examen, si le sujet demande un \emph{rapport}, termine par des suggestions ou des recommandations ; s'il demande un \emph{compte rendu}, reste dans la simple restitution des faits.
\end{attention}

\courssub{Les principaux types de rapports}

Selon la situation qui le provoque, le rapport prend des formes différentes. En voici les cinq types les plus fréquents dans la vie scolaire et professionnelle du technicien.

\begin{itemize}
\item \textbf{Le rapport de stage :} rédigé à la fin d'une période de formation en entreprise (par exemple un stage d'observation à la SODECOTON de Garoua ou dans un garage de Bafoussam), il présente l'entreprise, les tâches accomplies, les acquis et les difficultés rencontrées.
\item \textbf{Le rapport de visite d'étude :} il rend compte d'une sortie pédagogique (visite du port autonome de Douala, d'une usine, d'un musée) : déroulement, observations, enseignements tirés.
\item \textbf{Le rapport d'incident ou d'accident :} comme celui demandé par le proviseur de Nkongsamba, il décrit avec précision les circonstances, les causes apparentes, les dégâts et propose des mesures de prévention.
\item \textbf{Le rapport d'enquête :} il présente les résultats d'une investigation (enquête de satisfaction auprès des élèves, enquête sur l'absentéisme en atelier) : méthode, données recueillies, conclusions.
\item \textbf{Le rapport d'activité :} rédigé périodiquement (mois, trimestre, année) par un responsable (président de club, chef d'atelier, maire), il fait le bilan des actions menées et des perspectives.
\end{itemize}

Tous partagent la même exigence de fond : dire le vrai, complètement, et dans l'ordre.

\courssub{Les quatre qualités d'un bon rapport}

\begin{aretenir}[Les quatre qualités du rapport]
\begin{enumerate}
\item \textbf{L'objectivité :} le rédacteur efface ses sentiments personnels. On écrit « la raboteuse a pris feu à 14 h 30 », et non « heureusement que j'étais là, sinon c'était la catastrophe ! ».
\item \textbf{L'exactitude :} les faits doivent être vérifiables : dates précises, heures, chiffres, noms, lieux. « Deux établis endommagés » et non « plusieurs dégâts ».
\item \textbf{La clarté :} phrases simples, vocabulaire précis, ordre logique (en général chronologique puis thématique).
\item \textbf{La concision :} aller à l'essentiel, supprimer les détails inutiles et les répétitions. Un rapport n'est pas un roman.
\end{enumerate}
\end{aretenir}

\begin{exemplebox}[Objectivité ou subjectivité ?]
Compare ces deux phrases relatives au même incident :
\begin{itemize}
\item \emph{Version subjective :} « À cause de l'incurie habituelle et de la négligence de certains, l'atelier a failli partir en fumée, ce qui aurait été dramatique pour notre pauvre lycée. »
\item \emph{Version objective :} « Le 12 mars, à 14 h 30, un départ de feu s'est déclaré au niveau de la raboteuse, vraisemblablement à la suite d'une surchauffe du moteur. L'incendie a été maîtrisé à 14 h 45 à l'aide de l'extincteur de l'atelier. »
\end{itemize}
Seule la seconde version a sa place dans un rapport : faits, heures, causes prudentes (« vraisemblablement »), aucun jugement.
\end{exemplebox}

\courssub{La préparation du rapport : une étape décisive}

Un bon rapport ne s'écrit pas d'un jet : il se \textbf{prépare}. C'est sur le terrain, pendant le stage, la visite ou juste après l'incident, que se joue la qualité du document final. La préparation comporte quatre phases enchaînées.

\begin{popfigure}
\begin{center}
\begin{tikzpicture}[font=\small, node distance=0.55cm,
  etape/.style={draw=popblue, fill=popblueL, rounded corners=3pt,
    minimum width=2.9cm, minimum height=1.5cm, align=center, text width=2.7cm},
  fleche/.style={-{Stealth[length=3mm]}, very thick, poporange}]
\node[etape, align=center] (e1){\textbf{1. Définir}\\ l'objectif et le destinataire};
\node[etape, right=of e1, align=center] (e2){\textbf{2. Collecter}\\ les informations sur le terrain};
\node[etape, right=of e2, align=center] (e3){\textbf{3. Trier}\\ et classer les informations};
\node[etape, right=of e3, align=center] (e4){\textbf{4. Planifier}\\ élaborer un plan provisoire};
\draw[fleche] (e1) -- (e2);
\draw[fleche] (e2) -- (e3);
\draw[fleche] (e3) -- (e4);
\end{tikzpicture}
\end{center}
\legende{Les quatre phases de la préparation d'un rapport : définir, collecter, trier, planifier.}
\end{popfigure}

\begin{methode}[Préparer un rapport en quatre étapes]
\begin{enumerate}
\item \textbf{Définir l'objectif et le destinataire.} Pose-toi deux questions : \emph{pourquoi} ce rapport (informer ? expliquer un incident ? proposer une décision ?) et \emph{pour qui} (le proviseur ? le chef d'entreprise ?). L'objectif et le destinataire déterminent le ton, le niveau de détail et le vocabulaire.
\item \textbf{Collecter les informations.} Rassemble la matière brute : \textbf{observation} directe sur le terrain, \textbf{interviews} des témoins et des responsables, consultation de \textbf{documents} (registres, consignes, fiches techniques), relevé de \textbf{données chiffrées} (dates, heures, quantités, coûts). Prends des notes immédiatement, avant d'oublier.
\item \textbf{Trier et classer.} Mets chaque information sur une fiche ou une ligne de tableau, puis sépare le \textbf{pertinent} (utile à l'objectif) du \textbf{non pertinent} (anecdote, détail sans rapport avec la mission). Regroupe ensuite les informations retenues par thèmes.
\item \textbf{Planifier.} À partir des thèmes classés, rédige un \textbf{plan provisoire} : introduction (contexte et objectif), corps (constats organisés), conclusion (bilan et suggestions). Ce plan guidera la rédaction.
\end{enumerate}
\end{methode}

\courssub{Étude de cas : trier les notes d'un stage à la SODECOTON}

Pendant son stage d'observation à la SODECOTON (Société de développement du coton) de Garoua, un élève de Première technique a pris 25 notes. Le lendemain de la visite du port autonome de Douala, exercice analogue : sur 25 notes, seules 14 sont retenues après tri. Pourquoi ce tri ? Parce que certaines notes sont des anecdotes (« le chauffeur du bus chantait très bien »), des impressions personnelles (« j'ai trouvé le port immense et impressionnant ») ou des répétitions, qui n'apportent rien au destinataire.

\begin{exoresolu}[Trier des notes de visite]
\textbf{Énoncé.} Voici six notes prises par un élève lors de la visite du port autonome de Douala, en vue d'un rapport destiné au proviseur. Indique pour chacune si elle est pertinente ou non, et justifie.
\begin{enumerate}
\item « Le port reçoit environ 1\,000 navires par an. »
\item « Notre guide portait une chemise très élégante. »
\item « 95\% du commerce extérieur du Cameroun transite par ce port. »
\item « Le déjeuner au restaurant du port était excellent. »
\item « Les conteneurs sont déchargés par des portiques d'une capacité de 40 tonnes. »
\item « Je me suis ennuyé pendant les explications sur la douane. »
\end{enumerate}
\tcblower
\textbf{Solution détaillée.}
\begin{enumerate}
\item \textbf{Pertinente :} donnée chiffrée, vérifiable, qui informe sur l'activité du port.
\item \textbf{Non pertinente :} détail vestimentaire sans lien avec l'objectif du rapport.
\item \textbf{Pertinente :} chiffre essentiel qui montre l'importance économique du port pour le Cameroun.
\item \textbf{Non pertinente :} appréciation personnelle sur un repas ; hors sujet.
\item \textbf{Pertinente :} information technique précise sur les équipements, utile pour des élèves techniciens.
\item \textbf{Non pertinente :} sentiment personnel, contraire à l'exigence d'objectivité.
\end{enumerate}
On retient donc les notes 1, 3 et 5, que l'on regroupera sous le thème « activité et équipements du port ». Les critères de sélection sont : le lien avec l'objectif du rapport, la vérifiabilité du fait, et l'absence de jugement personnel.
\end{exoresolu}

\begin{exemplebox}[Exemple chiffré : le tri en pratique]
Sur les 25 notes prises lors de la visite du port de Douala, seules 14 sont retenues après tri, soit 56\% des notes. Les 11 notes écartées se répartissent ainsi : 4 anecdotes sans lien avec l'objectif, 4 impressions personnelles, 3 répétitions d'une même information. Ce tri n'est pas une perte : c'est un gain de clarté. Le plan provisoire du rapport s'appuiera uniquement sur les 14 notes retenues, regroupées en trois thèmes : présentation du port, observations techniques, enseignements de la visite.
\end{exemplebox}

\begin{exercice}[À toi de jouer]
À la suite de l'incendie de l'atelier de menuiserie de Nkongsamba, le délégué a noté : (a) « l'incendie a débuté à 14 h 30 » ; (b) « le responsable de l'atelier est un homme sévère mais juste » ; (c) « l'extincteur était en bon état de fonctionnement » ; (d) « deux établis et une raboteuse ont été endommagés » ; (e) « tout le monde était paniqué, c'était horrible ». Quelles notes conserveras-tu pour ton rapport au proviseur ? Justifie chaque choix, puis propose un plan provisoire en trois parties.
\end{exercice}

\begin{corrige}[Exercice 1]
Notes à conserver : (a), (c) et (d), car ce sont des faits exacts, vérifiables et utiles à la décision du proviseur. On écarte (b), jugement sur une personne, et (e), expression d'un sentiment : contraires à l'objectivité. Plan provisoire possible : 1. Introduction : contexte et commande du rapport ; 2. Corps : déroulement des faits (14 h 30, maîtrise du feu), bilan matériel (deux établis, une raboteuse), état du matériel de sécurité (extincteur fonctionnel) ; 3. Conclusion : suggestions (contrôle des machines, formation à la sécurité).
\end{corrige}

\begin{aretenir}[L'essentiel de la section]
Le rapport est un écrit utilitaire objectif, exact et structuré, commandé par une autorité et destiné à l'informer pour l'aider à décider. Il se distingue du compte rendu (simple restitution) et du procès-verbal (rédigé pendant une séance, à valeur officielle). Il existe plusieurs types de rapports (stage, visite, incident, enquête, activité), mais tous exigent objectivité, exactitude, clarté et concision. Avant toute rédaction, la préparation suit quatre phases : définir l'objectif et le destinataire, collecter les informations sur le terrain, trier le pertinent du non pertinent, et élaborer un plan provisoire. C'est cette préparation rigoureuse qui fera la qualité de la structure et de la rédaction, que nous étudierons dans les sections suivantes.
\end{aretenir}

\coursec{La structure du rapport : externe et interne}

Dans la section précédente, nous avons découvert ce qu'est un rapport : un écrit utilitaire par lequel une personne rend compte, de manière objective et ordonnée, d'une mission, d'un incident, d'une visite ou d'un stage à un destinataire déterminé. Nous avons aussi vu qu'avant d'écrire, il faut \cle{préparer} : rassembler les faits, les dates, les chiffres, les témoignages. Mais une fois les informations réunies, une question se pose immédiatement : \emph{comment les présenter sur la page ?} Un bon rapport n'est pas seulement un texte juste ; c'est un texte \cle{structuré}. C'est précisément l'objet de cette section : apprendre à construire le \textbf{cadre matériel} du document (structure externe) et l'\textbf{organisation du contenu} (structure interne).

\courssub{La structure externe : le cadre matériel du document}

\textbf{Intuition.} Imaginez que le proviseur du lycée technique de Bafoussam reçoive une feuille de papier couverte de phrases, sans nom, sans date, sans titre, sans signature. Comment saurait-il qui écrit, à quel sujet, et à quel moment ? Cette feuille serait inutilisable, même si son contenu est excellent. Avant même de lire le texte, le lecteur d'un rapport a besoin de \emph{repères visibles en un coup d'œil}. Ces repères forment la structure externe.

\begin{definition}[Structure externe]
La \cle{structure externe} d'un rapport est l'ensemble des éléments matériels et administratifs qui encadrent le texte proprement dit : l'\textbf{en-tête}, le \textbf{lieu et la date}, l'\textbf{objet}, le \textbf{destinataire}, le \textbf{titre}, les \textbf{références}, les \textbf{pièces jointes} et la \textbf{signature}. Elle permet d'identifier immédiatement qui écrit, à qui, quand, et à quel sujet.
\end{definition}

\textbf{Les éléments de la structure externe}\par
Présentons chaque élément, sa place habituelle sur la page et sa fonction.

\begin{center}
\begin{tabularx}{\linewidth}{|l|X|X|}\hline
\rowcolor{popblueL} \textbf{Élément} & \textbf{Place habituelle} & \textbf{Fonction} \\ \hline
En-tête (nom, fonction du rédacteur) & En haut à gauche & Identifier l'auteur du rapport et sa qualité (ex. : délégué de classe, chef d'atelier, stagiaire). \\ \hline
Lieu et date & En haut à droite & Situer le document dans le temps et dans l'espace (ex. : Bafoussam, le 12 mars 2026). \\ \hline
Destinataire & Sous l'en-tête ou en haut à droite & Nommer la personne ou le service qui reçoit le rapport (ex. : Monsieur le Proviseur). \\ \hline
Objet & Au-dessus du texte, souvent souligné & Annoncer en une phrase courte le sujet exact du rapport. \\ \hline
Références / numéro & Près de l'objet & Permettre le classement et le suivi administratif du document (ex. : N\textsuperscript{o} 15/LTB/AT/2026). \\ \hline
Titre & Au centre, en évidence & Résumer la nature du document (ex. : Rapport d'incident). \\ \hline
Pièces jointes & En bas, après la signature & Signaler les documents annexés (photos, factures, procès-verbaux). \\ \hline
Signature & En bas à droite, précédée du nom & Engager la responsabilité du rédacteur et donner valeur officielle au document. \\ \hline
\end{tabularx}
\end{center}

\begin{attention}[Deux erreurs qui coûtent cher]
Deux oublis sont fréquents et graves. Premièrement, \textbf{omettre l'objet} : le lecteur ne sait pas de quoi parle le document avant de l'avoir lu en entier, ce qui est contraire à la logique d'un écrit utilitaire. Deuxièmement, \textbf{signer avant d'avoir relu} : une faute de date ou de chiffre dans un rapport signé engage son auteur. Retenez surtout ceci : \emph{un rapport non signé n'a aucune valeur administrative}. Relisez toujours, puis signez.
\end{attention}

\begin{popfigure}
\begin{center}
\begin{tikzpicture}
\draw[thick, popdark] (0,0) rectangle (10.5,13.2);
% En-tête
\draw[fill=popblueL] (0.4,11.6) rectangle (5.2,12.8);
\node[anchor=west, font=\small] at (0.6,12.5) {NDOUMBE Alain};
\node[anchor=west, font=\small] at (0.6,12.0) {Chef d'atelier de mécanique};
\draw[->, thick, popblue] (5.2,12.2) -- (6.4,12.2) node[right, font=\footnotesize, text=popblue] {\textbf{1. En-tête}};
% Lieu et date
\draw[fill=poporangeL] (6.6,11.6) rectangle (10.1,12.8);
\node[font=\small] at (8.35,12.2) {Bafoussam, le 12/03/2026};
\draw[->, thick, poporange] (8.35,11.6) -- (8.35,11.1) node[below, font=\footnotesize, text=poporange] {\textbf{2. Lieu et date}};
% Destinataire
\draw[fill=popgreenL] (6.6,10.2) rectangle (10.1,11.0);
\node[font=\small] at (8.35,10.6) {À Monsieur le Proviseur};
\draw[->, thick, popgreen] (6.6,10.6) -- (5.5,10.6) node[left, font=\footnotesize, text=popgreen] {\textbf{3. Destinataire}};
% Objet
\draw[fill=poppinkL] (0.4,9.3) rectangle (10.1,10.0);
\node[font=\small] at (5.25,9.65) {Objet : Rapport d'incident survenu à l'atelier de mécanique};
\draw[->, thick, poppink] (5.25,9.3) -- (5.25,8.85) node[below, font=\footnotesize, text=poppink] {\textbf{4. Objet}};
% Titre
\node[font=\bfseries] at (5.25,8.35) {RAPPORT D'INCIDENT};
\draw[->, thick, poppurple] (7.6,8.35) -- (8.6,8.35) node[right, font=\footnotesize, text=poppurple] {\textbf{5. Titre}};
% Introduction
\draw[fill=popgoldL] (0.4,6.9) rectangle (10.1,7.9);
\node[font=\small] at (5.25,7.4) {Introduction : contexte, circonstances, objectif};
\draw[->, thick, popgold] (0.4,7.4) -- (-0.5,7.4) node[left, font=\footnotesize, text=popgold] {\textbf{6. Intro}};
% Corps
\draw[fill=popcream] (0.4,3.4) rectangle (10.1,6.5);
\node[font=\small] at (5.25,5.7) {Corps : partie 1 (constat)};
\node[font=\small] at (5.25,5.0) {partie 2 (causes)};
\node[font=\small] at (5.25,4.3) {partie 3 (solutions)};
\draw[->, thick, popdark] (0.4,4.95) -- (-0.5,4.95) node[left, font=\footnotesize, text=popdark] {\textbf{7. Corps}};
% Conclusion
\draw[fill=popgoldL] (0.4,2.2) rectangle (10.1,3.0);
\node[font=\small] at (5.25,2.6) {Conclusion : bilan et recommandations};
\draw[->, thick, popgold] (10.1,2.6) -- (11.0,2.6) node[right, font=\footnotesize, text=popgold] {\textbf{8. Concl.}};
% Signature
\draw[fill=popgreenL] (7.0,0.9) rectangle (10.1,1.8);
\node[font=\small] at (8.55,1.35) {Signature : NDOUMBE A.};
\draw[->, thick, popgreen] (8.55,0.9) -- (8.55,0.45) node[below, font=\footnotesize, text=popgreen] {\textbf{9. Signature}};
% Pièces jointes
\node[anchor=west, font=\footnotesize] at (0.4,1.35) {Pièces jointes : 2 photos, 1 devis};
\draw[->, thick, popmuted] (0.4,1.0) -- (0.4,0.5) node[below, font=\footnotesize, text=popmuted] {\textbf{10. P.J.}};
\end{tikzpicture}
\end{center}
\legende{Maquette d'une page de rapport : les dix rubriques de la structure externe et interne, numérotées de 1 à 10, dans leur disposition habituelle.}
\end{popfigure}

\courssub{La structure interne : l'organisation du contenu}

\textbf{Intuition.} Une fois le cadre posé, il faut remplir le document. Mais dans quel ordre raconter les faits ? Un lecteur pressé veut d'abord savoir \emph{pourquoi} on lui écrit, ensuite \emph{ce qui s'est passé}, enfin \emph{ce qu'on propose}. Cette logique en trois temps — introduction, corps, conclusion — est la structure interne.

\begin{definition}[Structure interne]
La \cle{structure interne} d'un rapport est l'organisation du contenu en trois grandes parties : l'\textbf{introduction} (contexte, circonstances de la mission, objectif, annonce du plan), le \textbf{corps} (constats et faits organisés en parties) et la \textbf{conclusion} (bilan, suggestions ou recommandations).
\end{definition}

\textbf{L'introduction.}\par
Elle répond à quatre questions : \emph{dans quel contexte} le rapport est-il rédigé ? \emph{dans quelles circonstances} (ordre de mission, demande du supérieur, événement survenu) ? \emph{quel est l'objectif} du document ? \emph{comment le propos est-il organisé} (annonce du plan) ? Une introduction de rapport est brève : quatre à huit lignes suffisent.

\textbf{Le corps.}\par
C'est le cœur du rapport. On y présente les \cle{constats}, c'est-à-dire les faits observés, décrits objectivement, sans jugement personnel. Ces constats sont regroupés en \textbf{parties} numérotées ou titrées, chacune développant un aspect du sujet.

\textbf{La conclusion.}\par
Elle dresse le \textbf{bilan} (ce qu'il faut retenir de l'essentiel) et formule des \textbf{suggestions} ou des \textbf{recommandations} concrètes à l'attention du destinataire. C'est souvent la partie que le lecteur administratif lit en premier : elle doit être claire et directement exploitable.

\begin{popfigure}
\begin{center}
\begin{tikzpicture}[
  level 1/.style={sibling distance=42mm, level distance=16mm},
  level 2/.style={sibling distance=26mm, level distance=16mm},
  every node/.style={draw, rounded corners, align=center, font=\small, thick},
  edge from parent/.style={draw=popmuted, thick}]
\node[fill=popblueL] {LE RAPPORT}
  child { node[fill=poppinkL, align=center]{Introduction\\ \footnotesize contexte, objectif} }
  child { node[fill=poporangeL] {Corps}
    child { node[fill=poporangeL!60, align=center]{\footnotesize Partie 1\\ \footnotesize constat} }
    child { node[fill=poporangeL!60, align=center]{\footnotesize Partie 2\\ \footnotesize causes} }
    child { node[fill=poporangeL!60, align=center]{\footnotesize Partie 3\\ \footnotesize solutions} } }
  child { node[fill=popgreenL, align=center]{Conclusion\\ \footnotesize bilan, recommandations} };
\end{tikzpicture}
\end{center}
\legende{Arbre de la structure interne du rapport : trois blocs obligatoires, le corps se subdivisant en parties.}
\end{popfigure}

\courssub{Choisir le plan du corps}

Tous les rapports ne s'organisent pas de la même façon. Le plan du corps dépend du \textbf{type de rapport} et de la nature des informations à présenter. Trois plans sont à connaître.

\begin{center}
\begin{tabularx}{\linewidth}{|l|X|X|}\hline
\rowcolor{popblueL} \textbf{Plan} & \textbf{Principe} & \textbf{Rapport concerné} \\ \hline
Chronologique & Les faits sont présentés dans l'ordre du temps : avant, pendant, après. & Rapport d'incident, d'accident, d'événement. \\ \hline
Thématique & Les informations sont regroupées par thèmes ou par aspects du sujet. & Rapport de stage, de visite, de réunion. \\ \hline
Constat -- causes -- solutions & On décrit d'abord la situation, puis on explique ses causes, enfin on propose des solutions. & Rapport de problème technique, d'inspection, de dysfonctionnement. \\ \hline
\end{tabularx}
\end{center}

\begin{methode}[Choisir le bon plan en trois questions]
\begin{enumerate}
\item \textbf{Mon sujet est-il un événement qui s'est déroulé dans le temps ?} Si oui, adoptez le plan chronologique (ex. : un incendie dans l'atelier : les faits avant, le déroulement, les suites).
\item \textbf{Mon sujet comporte-t-il plusieurs aspects à décrire sans lien de temps ?} Si oui, adoptez le plan thématique (ex. : un rapport de stage : l'accueil, les tâches effectuées, les acquis).
\item \textbf{Mon sujet est-il un problème à résoudre ?} Si oui, adoptez le plan constat -- causes -- solutions (ex. : les pannes répétées d'une machine).
\end{enumerate}
En cas de doute, le plan constat -- causes -- solutions est le plus sûr : il convient à la plupart des situations professionnelles de la vie technique.
\end{methode}

\begin{exemplebox}[Un exemple chiffré à retenir]
Un rapport de stage de Première technique comporte en moyenne \textbf{3 à 5 parties} dans le corps. En dessous de trois, le travail paraît pauvre ; au-delà de cinq, le lecteur se perd : il faut alors \emph{regrouper les idées} proches sous un même titre de partie. Par exemple, au lieu de six parties (« l'accueil », « l'atelier », « les machines », « les ouvriers », « les horaires », « la sécurité »), on regroupe en trois : « les conditions d'accueil et de travail », « les activités réalisées », « l'organisation et la sécurité ».
\end{exemplebox}

\courssub{Analyse guidée d'un rapport modèle}

Lisons ensemble un rapport d'incident rédigé au lycée technique de Bafoussam, et repérons chaque rubrique.

\begin{exemplebox}[Rapport d'incident — atelier de mécanique]
\textbf{M. FOTSO Jean, surveillant d'atelier} \hfill \textbf{Bafoussam, le 12 mars 2026}

\textbf{À Monsieur le Proviseur du Lycée technique de Bafoussam}

\textbf{N\textsuperscript{o} 15/LTB/AM/2026}

\textbf{Objet :} Rapport d'incident survenu à l'atelier de mécanique le 11 mars 2026.

\textbf{RAPPORT D'INCIDENT}

\emph{Introduction.} Suite à l'incident survenu le 11 mars 2026 vers 10 h 30 dans l'atelier de mécanique, et conformément à vos instructions, j'ai l'honneur de vous présenter le rapport suivant, organisé en trois parties : le déroulement des faits, les causes de l'incident et les mesures proposées.

\emph{I. Déroulement des faits.} À 10 h 30, pendant la séance de travaux pratiques de la classe de Première F4, la disqueuse utilisée par un élève a projeté des étincelles sur un chiffon imbibé d'huile, qui s'est enflammé. L'élève a été légèrement brûlé à la main gauche. Le feu a été éteint en deux minutes à l'aide de l'extincteur de l'atelier.

\emph{II. Causes de l'incident.} L'enquête menée sur place montre que les chiffons huileux étaient stockés à moins d'un mètre du poste de meulage, contrairement à la consigne affichée, et que l'élève ne portait pas ses gants de protection.

\emph{III. Mesures proposées.} Je suggère : 1\textsuperscript{o} l'installation d'un bac métallique fermé pour les chiffons usagés ; 2\textsuperscript{o} le rappel des consignes de sécurité en début de chaque séance ; 3\textsuperscript{o} le contrôle systématique du port des équipements de protection.

\emph{Conclusion.} L'incident, sans gravité pour l'élève blessé, révèle des manquements dans l'application des règles de sécurité. La mise en œuvre des mesures proposées devrait empêcher toute récidive.

\textbf{Pièces jointes :} 2 photos, 1 certificat médical. \hfill \textbf{Signature : FOTSO J.}
\end{exemplebox}

Repérage : l'\textbf{en-tête} (nom et fonction), le \textbf{lieu et la date}, le \textbf{destinataire}, le \textbf{numéro de référence} et l'\textbf{objet} forment la structure externe ; le \textbf{titre} annonce la nature du document ; l'\textbf{introduction} donne le contexte, la circonstance (« conformément à vos instructions »), l'objectif et annonce le plan ; le \textbf{corps} suit le plan constat -- causes -- solutions, adapté à un problème technique ; la \textbf{conclusion} fait le bilan et recommande ; les \textbf{pièces jointes} et la \textbf{signature} ferment le document.

\begin{savaistu}[L'objet et la signature à l'examen]
Au Probatoire et au Baccalauréat technique, la rubrique « objet » et la signature sont \emph{systématiquement} évaluées dans la grille de correction de la production d'écrits. Un candidat qui rédige un excellent texte mais oublie l'objet ou la signature perd des points de présentation, alors que ces deux lignes ne demandent que quelques secondes d'attention. Prenez l'habitude, dès maintenant, de vérifier ces rubriques avant de rendre votre copie.
\end{savaistu}

\courssub{Exercice de remise en ordre}

\begin{exoresolu}[Reconstituer un rapport mélangé]
\textbf{Énoncé.} Les rubriques suivantes d'un rapport de visite ont été mélangées. Remettez-les dans l'ordre correct et justifiez votre classement.

A. Conclusion : la visite montre que l'entreprise respecte les normes ; je recommande d'y envoyer d'autres classes.

B. Objet : Rapport de visite de l'entreprise SOCAPALM de Dibombari.

C. Signature : KAMGA L.

D. Douala, le 5 février 2026.

E. Introduction : dans le cadre de notre formation, notre classe a visité l'entreprise SOCAPALM le 3 février 2026 ; ce rapport présente l'organisation de l'entreprise, les activités observées et les enseignements tirés.

F. Corps, partie 1 : présentation de l'entreprise (effectifs, ateliers, production).

G. KAMGA Laure, déléguée de la classe de Première F3.

H. Corps, partie 2 : les activités observées et les enseignements tirés de la visite.

I. À Monsieur le Principal du collège technique.

\tcblower
\textbf{Solution détaillée.} On applique la logique de la structure externe (du haut vers le bas de la page), puis la structure interne (introduction, corps, conclusion) :

\begin{enumerate}
\item \textbf{G} : en-tête (identité et fonction du rédacteur), en haut à gauche.
\item \textbf{D} : lieu et date, en haut à droite.
\item \textbf{I} : destinataire.
\item \textbf{B} : objet, avant le texte.
\item \textbf{E} : introduction (contexte, circonstance, annonce du plan).
\item \textbf{F} puis \textbf{H} : le corps, en respectant l'ordre annoncé dans l'introduction (d'abord l'organisation de l'entreprise, ensuite les activités observées).
\item \textbf{A} : conclusion (bilan et recommandation).
\item \textbf{C} : signature, en bas à droite.
\end{enumerate}

Ordre final : \textbf{G -- D -- I -- B -- E -- F -- H -- A -- C}. La justification repose sur deux principes : les éléments externes précèdent toujours le texte, et l'ordre des parties du corps doit correspondre à l'annonce du plan faite dans l'introduction.
\end{exoresolu}

\begin{exercice}[À vous de structurer]
Le chef d'atelier d'électrotechnique de votre lycée vous demande un rapport sur la panne répétée du tableau électrique de l'atelier. 1) Rédigez les éléments de la structure externe (en-tête, lieu et date, destinataire, objet, titre). 2) Choisissez le plan du corps le plus adapté et justifiez votre choix. 3) Rédigez les titres des trois parties du corps.
\end{exercice}

\begin{corrige}[Exercice : À vous de structurer]
1) Structure externe attendue : nom et fonction du rédacteur (ex. : élève délégué de Première F2) ; lieu et date ; destinataire (Monsieur le chef d'atelier d'électrotechnique) ; objet : rapport sur les pannes répétées du tableau électrique de l'atelier ; titre : Rapport technique. 2) Le plan adapté est le plan \textbf{constat -- causes -- solutions}, car il s'agit d'un problème technique à résoudre et non d'un événement unique (plan chronologique) ni d'une description par aspects (plan thématique). 3) Titres possibles : I. Constat : description des pannes observées ; II. Causes : vétusté du matériel et surcharge des circuits ; III. Solutions : remplacement des disjoncteurs défectueux et répartition des charges.
\end{corrige}

\begin{aretenir}[L'essentiel de la section]
\begin{itemize}
\item La \cle{structure externe} est le cadre matériel du rapport : en-tête, lieu et date, destinataire, objet, références, titre, pièces jointes, signature.
\item La \cle{structure interne} organise le contenu : introduction (contexte, circonstances, objectif, annonce du plan), corps (constats en parties), conclusion (bilan et recommandations).
\item Le plan du corps se choisit selon le type de rapport : chronologique pour un incident, thématique pour un stage ou une visite, constat -- causes -- solutions pour un problème technique.
\item Un rapport sans objet ni signature est un rapport sans valeur : relisez, puis signez.
\end{itemize}
\end{aretenir}

Maintenant que le squelette du rapport est en place, il reste à lui donner sa chair linguistique : la section suivante étudiera le \emph{texte explicatif}, c'est-à-dire les procédés de style et de rhétorique qui permettent d'expliquer clairement les faits dans le corps du rapport.

\coursec{Le texte explicatif (stylistique et rhétorique)}

Dans la section précédente, nous avons appris à construire la \cle{structure} du rapport : sa structure externe (page de garde, sommaire, annexes) et sa structure interne (introduction, corps, conclusion). Mais une belle charpente ne suffit pas : encore faut-il que chaque paragraphe du corps \emph{explique} réellement quelque chose au lecteur. Or le lecteur d'un rapport — le chef d'atelier, le directeur, l'inspecteur — n'était pas sur place. Il veut comprendre \emph{pourquoi} un incident s'est produit et \emph{comment} fonctionne une situation. C'est précisément le rôle du \cle{texte explicatif}, que nous étudions maintenant.

\courssub{Qu'est-ce qu'un texte explicatif ?}

Imaginons deux phrases tirées de deux rapports d'atelier :

\begin{itemize}
\item Rapport A : « La scie circulaire s'est arrêtée mardi à 10 h. »
\item Rapport B : « La scie circulaire s'est arrêtée mardi à 10 h \emph{parce que} le fusible de protection a fondu, \emph{à la suite d'} une surtension du réseau électrique. »
\end{itemize}

Le rapport A se contente de \emph{raconter} un fait. Le rapport B \emph{explique} ce fait : il donne la cause (le fusible fondu) et l'enchaînement (la surtension). Le lecteur du rapport B comprend, et il peut agir (faire vérifier l'installation électrique). C'est toute la différence entre un texte narratif ou descriptif et un texte explicatif.

\begin{definition}[Le texte explicatif]
Un \cle{texte explicatif} est un texte qui donne à \emph{comprendre} un phénomène, un fonctionnement ou une situation. Il répond à deux questions essentielles : \textbf{pourquoi ?} (les causes) et \textbf{comment ?} (les mécanismes, les étapes). Il s'appuie sur des \cle{procédés d'explication} :
\begin{itemize}
\item \textbf{définir} : donner le sens précis d'un terme (« Un court-circuit est un contact accidentel entre deux conducteurs ») ;
\item \textbf{exemplifier} : donner un cas concret (« par exemple, le non-port des gants ») ;
\item \textbf{illustrer} : décrire une situation qui rend l'idée visible ;
\item \textbf{comparer} : rapprocher l'inconnu du connu (« le fusible agit comme un garde du corps qui se sacrifie pour protéger la machine ») ;
\item \textbf{chiffrer} : citer des données, des statistiques, des mesures.
\end{itemize}
\end{definition}

\begin{attention}[Le piège de l'énumération sans liens]
L'erreur la plus fréquente dans les rapports d'élèves est d'\emph{enchaîner les faits} les uns après les autres, sans aucun lien logique : « Il y a eu un accident. L'apprenti n'avait pas de lunettes. La machine était en marche. Le chef était absent. » Le lecteur ne comprend ni les \emph{causes} ni les \emph{conséquences} : qui a provoqué quoi ? Un fait sans explication est une information morte. Dans un rapport, chaque constat doit être relié à sa cause ou à sa conséquence par un \cle{connecteur logique}.
\end{attention}

\courssub{Les relations logiques : cause, conséquence, but}

Expliquer, c'est établir des \cle{relations logiques} entre les faits. Trois relations dominent dans le texte explicatif :

\begin{itemize}
\item la \textbf{cause} (pourquoi ?) : \emph{parce que, puisque, car, à cause de, en raison de, grâce à, à la suite de} ;
\item la \textbf{conséquence} (avec quel résultat ?) : \emph{donc, c'est pourquoi, si bien que, par conséquent, de sorte que, ainsi} ;
\item le \textbf{but} (dans quelle intention ?) : \emph{pour, afin de, dans le but de, de manière à, en vue de}.
\end{itemize}

\begin{popfigure}
\begin{center}
\begin{tikzpicture}[font=\small]
\node[draw=popblue, fill=popblueL, rounded corners, minimum width=3.4cm, minimum height=2.2cm, align=center] (cause) at (0,0) {\textbf{CAUSE}\\ \emph{pourquoi ?}\\[2pt] parce que\\ puisque, car\\ à cause de\\ en raison de};
\node[draw=poporange, fill=poporangeL, rounded corners, minimum width=3.4cm, minimum height=2.2cm, align=center] (cons) at (4.6,0) {\textbf{CONSÉQUENCE}\\ \emph{quel résultat ?}\\[2pt] donc\\ c'est pourquoi\\ si bien que\\ par conséquent};
\node[draw=popgreen, fill=popgreenL, rounded corners, minimum width=3.4cm, minimum height=2.2cm, align=center] (but) at (9.2,0) {\textbf{BUT}\\ \emph{quelle intention ?}\\[2pt] pour\\ afin de\\ dans le but de\\ de manière à};
\end{tikzpicture}
\end{center}
\legende{Les trois relations logiques de l'explication et leurs connecteurs principaux.}
\end{popfigure}

\begin{exemplebox}[Une même situation, trois relations]
Situation : un apprenti porte des lunettes de protection.
\begin{itemize}
\item \textbf{Cause} : L'apprenti porte des lunettes \emph{parce que} des copeaux de métal peuvent jaillir.
\item \textbf{But} : L'apprenti porte des lunettes \emph{afin de} protéger ses yeux des projections.
\item \textbf{Conséquence} : Des copeaux peuvent jaillir ; \emph{c'est pourquoi} l'apprenti porte des lunettes.
\end{itemize}
Les trois phrases expriment la même réalité, mais organisent les faits différemment. Le rédacteur du rapport choisit la relation selon ce qu'il veut mettre en avant.
\end{exemplebox}

\courssub{Les marques d'objectivité}

Un rapport n'est pas un journal intime : il doit être \cle{objectif}, c'est-à-dire présenter les faits et leurs explications sans laisser passer les émotions ou les opinions personnelles du rédacteur. L'objectivité se reconnaît à plusieurs marques linguistiques :

\begin{itemize}
\item le \textbf{vocabulaire neutre et précis} : on écrit « l'apprenti a été légèrement blessé à la main gauche » et non « le pauvre garçon s'est affreusement massacré la main » ;
\item le \textbf{présent de vérité générale} pour les explications permanentes : « la surchauffe provoque la fonte du fusible » ;
\item le \textbf{passé factuel} (passé composé, passé simple) pour rapporter les événements constatés : « la machine s'est arrêtée à 10 h » ;
\item l'\textbf{absence du « je » affectif} : on préfère « l'équipe a constaté que... », « il apparaît que... », « on observe que... » aux formules comme « je pense que », « je trouve que », « j'ai été choqué ».
\end{itemize}

\begin{propriete}[La règle d'or du paragraphe de rapport]
Dans un rapport, \textbf{chaque constat doit être suivi d'une explication ou d'une preuve} : un chiffre, une observation, un témoignage, une mesure. Un constat sans preuve est une rumeur ; un constat expliqué et prouvé est un fait établi. Cette règle est exigible dans toute production de rapport au Probatoire et au Baccalauréat technique.
\end{propriete}

\courssub{L'organisation du paragraphe explicatif}

Le paragraphe explicatif est la brique de base du corps du rapport. Il se construit en trois temps : on énonce d'abord l'\textbf{idée principale} (le constat), puis on développe l'\textbf{explication} (les causes, le mécanisme), enfin on apporte l'\textbf{illustration} (exemple, chiffre, observation précise).

\begin{popfigure}
\begin{center}
\begin{tikzpicture}[font=\small]
\node[draw=popblue, fill=popblueL, rounded corners, minimum width=10.5cm, minimum height=1.1cm, align=center] (idee) at (0,3.4) {\textbf{1. IDÉE PRINCIPALE} : le constat — « Les accidents de l'atelier ont augmenté ce trimestre. »};
\node[draw=poporange, fill=poporangeL, rounded corners, minimum width=9cm, minimum height=1.1cm, align=center] (expl) at (0,1.9) {\textbf{2. EXPLICATION} : les causes — « \emph{parce que} les consignes de sécurité ne sont pas toujours respectées »};
\node[draw=popgreen, fill=popgreenL, rounded corners, minimum width=7.5cm, minimum height=1.1cm, align=center] (ex) at (0,0.4) {\textbf{3. ILLUSTRATION} : la preuve — « \emph{par exemple}, 60 \% des incidents sont dus au non-port des équipements »};
\draw[-{Stealth}, thick, popdark] (idee.south) -- (expl.north);
\draw[-{Stealth}, thick, popdark] (expl.south) -- (ex.north);
\end{tikzpicture}
\end{center}
\legende{Schéma du paragraphe explicatif : trois blocs emboîtés, de l'idée générale vers la preuve concrète.}
\end{popfigure}

\begin{methode}[Rédiger un paragraphe explicatif en trois temps]
\begin{enumerate}
\item \textbf{Formuler l'idée principale} en une phrase claire : c'est le constat ou le phénomène à expliquer. Cette phrase doit pouvoir se lire seule.
\item \textbf{Développer l'explication} : répondre à la question « pourquoi ? » ou « comment ? » en utilisant les connecteurs de cause, de conséquence ou de but (\emph{parce que, donc, afin de}...). Employer le présent de vérité générale pour les mécanismes.
\item \textbf{Apporter l'illustration} : un exemple concret, un chiffre, une observation datée et localisée qui prouve l'explication. Introduire l'exemple par \emph{par exemple, ainsi, c'est le cas de}.
\end{enumerate}
Vérification finale : relire le paragraphe en supprimant les connecteurs. Si le lecteur ne comprend plus les liens entre les faits, c'est que les connecteurs faisaient bien leur travail — et qu'ils étaient indispensables.
\end{methode}

\courssub{Analyse d'un extrait de texte explicatif : les accidents en atelier}

Lisons le paragraphe suivant, extrait du rapport de fin de stage d'un élève de Première technique en atelier de mécanique :

\begin{exemplebox}[Texte à analyser]
« Les accidents de travail dans l'atelier ont augmenté de 20 \% au cours du dernier trimestre. Cette hausse s'explique \emph{principalement par} le non-respect des consignes de sécurité : \emph{en effet}, 60 \% des incidents sont dus au non-port des équipements de protection individuelle. \emph{Par exemple}, le 14 mars, un apprenti a été blessé à l'œil \emph{parce qu'}il meulait sans lunettes de protection. \emph{Par conséquent}, le chef d'atelier a décidé de renforcer le contrôle des équipements \emph{afin de} réduire le nombre d'accidents. »
\end{exemplebox}

Repérons ensemble les procédés explicatifs :

\begin{center}
\begin{tabularx}{\linewidth}{|l|X|}\hline
\rowcolor{popblueL} \textbf{Élément du texte} & \textbf{Procédé ou marque repérée} \\ \hline
« Les accidents... ont augmenté de 20 \% » & Idée principale (constat) appuyée par un \emph{chiffre} \\ \hline
« s'explique principalement par » & Connecteur de \emph{cause} \\ \hline
« 60 \% des incidents sont dus au non-port... » & Explication par une \emph{donnée chiffrée} \\ \hline
« Par exemple, le 14 mars, un apprenti... » & \emph{Exemple} daté et précis (illustration) \\ \hline
« parce qu'il meulait sans lunettes » & Connecteur de \emph{cause} \\ \hline
« Par conséquent, le chef d'atelier a décidé... » & Connecteur de \emph{conséquence} \\ \hline
« afin de réduire le nombre d'accidents » & Connecteur de \emph{but} \\ \hline
Vocabulaire : « incidents », « équipements de protection individuelle » & Lexique \emph{neutre et précis}, aucun « je » affectif \\ \hline
\end{tabularx}
\end{center}

Le paragraphe suit exactement le schéma idée $\rightarrow$ explication $\rightarrow$ illustration, et chaque fait est relié au suivant par une relation logique explicite. C'est un modèle à imiter.

\courssub{Application au rapport : des notes télégraphiques aux paragraphes explicatifs}

Sur le terrain, l'élève en stage ou en visite prend des \cle{notes télégraphiques} : courtes, sans liens, sans phrases complètes. Le travail de rédaction consiste à transformer ces notes en paragraphes explicatifs complets.

\begin{exoresolu}[Transformer des notes en paragraphe explicatif]
\textbf{Énoncé.} Voici les notes prises par un élève lors d'une visite d'atelier : « Panne machine à coudre industrielle. Mardi 14 h. Courroie usée. Machine pas entretenue depuis 6 mois. Production arrêtée 2 heures. Réparateur appelé. Chef atelier : carnet d'entretien obligatoire désormais. » Rédigez un paragraphe explicatif de rapport à partir de ces notes.
\tcblower
\textbf{Solution détaillée.} On applique la méthode en trois temps.

\emph{Étape 1 — l'idée principale} : quel est le fait central ? La panne de la machine à coudre industrielle, mardi à 14 h.

\emph{Étape 2 — l'explication} : pourquoi la panne ? Parce que la courroie était usée, elle-même conséquence d'un défaut d'entretien (six mois sans maintenance). On relie avec des connecteurs de cause.

\emph{Étape 3 — l'illustration et les conséquences} : les faits précis (arrêt de deux heures, intervention du réparateur) et la mesure prise (carnet d'entretien), reliés par des connecteurs de conséquence et de but.

\textbf{Paragraphe rédigé :}

« Mardi à 14 h, la machine à coudre industrielle de l'atelier est tombée en panne. Cette panne s'explique \emph{par} l'usure de la courroie de transmission, \emph{car} la machine n'avait pas été entretenue depuis six mois. \emph{Par conséquent}, la production a été interrompue pendant deux heures, le temps que le réparateur remplace la pièce défectueuse. \emph{Afin d'}éviter qu'un tel incident ne se reproduise, le chef d'atelier a décidé de rendre obligatoire la tenue d'un carnet d'entretien, \emph{de sorte que} chaque machine soit désormais vérifiée régulièrement. »

On vérifie : vocabulaire neutre, passé factuel pour les événements, aucun « je », quatre connecteurs logiques, et chaque fait est expliqué ou relié à sa conséquence.
\end{exoresolu}

\begin{exercice}[À vous de rédiger]
Voici des notes prises lors d'une visite de chantier : « Retard livraison ciment. Pluies abondantes en septembre. Route de terre impraticable. Camion bloqué 3 jours. Travaux fondations retardés. Entreprise : commander matériaux à l'avance saison sèche. » Rédigez un paragraphe explicatif de cinq à sept phrases, en respectant le schéma idée $\rightarrow$ explication $\rightarrow$ illustration et en employant au moins quatre connecteurs logiques différents.
\end{exercice}

\begin{corrige}[Exercice : le retard du chantier]
Proposition de corrigé : « Les travaux de fondation du chantier ont accusé un retard important au mois de septembre. Ce retard s'explique \emph{par} le blocage de la livraison de ciment : \emph{en effet}, les pluies abondantes de la saison ont rendu la route de terre impraticable, \emph{si bien que} le camion de livraison est resté bloqué pendant trois jours. \emph{Par conséquent}, les travaux de fondation n'ont pas pu démarrer à la date prévue. \emph{Afin d'}éviter une telle situation à l'avenir, l'entreprise a décidé de commander les matériaux à l'avance pendant la saison sèche, \emph{de manière à} constituer un stock suffisant avant les pluies. » Tout paragraphe respectant la structure en trois temps, l'objectivité et la présence de connecteurs variés est acceptable.
\end{corrige}

\begin{aretenir}[L'essentiel du texte explicatif]
\begin{itemize}
\item Le texte explicatif répond à \emph{pourquoi ?} et \emph{comment ?} ; il fait \emph{comprendre}, il ne se contente pas de raconter.
\item Ses procédés : \cle{définir}, \cle{exemplifier}, \cle{illustrer}, \cle{comparer}, \cle{chiffrer}.
\item Ses connecteurs : cause (\emph{parce que, car, puisque}), conséquence (\emph{donc, c'est pourquoi, par conséquent}), but (\emph{pour, afin de}).
\item Le paragraphe explicatif se construit en trois temps : \textbf{idée principale} $\rightarrow$ \textbf{explication} $\rightarrow$ \textbf{illustration}.
\item L'objectivité impose un vocabulaire neutre, le présent de vérité générale ou le passé factuel, et l'absence du « je » affectif.
\item Règle d'or du rapport : tout constat appelle une explication ou une preuve.
\end{itemize}
\end{aretenir}

Ces outils stylistiques et rhétoriques nous serviront directement dans la suite de la leçon : pour rédiger le corps du rapport, il faudra aussi choisir le bon vocabulaire — c'est pourquoi nous étudierons ensuite la différence entre \emph{lexique commun} et \emph{lexique spécialisé}.

\coursec{Lexique commun et lexique spécialisé (sémantique et lexicologie)}

Dans la section précédente, nous avons vu que le texte explicatif repose sur des procédés précis : définitions, exemples, comparaisons, reformulations. Mais expliquer, c'est d'abord \emph{choisir ses mots}. Or le rapporteur professionnel — technicien, infirmier, comptable — ne parle pas tout à fait la même langue que le quidam : il dispose d'un vocabulaire de métier. Faut-il l'employer dans un rapport ? Jusqu'à quel point ? C'est toute la question de cette section : distinguer le \cle{lexique commun} du \cle{lexique spécialisé}, et apprendre à doser l'un et l'autre selon le destinataire du rapport.

\courssub{Le lexique commun : la langue de tous}

\textbf{Intuition.} Imaginez qu'un incendie se déclare dans l'atelier de votre lycée. Le premier témoin alerté crie : « Il y a le feu ! Une machine brûle ! Un élève est blessé ! » Ces mots — \emph{feu}, \emph{machine}, \emph{blessé} — sont compris instantanément par tout le monde, du directeur au gardien, du professeur à l'élève de seconde. Personne n'a besoin de dictionnaire. C'est cela, le lexique commun : le socle de mots que tous les locuteurs d'une langue partagent, quels que soient leur métier, leur âge ou leur niveau d'études.

\begin{definition}[Lexique commun et lexique spécialisé]
Le \cle{lexique commun} (ou vocabulaire courant) est l'ensemble des mots compris et utilisés par tous les locuteurs d'une langue, indépendamment de leur profession ou de leur spécialité. Exemples : \emph{machine}, \emph{feu}, \emph{blessé}, \emph{couper}, \emph{outil}.

Le \cle{lexique spécialisé} (ou lexique technique) est l'ensemble des termes propres à un domaine d'activité, à un métier, à une science ou à une technique, et qui ne sont pleinement compris que par les spécialistes de ce domaine. Exemples : \emph{tour à métaux} (usinage), \emph{extincteur à poudre} (sécurité incendie), \emph{EPI} (sécurité du travail), \emph{lésion} (médecine).
\end{definition}

Le lexique commun présente trois caractéristiques essentielles :

\begin{itemize}
\item \textbf{Universalité} : il est compris par tous, sans formation particulière ;
\item \textbf{Fréquence} : ce sont les mots les plus employés dans la conversation quotidienne ;
\item \textbf{Polysémie fréquente} : ses mots ont souvent plusieurs sens (le \emph{feu} peut être une flamme, un feu de circulation, le feu d'une cuisinière).
\end{itemize}

\begin{exemplebox}[Lexique commun dans une phrase de rapport]
« Lors de la visite de l'atelier, nous avons vu une \emph{machine} qui \emph{coupe} le bois. Un ouvrier s'est fait \emph{mal} à la main. »

Tous les mots soulignés appartiennent au lexique commun : n'importe quel lecteur comprend la phrase. Mais elle est vague et imprécise — c'est sa limite, comme nous le verrons.
\end{exemplebox}

\courssub{Le lexique spécialisé : la langue des métiers}

\textbf{Observation préalable.} Reprenons la même scène d'atelier, mais écoutons maintenant le chef d'atelier rédiger son rapport d'incident : « Lors de l'opération de \emph{tronçonnage} sur la \emph{scie circulaire}, l'opérateur a subi une \emph{lésion} à la main gauche, faute de port des \emph{EPI}. L'\emph{arrêt d'urgence} a été actionné et l'\emph{extincteur à poudre} tenu prêt. » Le sens est le même, mais chaque terme est précis, sans ambiguïté : on sait exactement quelle machine, quelle opération, quelle blessure, quel matériel de sécurité.

\begin{aretenir}[Pourquoi le lexique spécialisé existe]
Chaque métier crée ses propres mots pour trois raisons :
\begin{enumerate}
\item la \textbf{précision} : \emph{lésion} désigne médicalement une altération des tissus, ce que « blessure » ne précise pas ;
\item l'\textbf{économie} : dire « EPI » évite de répéter « casque, gants, lunettes, chaussures de sécurité » ;
\item la \textbf{communication professionnelle} : les spécialistes d'un même domaine se comprennent immédiatement, sans risque de malentendu.
\end{enumerate}
\end{aretenir}

Le lexique spécialisé se présente sous plusieurs formes :

\begin{itemize}
\item des \textbf{mots savants ou techniques} : \emph{lésion}, \emph{usinage}, \emph{tronçonner}, \emph{coulée} (maçonnerie) ;
\item des \textbf{sigles et abréviations} : EPI, TTC, CV (curriculum vit\ae), BTS ;
\item des \textbf{expressions figées de métier} : \emph{arrêt d'urgence}, \emph{plan de masse}, \emph{traitement de texte} ;
\item des \textbf{mots du lexique commun employés avec un sens technique précis} : c'est le phénomène de la \emph{double appartenance}, que nous étudions ci-dessous.
\end{itemize}

\begin{popfigure}
\begin{center}
\begin{tikzpicture}[scale=1]
\draw[fill=popblueL,opacity=0.75] (-1.6,0) circle (2.4);
\draw[fill=poporangeL,opacity=0.75] (1.6,0) circle (2.4);
\node[font=\bfseries\large,color=popblue] at (-3.1,1.7) {Lexique commun};
\node[font=\bfseries\large,color=poporange] at (3.1,1.7) {Lexique spécialisé};
\node[align=center] at (-2.6,0.3) {\emph{machine}\\ \emph{feu}\\ \emph{blessé}\\ \emph{outil}};
\node[align=center] at (2.6,0.3) {\emph{tronçonner}\\ \emph{lésion}\\ \emph{EPI}\\ \emph{plan de masse}};
\node[align=center,font=\bfseries] at (0,0.3) {\emph{vis}\\ \emph{courant}\\ \emph{champ}\\ \emph{tour}};
\node[align=center,font=\small\itshape] at (0,-1.1) {mots à double\\ appartenance};
\end{tikzpicture}
\end{center}
\legende{Diagramme de Venn : le lexique commun et le lexique spécialisé se recoupent. Certains mots du langage courant (\emph{vis}, \emph{courant}) possèdent en plus un sens technique précis dans un domaine (mécanique, électricité).}
\end{popfigure}

\textbf{Le cas des mots à double appartenance.} Le mot \emph{vis} appartient au lexique commun (tout le monde connaît la petite tige métallique filetée), mais aussi au lexique spécialisé de la mécanique, où il désigne un élément de fixation normalisé, avec des types précis (vis à tête hexagonale, vis sans fin). De même, \emph{courant} est un mot courant (le courant d'une rivière), mais en électricité il désigne rigoureusement la circulation de charges électriques, mesurée en ampères. Dans un rapport technique, c'est toujours le \emph{sens spécialisé} qui est visé : le contexte lève l'ambiguïté.

\courssub{Champ lexical, champ sémantique et lexique spécialisé}

Rappelons deux notions déjà étudiées en classe de seconde et réactivées dans cette unité.

\begin{definition}[Champ lexical et champ sémantique]
Le \cle{champ lexical} d'une notion est l'ensemble des mots (noms, verbes, adjectifs) qui se rapportent à une même idée ou à un même thème. Exemple — champ lexical du feu : \emph{flamme, incendie, brûler, fumée, extincteur, sinistre}.

Le \cle{champ sémantique} d'un mot est l'ensemble des sens que ce mot peut prendre selon les contextes. Exemple — champ sémantique de \emph{feu} : combustion, signal lumineux (feu rouge), passion (le feu de l'amour), autorisation (donner le feu vert).
\end{definition}

Quel est le lien avec notre sujet ? Le lexique spécialisé d'un métier n'est rien d'autre qu'un \textbf{champ lexical organisé autour d'une activité professionnelle}. Le champ lexical de la maçonnerie (\emph{béton, mortier, parpaing, coffrage, couler, ferraillage, taloche}) constitue précisément le lexique spécialisé du maçon. Ainsi, quand vous construisez le champ lexical d'un thème pour rédiger un rapport, vous êtes en train d'assembler le lexique spécialisé dont vous aurez besoin.

\begin{exemplebox}[Du champ lexical au lexique spécialisé]
Thème du rapport : la sécurité incendie dans un atelier.
\begin{itemize}
\item Champ lexical général : \emph{feu, flamme, fumée, brûlure, alarme, évacuation} ;
\item Lexique spécialisé correspondant : \emph{extincteur à poudre, RIA (robinet d'incendie armé), consigne de sécurité, point de rassemblement, déclencheur manuel, classe de feu}.
\end{itemize}
Le passage du premier au second niveau marque le passage du vocabulaire de tout le monde au vocabulaire du professionnel.
\end{exemplebox}

\courssub{Correspondances entre mot commun et terme spécialisé}

Pour bien sentir la différence de niveau de langue professionnelle, voici un tableau de correspondances utiles dans les rapports techniques de votre niveau.

\begin{popfigure}
\begin{center}
\begin{tabularx}{\linewidth}{|>{\raggedright\arraybackslash}X|>{\raggedright\arraybackslash}X|}
\hline
\rowcolor{popblueL}
\textbf{Mot du lexique commun} & \textbf{Terme spécialisé} \\
\hline
couper (métal, bois) & tronçonner, scier, usiner \\
\hline
blessure & lésion \\
\hline
machine à tourner le métal & tour à métaux \\
\hline
mal au dos professionnel & lombalgie professionnelle \\
\hline
papier qui résume une visite & compte rendu de visite \\
\hline
mettre le feu sans le vouloir & départ de feu accidentel \\
\hline
\end{tabularx}
\end{center}
\legende{Correspondances entre lexique commun et lexique spécialisé : à chaque formulation courante correspond un terme technique plus précis, adapté au rapport professionnel.}
\end{popfigure}

\courssub{Adapter le vocabulaire au destinataire : le critère décisif}

Voici le cœur pratique de cette leçon. Un même fait peut être rapporté de deux façons :

\begin{itemize}
\item Pour un collègue technicien : « La pompe centrifuge a été désaccouplée pour contrôle du joint d'étanchéité. »
\item Pour le proviseur, non technicien : « La pompe (machine qui fait circuler l'eau) a été séparée de son moteur pour vérifier la pièce qui empêche les fuites. »
\end{itemize}

Ni l'une ni l'autre version n'est « meilleure » en soi : tout dépend de \emph{qui lit}.

\begin{methode}[Choisir le bon niveau de vocabulaire selon le destinataire]
\begin{enumerate}
\item \textbf{Identifier le destinataire} du rapport : est-ce un spécialiste du domaine (chef d'atelier, ingénieur, médecin) ou un non-spécialiste (proviseur, parent d'élève, autorité administrative, client) ?
\item \textbf{Destinataire spécialiste} : employer librement le lexique technique du domaine. Il garantit la précision et fait gagner du temps. Inutile de définir les termes courants du métier.
\item \textbf{Destinataire non spécialiste (profane)} : employer le lexique technique \emph{quand il est indispensable à la précision}, mais \textbf{définir chaque terme difficile} entre parenthèses, dans une note ou dans un glossaire en fin de rapport.
\item \textbf{Vérifier l'équilibre} : relire en se demandant, phrase par phrase : « Mon lecteur comprendra-t-il ce mot sans explication ? » Si la réponse est non, définir ou remplacer.
\item \textbf{Rester cohérent} : ne pas alterner sans raison « blessure » puis « lésion » pour désigner la même chose ; une fois le terme choisi et défini, s'y tenir.
\end{enumerate}
\end{methode}

\begin{exemplebox}[Un dosage réussi : le rapport de visite à la SONARA]
Dans un rapport de visite pédagogique à la raffinerie SONARA de Limbé, destiné au proviseur (non technicien), un élève de Première technique a employé 40 termes relevant du domaine pétrochimique. Sur ces 40 termes, \textbf{12 termes spécialisés ont été définis entre parenthèses} dès leur première occurrence : « le brut (pétrole non raffiné) », « la distillation fractionnée (séparation des produits selon leur température d'ébullition) », « le torchage (brûlage des gaz excédentaires) ». Les 28 autres termes (\emph{usine, réservoir, pipeline, visite, atelier}) appartenaient au lexique commun ou étaient transparents. Résultat : un rapport précis \emph{et} lisible par tous les membres de l'administration.
\end{exemplebox}

\begin{attention}[Erreur fréquente : abuser du jargon pour « faire savant »]
Beaucoup d'élèves croient qu'un bon rapport est un rapport bourré de termes techniques. C'est faux. Empiler les sigles et les mots savants que le destinataire ne comprend pas, c'est \emph{échouer à communiquer} — or la fonction première du rapport est d'informer clairement. Trois pièges à éviter :
\begin{itemize}
\item utiliser un sigle sans le développer à sa première occurrence (écrire d'abord « Équipement de Protection Individuelle (EPI) », puis « EPI » ensuite) ;
\item employer un terme technique alors qu'un mot commun suffirait (« contusionner » pour « se faire un bleu » dans un récit simple) ;
\item définir des mots que tout le monde connaît, ce qui alourdit le texte et peut sembler méprisant pour le lecteur.
\end{itemize}
Le bon rapporteur n'est pas celui qui en sait le plus de mots, mais celui qui se fait comprendre sans rien sacrifier à la précision.
\end{attention}

\begin{savaistu}[Le sigle EPI, un mot de métier devenu obligatoire]
« EPI » signifie \textbf{Équipement de Protection Individuelle} : casque, gants, lunettes, chaussures de sécurité, gilet, bouchons d'oreilles... Ce sigle appartient au lexique spécialisé de la sécurité du travail. Dans tout rapport d'incident ou d'accident survenu en atelier — situation que vous rencontrerez dans votre vie professionnelle — la mention des EPI (portés ou non portés par la victime) est \emph{exigible} : c'est un élément juridique qui détermine les responsabilités. Un rapport d'incident qui oublie de préciser les EPI est un rapport incomplet.
\end{savaistu}

\courssub{Atelier : construire le lexique spécialisé de trois métiers}

\begin{experience}[Atelier de construction de lexiques spécialisés]
\textbf{Matériel} : affiches de sécurité de l'atelier du lycée, catalogues d'outillage, fiches métier, dictionnaire, documents fournis par les enseignants de matières techniques.

\textbf{Protocole} (travail en groupes de quatre, une séance) :
\begin{enumerate}
\item Chaque groupe choisit un métier : \emph{menuiserie}, \emph{maçonnerie} ou \emph{secrétariat bureautique}.
\item À partir des affiches et documents, relever 10 termes spécialisés du métier (outils, gestes, matériaux, documents, sigles).
\item Pour chaque terme, rédiger une définition courte en lexique commun, comme si l'on s'adressait à un profane.
\item Classer les termes en sous-catégories : outils / gestes professionnels / matériaux ou documents / sécurité.
\item Présenter le lexique à la classe ; les autres groupes complètent.
\end{enumerate}

\textbf{Observations attendues} (exemples de résultats) :
\begin{itemize}
\item \textbf{Menuiserie} : \emph{rabot, varlope, mortaise, tenon, dégauchir, abouter, panneau de particules, chant, assemblage, EPI} ;
\item \textbf{Maçonnerie} : \emph{parpaing, mortier, béton armé, ferraillage, coffrage, couler, taloche, niveau à bulle, chaînage, plan de masse} ;
\item \textbf{Secrétariat bureautique} : \emph{traitement de texte, tableur, classement chronologique, note de service, agenda partagé, numérotation, archivage, courrier arrivée, CV, TTC (Toutes Taxes Comprises)}.
\end{itemize}

\textbf{Interprétation} : chaque métier possède un champ lexical propre, organisé autour de ses outils, de ses gestes et de ses objets. Maîtriser ce lexique, c'est déjà maîtriser une partie du métier — et c'est la condition pour rédiger un rapport crédible dans ce domaine.
\end{experience}

\courssub{Application : réécrire un passage de rapport}

\begin{exoresolu}[Du lexique commun au lexique spécialisé]
\textbf{Énoncé.} Voici un extrait du rapport d'un élève sur un incident survenu à l'atelier de menuiserie, rédigé en lexique commun. Réécrivez ce passage en remplaçant les mots soulignés par les termes spécialisés appropriés, puis indiquez quels termes vous définiriez pour un lecteur profane.

« Pendant le travail, un élève \emph{coupait} une planche avec la \emph{machine à scier}. Il ne portait pas ses \emph{affaires de protection}. La lame a touché sa main et lui a fait une \emph{blessure}. Le professeur a appuyé sur le \emph{bouton rouge} pour arrêter la machine et a prévenu l'infirmière. »

\tcblower
\textbf{Solution détaillée.}

Version réécrite en lexique spécialisé :

« Pendant les travaux pratiques, un élève \emph{tronçonnait} une planche sur la \emph{scie circulaire}. Il ne portait pas ses \emph{EPI} (Équipements de Protection Individuelle). La lame a atteint sa main gauche et provoqué une \emph{lésion}. Le professeur a actionné l'\emph{arrêt d'urgence} et a alerté l'infirmière de l'établissement. »

Correspondances opérées :
\begin{itemize}
\item \emph{coupait} $\rightarrow$ \emph{tronçonnait} (geste technique précis) ;
\item \emph{machine à scier} $\rightarrow$ \emph{scie circulaire} (désignation exacte de la machine) ;
\item \emph{affaires de protection} $\rightarrow$ \emph{EPI} (terme normalisé de la sécurité) ;
\item \emph{blessure} $\rightarrow$ \emph{lésion} (terme médical) ;
\item \emph{bouton rouge} $\rightarrow$ \emph{arrêt d'urgence} (désignation fonctionnelle du dispositif).
\end{itemize}

Pour un lecteur profane (par exemple le proviseur), on définirait entre parenthèses : \emph{EPI} (gants, lunettes, vêtements de protection) et éventuellement \emph{arrêt d'urgence} (bouton qui stoppe immédiatement la machine). Les termes \emph{tronçonner}, \emph{scie circulaire} et \emph{lésion} restent compréhensibles par le contexte.
\end{exoresolu}

\begin{exercice}[À vous de jouer]
Rédigez un court paragraphe de rapport (5 à 7 lignes) sur la visite d'un chantier de construction, destiné au proviseur de votre lycée. Contraintes : employer au moins cinq termes du lexique spécialisé de la maçonnerie, en définir deux entre parenthèses, et utiliser le lexique commun pour tout le reste. Soulignez les termes spécialisés employés.
\end{exercice}

\begin{corrige}[Exercice : éléments de correction]
Exemple de production attendue :

« Le mardi 14 mars, notre classe a visité le chantier du nouveau bloc administratif. Les maçons achevaient le \emph{coffrage} (structure provisoire qui donne sa forme au béton) du deuxième niveau. Nous avons observé le \emph{ferraillage} des poteaux avant la \emph{coulée} du béton. Le chef de chantier nous a expliqué que le \emph{béton armé} (béton renforcé de barres d'acier) garantit la solidité de l'ouvrage. Tous les ouvriers portaient leurs EPI. La visite s'est déroulée sans incident, dans le respect des consignes de sécurité. »

Barème indicatif : 5 termes spécialisés correctement employés (5 points), 2 définitions entre parenthèses adaptées à un profane (2 points), correction de la langue et cohérence du récit (3 points).
\end{corrige}

\begin{aretenir}[L'essentiel de la section]
\begin{itemize}
\item Le \cle{lexique commun} est compris de tous (\emph{machine, feu, blessé}) ; le \cle{lexique spécialisé} est propre à un métier ou à une science (\emph{tour à métaux, extincteur à poudre, EPI}).
\item Certains mots ont une double appartenance (\emph{vis, courant}) : dans un rapport technique, c'est le sens spécialisé qui compte.
\item Le lexique spécialisé d'un métier est un champ lexical organisé autour de ses outils, gestes et objets.
\item Dans un rapport : employer le lexique technique du domaine pour la précision, mais \textbf{définir les termes difficiles} si le destinataire est profane.
\item Le jargon excessif est une faute de communication : le bon niveau de vocabulaire est celui du destinataire, pas celui de l'auteur.
\end{itemize}
\end{aretenir}

Ce travail sur le vocabulaire trouvera son application immédiate dans la suite de la leçon : avant de rédiger un rapport, il faut savoir \emph{lire} les documents qui le nourrissent — notices, consignes, affiches de sécurité. C'est l'objet de la section suivante, consacrée à la réception des textes fonctionnels et des affiches d'indications sécuritaires.

\coursec{Lire les textes fonctionnels et les affiches de sécurité (réception)}

Dans la section précédente, nous avons distingué le \cle{lexique commun} du \cle{lexique spécialisé} : un même objet peut être nommé différemment selon que l'on parle dans la vie courante ou dans un milieu professionnel. Mais le vocabulaire ne suffit pas à communiquer : il faut encore choisir le bon \emph{type de texte}. Or, dans la vie professionnelle du technicien — à l'atelier, au chantier, au laboratoire —, on ne lit pas des romans ni des poèmes : on lit des \cle{notices}, des \cle{consignes}, des \cle{règlements}, des \cle{affiches}, des \cle{formulaires}. Ces textes ont un point commun : ils ne cherchent pas à plaire, ils cherchent à \emph{faire agir} ou à \emph{informer pour agir}. Ce sont des textes fonctionnels. Savoir les lire vite et bien est une compétence de sécurité, parfois une question de vie ou de mort. C'est l'objet de cette section, qui nous donnera aussi un outil précieux pour la rédaction du rapport : citer une consigne comme preuve.

\courssub{Qu'est-ce qu'un texte fonctionnel ?}

\begin{definition}[Le texte fonctionnel]
Un \cle{texte fonctionnel} (ou texte utilitaire) est un texte dont la fonction principale est de \emph{guider une action} ou d'\emph{informer le lecteur pour qu'il agisse correctement}. Il ne vise ni le plaisir esthétique ni l'émotion : il vise l'\emph{efficacité pratique}.
\end{definition}

\begin{definition}[Le texte iconique]
Un \cle{texte iconique} est une image porteuse d'information : affiche, pictogramme, panneau, schéma, logo. Il communique un message \emph{sans phrase} ou avec très peu de mots. Dans les documents professionnels, le texte iconique accompagne souvent le texte verbal : on parle alors de lecture \emph{croisée} image-texte.
\end{definition}

L'intuition est simple : quand tu arrives dans un atelier pour la première fois, personne ne te fait un discours. On te remet une \emph{notice}, on te montre une \emph{affiche} au mur, on te fait signer un \emph{règlement intérieur}. Ces trois documents sont des textes fonctionnels. Leur lecteur est pressé, parfois peu lettré, parfois étranger : le message doit donc être immédiat, clair et sans ambiguïté.

\begin{aretenir}[Exemples de textes fonctionnels]
\begin{itemize}
\item la \cle{notice} d'utilisation d'une machine (meuleuse, perceuse, cuisinière) ;
\item la \cle{consigne} affichée : « Porter les lunettes de protection » ;
\item le \cle{règlement} intérieur de l'atelier ou de l'entreprise ;
\item l'\cle{affiche} de sécurité ou d'information ;
\item le \cle{formulaire} à remplir (fiche d'inscription, fiche d'incident, bon de commande).
\end{itemize}
\end{aretenir}

\courssub{Les caractéristiques du texte fonctionnel}

Observons une consigne réelle d'atelier : « AVANT DE DÉMARRER LA MACHINE, METTRE LES LUNETTES ET VÉRIFIER LE CARTER DE PROTECTION. » Que remarque-t-on ?

\begin{enumerate}
\item La \cle{brièveté} : une ou deux phrases, aucun développement. Le lecteur doit comprendre en quelques secondes.
\item L'emploi de l'\cle{infinitif} (« mettre », « vérifier ») ou de l'\cle{impératif} (« Mettez », « Ne touchez pas ») : ce sont les modes de l'ordre et de la prescription.
\item Des \cle{phrases courtes}, sans subordonnées inutiles, au vocabulaire précis (lexique spécialisé : « carter de protection »).
\item Une \cle{typographie hiérarchisée} : majuscules, gras, tailles différentes pour distinguer le titre, l'avertissement, le détail.
\item La présence fréquente de \cle{pictogrammes} qui renforcent ou remplacent le message verbal.
\end{enumerate}

\begin{attention}[Piège fréquent]
L'erreur la plus courante, à l'examen comme dans la vie, est de \emph{décrire} l'affiche (« il y a un rond rouge avec une cigarette ») sans \emph{interpréter} la relation entre l'image et le texte. Or les deux se complètent : le pictogramme attire l'œil et donne le sens général, la consigne écrite précise la conduite à tenir. Une bonne lecture dit toujours \emph{ce que l'image ajoute au texte} et \emph{ce que le texte ajoute à l'image}.
\end{attention}

\courssub{Les affiches d'indications sécuritaires : le code forme-couleur}

Les affiches de sécurité obéissent à un code international très simple, fondé sur la \emph{forme} et la \emph{couleur}. Ce code permet de comprendre le message avant même de lire les mots : c'est précisément ce qui le rend efficace auprès de tout public, même non francophone.

\begin{exemplebox}[Le code couleur international]
\begin{itemize}
\item \textbf{Rouge, rond barré} : \cle{interdiction} (ex. : « Défense de fumer »).
\item \textbf{Bleu, rond} : \cle{obligation} (ex. : « Port des lunettes obligatoire »).
\item \textbf{Jaune, triangle} : \cle{avertissement}, danger (ex. : « Attention, pièces en rotation »).
\item \textbf{Vert, rectangle ou carré} : \cle{secours et évacuation} (ex. : « Sortie de secours », « Poste de secours »).
\end{itemize}
\end{exemplebox}

\begin{popfigure}
\begin{center}
\begin{tikzpicture}[scale=0.9]
% Interdiction : cercle barré rouge
\begin{scope}[shift={(0,0)}]
\draw[line width=2.5pt, poppink] (0,0) circle (0.8);
\draw[line width=2.5pt, poppink] (-0.57,0.57) -- (0.57,-0.57);
\node at (0,-1.3) {\small \textbf{Interdiction}};
\node at (0,-1.7) {\small rond rouge barré};
\end{scope}
% Obligation : disque bleu
\begin{scope}[shift={(3.5,0)}]
\fill[popblue] (0,0) circle (0.8);
\node[white] at (0,0) {\small \textbf{!}};
\node at (0,-1.3) {\small \textbf{Obligation}};
\node at (0,-1.7) {\small disque bleu};
\end{scope}
% Danger : triangle jaune
\begin{scope}[shift={(7,0)}]
\fill[popgold, draw=popdark, line width=1.5pt] (-0.9,-0.7) -- (0.9,-0.7) -- (0,0.85) -- cycle;
\node at (0,-0.25) {\small \textbf{!}};
\node at (0,-1.3) {\small \textbf{Danger}};
\node at (0,-1.7) {\small triangle jaune};
\end{scope}
% Évacuation : rectangle vert
\begin{scope}[shift={(10.5,0)}]
\fill[popgreen] (-0.9,-0.6) rectangle (0.9,0.6);
\draw[white, line width=2pt, -{Stealth}] (-0.5,0) -- (0.5,0);
\node at (0,-1.3) {\small \textbf{Évacuation}};
\node at (0,-1.7) {\small rectangle vert};
\end{scope}
\end{tikzpicture}
\end{center}
\legende{Les quatre familles de signalétique de sécurité : la forme et la couleur annoncent le type de message avant toute lecture.}
\end{popfigure}

\begin{savaistu}[Une norme mondiale]
Les pictogrammes de sécurité sont \cle{normalisés} par la norme internationale \textbf{ISO 7010}. Cela signifie qu'un panneau « sortie de secours » ou « port du casque obligatoire » est \emph{identique} dans une entreprise de Douala, de Yaoundé, de Lyon ou de Tokyo. Un technicien camerounais peut donc travailler partout dans le monde en comprenant immédiatement la signalétique : c'est une langue universelle sans mots.
\end{savaistu}

\courssub{Méthode : lire une affiche de sécurité en quatre points}

\begin{methode}[Lire une affiche de sécurité]
\begin{enumerate}
\item \textbf{La forme} : rond, triangle ou rectangle ? La forme indique déjà la nature du message (interdiction, obligation, danger, secours).
\item \textbf{La couleur} : rouge, bleu, jaune ou vert ? Elle confirme le type de prescription.
\item \textbf{Le pictogramme} : que représente le dessin (cigarette, lunettes, flamme, flèche, bonhomme qui court) ? Il donne le contenu concret.
\item \textbf{Le message verbal} : lire la consigne écrite, noter son mode (infinitif, impératif) et son vocabulaire, puis vérifier qu'image et texte disent la même chose.
\end{enumerate}
\end{methode}

\begin{popfigure}
\begin{center}
\begin{tikzpicture}[node distance=0.6cm and 1.1cm,
box/.style={draw=popblue, fill=popblueL, rounded corners, thick, text width=2.6cm, align=center, minimum height=1cm, font=\small}]
\node[box, align=center] (picto){\textbf{1. Pictogramme}\\ que voit-on ?};
\node[box, right=of picto, align=center] (consigne){\textbf{2. Consigne}\\ que lit-on ?};
\node[box, right=of consigne, align=center] (couleur){\textbf{3. Forme et couleur}\\ quel type ?};
\node[box, right=of couleur, align=center] (lieu){\textbf{4. Emplacement}\\ où est-elle ?};
\draw[-{Stealth}, thick, poporange] (picto) -- (consigne);
\draw[-{Stealth}, thick, poporange] (consigne) -- (couleur);
\draw[-{Stealth}, thick, poporange] (couleur) -- (lieu);
\node[below=0.5cm of consigne, text width=6cm, align=center, font=\small\itshape] {Lecture croisée : le sens complet naît de la réunion des quatre éléments.};
\end{tikzpicture}
\end{center}
\legende{Schéma de lecture d'une affiche de sécurité : du pictogramme vers la consigne, la couleur, puis l'emplacement.}
\end{popfigure}

\courssub{Analyse guidée : trois affiches de l'atelier du lycée}

Appliquons la méthode à trois affiches réelles de l'atelier.

\textbf{Affiche 1 : « Port des lunettes obligatoire ».} Forme : disque. Couleur : bleu. Pictogramme : un visage stylisé portant des lunettes. Message verbal : syntagme nominal exprimant une obligation (l'adjectif « obligatoire » vaut un impératif). Interprétation : le bleu impose une conduite ; le pictogramme montre l'équipement concerné ; l'affiche est placée à l'entrée de l'atelier, là où la décision de se protéger doit être prise.

\textbf{Affiche 2 : « Sortie de secours ».} Forme : rectangle. Couleur : vert. Pictogramme : un bonhomme qui court vers une porte, avec une flèche. Interprétation : le vert n'interdit ni n'oblige, il \emph{indique un secours} ; la flèche oriente le déplacement en cas d'incendie.

\textbf{Affiche 3 : « Défense de fumer ».} Forme : rond barré. Couleur : rouge. Pictogramme : une cigarette allumée. Interprétation : le rouge barré interdit ; la barre oblique matérialise visuellement la négation, même pour un lecteur qui ne lit pas le français.

\begin{exoresolu}[Lecture croisée d'une affiche]
Énoncé. On observe dans l'atelier un panneau triangulaire jaune représentant un engrenage, avec la mention « Pièces en rotation — ne pas approcher les mains ». Analyse cette affiche en quatre points et dégage son sens complet.
\tcblower
Solution détaillée.
\begin{enumerate}
\item \emph{Forme} : triangle, donc message d'avertissement.
\item \emph{Couleur} : jaune, qui confirme la présence d'un danger.
\item \emph{Pictogramme} : un engrenage, c'est-à-dire des pièces mécaniques en mouvement qui peuvent happer les doigts ou les vêtements.
\item \emph{Message verbal} : « ne pas approcher les mains », infinitif négatif, mode de la consigne.
\end{enumerate}
Sens complet : l'image identifie la source du danger (la machine), la couleur et la forme signalent qu'il s'agit d'un avertissement, et le texte prescrit la conduite exacte (garder les mains à distance). Image et texte se complètent : séparément, chacun serait moins précis.
\end{exoresolu}

\courssub{Le lien avec le rapport : citer une consigne comme preuve}

Pourquoi étudier les affiches dans une leçon consacrée au \cle{rapport} ? Parce qu'un rapport d'incident ou d'accident doit s'appuyer sur des \emph{preuves}, et les consignes affichées en font partie. Si un élève se blesse à la meuleuse, le rapport devra préciser si la consigne existait, où elle était affichée et si elle a été respectée.

Pour citer correctement une affiche dans un rapport, on indique : sa \emph{nature} (panneau d'obligation), son \emph{emplacement}, sa \emph{formulation exacte} entre guillemets.

\begin{exemplebox}[Citation d'une consigne dans un rapport d'incident]
« Le panneau d'obligation (disque bleu) fixé à l'entrée de l'atelier de mécanique prescrit : "Port des lunettes obligatoire". Or, au moment de l'incident, l'élève X. ne portait pas ses lunettes de protection. »
\end{exemplebox}

\begin{attention}[Citer, c'est être exact]
Dans un rapport, on ne paraphrase pas une consigne : on la cite \emph{mot pour mot}, entre guillemets, avec son emplacement. Écrire « il était marqué quelque part qu'il fallait des lunettes » est imprécis et sans valeur probante. L'exactitude de la citation fait la crédibilité du rapport.
\end{attention}

\courssub{Mini-TP : concevoir une affiche sécuritaire pour l'atelier}

\begin{exercice}[Conception d'une affiche de sécurité]
Par groupes de trois, concevez une affiche de sécurité pour l'atelier de votre lycée (exemples : port des gants, interdiction des vêtements amples près des machines, lavage des mains après manipulation, rassemblement en cas d'incendie).
\begin{enumerate}
\item Choisissez la \emph{forme} et la \emph{couleur} adaptées au type de message (interdiction, obligation, danger, évacuation).
\item Dessinez un \emph{pictogramme} simple et lisible de loin.
\item Rédigez la \emph{consigne} à l'infinitif, en une phrase courte, avec le lexique spécialisé exact.
\item Présentez votre affiche à la classe en justifiant chaque choix selon la méthode en quatre points.
\end{enumerate}
\end{exercice}

\begin{corrige}[Mini-TP : exemple de production attendue]
Affiche « Port des gants obligatoire » : disque bleu (obligation) ; pictogramme : une main gantée ; consigne rédigée : « Porter les gants de protection avant toute manipulation des pièces coupantes » (infinitif, phrase courte, lexique spécialisé). Justification : le bleu impose, le pictogramme montre l'équipement, l'infinitif prescrit sans ambiguïté ; l'affiche sera placée à l'entrée de l'atelier, point de décision.
\end{corrige}

\begin{aretenir}[L'essentiel de la section]
Un texte fonctionnel informe pour faire agir : il est bref, rédigé à l'infinitif ou à l'impératif, hiérarchisé typographiquement, souvent accompagné d'un pictogramme. Les affiches de sécurité obéissent à un code international (rouge = interdiction, bleu = obligation, jaune = danger, vert = secours) normalisé par l'ISO 7010. On les lit en quatre points : forme, couleur, pictogramme, message verbal. Enfin, dans un rapport, une consigne se cite exactement, entre guillemets, avec son emplacement : elle constitue une preuve.
\end{aretenir}

\coursec{Rédaction du rapport : introduction, corps, conclusion}

Dans la section précédente, nous avons appris à \emph{lire} les textes fonctionnels (notes de service, comptes rendus, affiches de sécurité) : nous savons désormais repérer leur organisation, leur vocabulaire et leur ton objectif. Il est temps de passer de la réception à la \cle{production} : vous allez maintenant \emph{rédiger} vous-même un rapport complet. C'est l'aboutissement de toute l'unité : les notes prises pendant une visite ou un stage, le plan établi, le lexique spécialisé repéré et les procédés du texte explicatif vont enfin se combiner dans un texte fini, prêt à être remis à un destinataire.

\courssub{Vue d'ensemble : le processus de rédaction}

On ne rédige jamais un rapport directement « au propre ». Comme un technicien qui suit une procédure de fabrication, le rédacteur suit des étapes ordonnées. Sauter une étape, c'est risquer un texte désorganisé, truffé de ratures et de fautes.

\begin{popfigure}
\begin{center}
\begin{tikzpicture}[font=\small, >=Stealth,
  etape/.style={draw=popblue, fill=popblueL, rounded corners=2pt, minimum height=0.9cm, align=center, text width=1.9cm}]
  \node[etape, align=center] (n1) at (0,0){\textbf{1. Notes}\\(prises sur le terrain)};
  \node[etape, right=0.55cm of n1, align=center] (n2){\textbf{2. Plan}\\(notes triées et ordonnées)};
  \node[etape, right=0.55cm of n2, align=center] (n3){\textbf{3. Brouillon}\\(rédaction libre)};
  \node[etape, right=0.55cm of n3, align=center] (n4){\textbf{4. Rédaction}\\(version structurée)};
  \node[etape, right=0.55cm of n4, align=center] (n5){\textbf{5. Relecture}\\(correction)};
  \node[etape, right=0.55cm of n5, fill=popgreenL, draw=popgreen, align=center] (n6){\textbf{6. Version}\\\textbf{finale}};
  \draw[->, thick, popdark] (n1) -- (n2);
  \draw[->, thick, popdark] (n2) -- (n3);
  \draw[->, thick, popdark] (n3) -- (n4);
  \draw[->, thick, popdark] (n4) -- (n5);
  \draw[->, thick, popdark] (n5) -- (n6);
  \draw[->, thick, poporange, dashed] (n5.south) to[bend right=25] node[below, font=\scriptsize]{fautes détectées : on reprend} (n3.south);
\end{tikzpicture}
\end{center}
\legende{Frise du processus de rédaction d'un rapport : six étapes successives, avec possibilité de retour au brouillon après la relecture.}
\end{popfigure}

\begin{aretenir}[Principe directeur]
Un rapport se construit en \cle{trois parties} : une \textbf{introduction} (contexte, circonstances, objectif, plan), un \textbf{corps} (constats objectifs articulés par des connecteurs) et une \textbf{conclusion} (bilan, suggestions, formule de courtoisie). Chaque partie a une fonction précise qu'il ne faut ni confondre ni omettre.
\end{aretenir}

\courssub{L'introduction : situer et annoncer}

\textbf{Intuition.} Imaginez que le proviseur reçoive un texte commençant brutalement par : « L'atelier de soudure compte 12 postes. » Sa première réaction sera : \emph{de quoi parle-t-on ? qui écrit ? pourquoi ?} L'introduction répond à ces questions avant même que le lecteur ne les pose. Elle est la porte d'entrée du rapport.

\begin{definition}[L'introduction du rapport]
L'\cle{introduction} est la partie liminaire du rapport. Elle présente le \textbf{contexte} (qui rédige, quand, où), les \textbf{circonstances} de la mission (qui a donné l'ordre, à quelle occasion), l'\textbf{objectif} du rapport (ce qu'on est allé observer ou vérifier) et l'\textbf{annonce du plan} (comment le corps sera organisé).
\end{definition}

\begin{methode}[Rédiger l'introduction en quatre phrases]
\begin{enumerate}
\item \textbf{Phrase de contexte} : qui êtes-vous, quand et où la mission a-t-elle eu lieu ? Exemple : « Le mardi 14 janvier, les élèves de la classe de Première technique se sont rendus à l'entreprise SOCATER de Douala. »
\item \textbf{Phrase de circonstances} : qui a ordonné ou autorisé la mission ? Exemple : « Cette visite s'inscrit dans le cadre des activités pédagogiques prévues par le programme de Français. »
\item \textbf{Phrase d'objectif} : que deviez-vous observer ou vérifier ? Exemple : « Elle avait pour objectif d'observer l'organisation de la production et l'application des consignes de sécurité. »
\item \textbf{Phrase d'annonce du plan} : comment le rapport est-il organisé ? Exemple : « Après avoir présenté l'entreprise, nous décrirons le déroulement de la visite, puis nous dégagerons les enseignements tirés de cette sortie. »
\end{enumerate}
\end{methode}

\textbf{Les formules d'ouverture consacrées.} L'administration et le monde professionnel utilisent des formules stéréotypées qu'il faut connaître et employer à bon escient :

\begin{center}
\begin{tabularx}{\linewidth}{|l|X|}\hline
\rowcolor{popblueL} \textbf{Formule} & \textbf{Emploi} \\ \hline
« À la suite de... » & La mission découle d'un événement : « À la suite de l'accident survenu le 3 mars, j'ai été chargé de... » \\ \hline
« Conformément à votre instruction du... » & La mission répond à un ordre précis et daté du supérieur. \\ \hline
« J'ai l'honneur de... » & Ouverture solennelle d'un rapport adressé à une autorité hiérarchique. \\ \hline
« Dans le cadre de... » & La mission s'insère dans un programme, un projet, une formation. \\ \hline
\end{tabularx}
\end{center}

\begin{exemplebox}[Introduction complète d'un rapport de visite]
« Conformément à votre instruction du 10 janvier, j'ai effectué, le mardi 14 janvier, une visite d'étude à l'entreprise SOCATER de Douala, accompagné des élèves de la classe de Première technique. Cette visite avait pour objectif d'observer l'organisation de la chaîne de production et le respect des consignes de sécurité affichées dans les ateliers. Après une brève présentation de l'entreprise, le présent rapport décrit le déroulement de la visite, expose les constats effectués et formule quelques suggestions. »
\end{exemplebox}

\courssub{Le corps : rapporter des faits avec objectivité}

\textbf{Intuition.} Le corps est le cœur du rapport : c'est là que se trouve l'information attendue par le destinataire. Mais attention : un rapport n'est ni un récit de voyage ni un journal intime. Le rédacteur doit se faire \emph{témoin impartial} : il rapporte ce qu'il a vu, mesuré, entendu — rien de plus.

\begin{definition}[Le corps du rapport]
Le \cle{corps} est la partie développée du rapport. Il présente les \textbf{constats} organisés en paragraphes, chacun correspondant à une partie du plan. Les paragraphes sont articulés par des \textbf{connecteurs chronologiques ou logiques} : \emph{d'abord, ensuite, puis, en outre, par ailleurs, enfin}.
\end{definition}

\textbf{Les trois garanties de l'objectivité.}

\begin{enumerate}
\item \textbf{Des faits et des chiffres} : préférer « l'atelier compte 12 postes de soudure, dont 3 hors service » à « l'atelier est mal équipé ». Le chiffre est vérifiable, le jugement ne l'est pas.
\item \textbf{Des témoignages cités} : rapporter les paroles recueillies au discours direct ou indirect, en nommant la source par sa fonction : « Selon le chef d'atelier, la machine n° 2 est arrêtée depuis trois semaines. »
\item \textbf{Des références aux documents} : s'appuyer sur ce qui est écrit et visible : « Conformément à l'affiche de sécurité placée à l'entrée, le port du casque est obligatoire ; or 4 ouvriers sur 15 n'en portaient pas au moment de la visite. »
\end{enumerate}

\begin{attention}[Erreur fréquente : la subjectivité]
N'introduisez \textbf{jamais} d'opinions personnelles ni de jugements de valeur dans le corps d'un rapport. La phrase « \emph{je pense que le chef d'atelier est négligent} » est inacceptable : elle juge une personne au lieu de décrire des faits. Écrivez plutôt : « \emph{le registre de maintenance n'a pas été rempli depuis le 2 février} ». Le lecteur tirera lui-même ses conclusions. Règle d'or : \cle{rapporter des faits, pas des jugements}.
\end{attention}

\begin{exemplebox}[Paragraphe de corps bien articulé]
« D'abord, le groupe a été reçu par le directeur technique, qui a présenté l'historique de l'entreprise. Ensuite, la visite s'est poursuivie dans l'atelier de découpe, où fonctionnent 8 machines commandées numériquement. En outre, les élèves ont relevé les consignes portées sur les affiches de sécurité : port obligatoire des lunettes, interdiction de fumer, consignes d'évacuation. Enfin, un échange avec le chef d'atelier a permis de préciser les rythmes de production : 150 pièces par jour en moyenne. »
\end{exemplebox}

\courssub{La conclusion : bilan, suggestions et courtoisie}

\textbf{Intuition.} Un rapport qui s'arrête brutalement après le dernier constat laisse le destinataire sur sa faim : \emph{que faut-il retenir ? que faut-il faire ?} La conclusion répond à ces deux questions, puis prend congé poliment.

\begin{definition}[La conclusion du rapport]
La \cle{conclusion} comprend trois éléments : un \textbf{bilan synthétique} (l'essentiel des constats en deux ou trois phrases), des \textbf{suggestions ou recommandations} adressées au destinataire, et une \textbf{formule de courtoisie administrative} adaptée au destinataire.
\end{definition}

\begin{methode}[Conclure : bilan + suggestions réalistes et polies]
\begin{enumerate}
\item Rédiger le \textbf{bilan} en reprenant les grandes lignes du corps, sans ajouter de fait nouveau : « Cette visite a montré une entreprise bien organisée, dont l'outil de production souffre toutefois de quelques pannes. »
\item Formuler les \textbf{suggestions} au conditionnel ou avec des tournures atténuées : « Nous suggérons que... », « Il serait souhaitable de... », « Nous recommandons la remise en service de... ». Les suggestions doivent être \textbf{réalistes} (ni irréalisables, ni hors de votre mission).
\item Terminer par la \textbf{formule de courtoisie} : « Nous vous prions d'agréer, Monsieur le Proviseur, l'expression de notre respectueuse considération. »
\end{enumerate}
\end{methode}

\begin{savaistu}[« J'ai l'honneur de... » : une formule hiérarchique]
La formule finale « \emph{J'ai l'honneur de vous présenter...} » est réservée aux rapports adressés à une \textbf{autorité hiérarchique} (proviseur, directeur, ministre). Elle ne s'emploie \textbf{jamais} entre camarades ni entre collègues de même rang : dans ce cas, on préfère « Cordialement » ou « Bien à vous ». Employer « J'ai l'honneur » dans un message à un ami serait aussi déplacé que de le vouvoyer solennellement dans la cour de récréation !
\end{savaistu}

\courssub{Correction guidée d'un brouillon}

Avant la version finale, tout brouillon doit subir une relecture méthodique sur trois plans : la \textbf{cohérence} (les idées se suivent-elles ? le plan annoncé est-il respecté ?), l'\textbf{objectivité} (y a-t-il des jugements personnels ?) et la \textbf{structure} (introduction, corps et conclusion sont-ils présents et complets ?).

\begin{exoresolu}[Corriger un extrait de brouillon]
\textbf{Énoncé.} Voici le brouillon d'un élève. Repérez les fautes de cohérence, de subjectivité et de structure, puis proposez une version corrigée.

« Notre visite à la laiterie du bon coin s'est très bien passée, le directeur est un homme vraiment gentil. Je pense que les ouvriers sont paresseux car beaucoup ne faisaient rien. Nous avons vu les machines. Je suggère de licencier les ouvriers paresseux. »

\tcblower
\textbf{Corrections repérées :}
\begin{itemize}
\item \textbf{Structure} : pas d'introduction (ni contexte, ni objectif, ni plan) ; pas de bilan ni de formule de courtoisie.
\item \textbf{Subjectivité} : « un homme vraiment gentil », « je pense que les ouvriers sont paresseux » sont des jugements personnels.
\item \textbf{Cohérence et imprécision} : « beaucoup ne faisaient rien » est vague ; « nous avons vu les machines » n'apporte aucune information ; la suggestion (licencier) est déplacée et irréaliste pour un élève.
\end{itemize}
\textbf{Version corrigée :} « Conformément à votre instruction du 5 mars, nous avons visité, le 8 mars, la laiterie de la ville. Cette visite avait pour but d'observer la chaîne de pasteurisation. D'abord, le directeur a présenté l'entreprise ; ensuite, nous avons observé les 4 machines de la chaîne, dont une était à l'arrêt pour maintenance. Enfin, nous avons relevé les consignes d'hygiène affichées. En guise de bilan, la visite a été instructive ; nous suggérons qu'une seconde visite soit organisée pour les classes de Terminale. Nous vous prions d'agréer, Monsieur le Proviseur, l'expression de notre respectueuse considération. »
\end{exoresolu}

\courssub{Séance de synthèse : produire un rapport complet}

À partir des notes triées lors des séances précédentes, rédigez intégralement un rapport de visite d'étude ou de stage, en respectant toutes les étapes de la frise : brouillon, rédaction, relecture, version finale. L'évaluation au Probatoire / Baccalauréat technique suit un barème précis qu'il faut avoir en tête pendant la rédaction.

\begin{exemplebox}[Barème type au Probatoire / Baccalauréat technique]
\begin{center}
\begin{tabularx}{\linewidth}{|l|X|c|}\hline
\rowcolor{popblueL} \textbf{Critère} & \textbf{Ce qui est évalué} & \textbf{Poids} \\ \hline
Respect de la structure & Introduction complète, corps articulé, conclusion avec courtoisie & 25 \% \\ \hline
Pertinence du contenu & Faits exacts, chiffres, témoignages, suggestions réalistes & 35 \% \\ \hline
Correction de la langue & Orthographe, grammaire, conjugaison, ponctuation & 25 \% \\ \hline
Présentation & Lisibilité, mise en page, en-tête, propreté & 15 \% \\ \hline
\end{tabularx}
\end{center}
Le contenu pèse le plus lourd (35 \%) : soigner d'abord l'exactitude des faits rapportés.
\end{exemplebox}

\begin{popfigure}
\begin{center}
\begin{tikzpicture}[font=\small]
  \draw[draw=popdark, thick, fill=popblueL] (0,0) rectangle (8,1);
  \node at (4,0.5) {\textbf{Grille d'évaluation d'un rapport} (note sur 16)};
  \draw[draw=popdark, fill=popgreenL] (0,-1) rectangle (2,-1.9);
  \draw[draw=popdark] (2,-1) rectangle (8,-1.9);
  \node at (1,-1.45) {\textbf{Structure} /4};
  \node[align=left, text width=5.4cm] at (5,-1.45) {intro 4 phrases, connecteurs, bilan, courtoisie};
  \draw[draw=popdark, fill=poporangeL] (0,-1.9) rectangle (2,-2.8);
  \draw[draw=popdark] (2,-1.9) rectangle (8,-2.8);
  \node at (1,-2.35) {\textbf{Exactitude} /4};
  \node[align=left, text width=5.4cm] at (5,-2.35) {faits vérifiables, chiffres, sources citées};
  \draw[draw=popdark, fill=poppurpleL] (0,-2.8) rectangle (2,-3.7);
  \draw[draw=popdark] (2,-2.8) rectangle (8,-3.7);
  \node at (1,-3.25) {\textbf{Langue} /4};
  \node[align=left, text width=5.4cm] at (5,-3.25) {orthographe, accords, temps, ponctuation};
  \draw[draw=popdark, fill=popgoldL] (0,-3.7) rectangle (2,-4.6);
  \draw[draw=popdark] (2,-3.7) rectangle (8,-4.6);
  \node at (1,-4.15) {\textbf{Présentation} /4};
  \node[align=left, text width=5.4cm] at (5,-4.15) {en-tête, lisibilité, absence de ratures};
\end{tikzpicture}
\end{center}
\legende{Grille-critères d'évaluation d'un rapport : quatre critères notés chacun sur 4 points.}
\end{popfigure}

\courssub{Compte rendu collectif et remédiation}

La dernière séance est consacrée à la \emph{mise en commun} : plusieurs rapports d'élèves sont lus à voix haute, puis commentés collectivement à l'aide de la grille ci-dessus. Chaque élève procède d'abord à une \textbf{auto-évaluation} honnête de sa copie (structure, exactitude, langue, présentation), avant la confrontation avec l'appréciation du groupe et du professeur. Les points faibles récurrents — introduction incomplète, connecteurs absents, traces de subjectivité, fautes d'accord — font l'objet d'exercices de remédiation ciblés.

\begin{exercice}[À vous de rédiger]
À partir des notes suivantes, rédigez un rapport complet (introduction, corps, conclusion) d'environ vingt lignes : visite de l'entreprise CAMBOIS le 21 mai ; instruction du directeur des études du 15 mai ; objectif : observer la fabrication des portes en bois ; 3 ateliers visités (débit, assemblage, vernissage) ; 25 ouvriers ; affiche de sécurité « port de gants obligatoire » respectée par 20 ouvriers sur 25 ; suggestion : renforcer le respect des consignes.
\end{exercice}

\begin{corrige}[Exercice : éléments attendus]
\textbf{Introduction} : formule « Conformément à votre instruction du 15 mai... » ; date et lieu (21 mai, CAMBOIS) ; objectif (fabrication des portes) ; annonce du plan. \textbf{Corps} : trois paragraphes articulés (\emph{d'abord, ensuite, enfin}) pour les trois ateliers ; chiffres exacts (25 ouvriers, 20 sur 25) ; référence explicite à l'affiche de sécurité, sans jugement sur les 5 ouvriers sans gants. \textbf{Conclusion} : bilan synthétique, suggestion formulée poliment (« Nous suggérons que le respect des consignes soit renforcé... »), formule de courtoisie adaptée au directeur des études. La présentation doit comporter un en-tête (lieu, date, objet, destinataire).
\end{corrige}

\begin{aretenir}[L'essentiel de la section]
Un rapport se rédige en trois parties codifiées : une \textbf{introduction en quatre phrases} (contexte, circonstances, objectif, plan), un \textbf{corps objectif} (faits, chiffres, témoignages, connecteurs), une \textbf{conclusion en trois temps} (bilan, suggestions polies, courtoisie). La relecture traque la subjectivité, les ruptures de cohérence et les fautes de langue. Au Probatoire / Baccalauréat technique, la pertinence du contenu pèse 35 \% de la note : rapportez des faits, jamais des jugements.
\end{aretenir}

\newpage

\coursec{Activité d'intégration}

\coursec{Produire des écrits utilitaires : le rapport}

\courssub{Objectifs de l'activité}

À l'issue de cette activité d'intégration, l'élève sera capable de :
\begin{itemize}
\item trier et organiser des notes de terrain désordonnées pour établir le \cle{plan} d'un rapport ;
\item interpréter un \cle{texte fonctionnel} (tableau de données chiffrées) et un \cle{texte iconique} (affiche de sécurité) pour en tirer des constats ;
\item rédiger un \cle{rapport} complet en respectant la \cle{structure externe} (en-tête, objet, annonces, signature) et la \cle{structure interne} (introduction, corps en trois parties, conclusion avec suggestions) ;
\item employer le \cle{texte explicatif} (définitions, exemples, comparaisons, connecteurs logiques) et distinguer le \cle{lexique commun} du \cle{lexique spécialisé} de l'industrie chocolatière ;
\item évaluer un rapport à l'aide d'une grille-critères et produire une version corrigée en remédiation.
\end{itemize}

\courssub{Situation}

Le lycée technique de Bonabéri organise chaque année des visites d'étude pour rapprocher les élèves du monde de l'entreprise. Le vendredi 14 mars 2026, la classe de Première technique, accompagnée de son professeur de Français, M. Etoundi, et de son professeur de technologie, Mme Ngo Bell, a visité l'usine \cle{CHOCOCAM} (Chocolaterie Confiserie Camerounaise), située dans la zone industrielle de Douala-Bassa. De retour en classe, le proviseur demande à la classe de lui adresser un \textbf{rapport de mission} complet, qui servira à décider du renouvellement de ce partenariat pédagogique l'année prochaine.

Le professeur de Français remet à chaque groupe un \textbf{dossier documentaire} contenant :

\textbf{Document 1 : notes de terrain désordonnées} (prises par quatre élèves pendant la visite)
\begin{itemize}
\item « arrivée à l'usine à 9 h 10, accueil par M. Kamga, responsable de production » ;
\item « l'affiche près de la fèveuse : port de la charlotte, des gants et des chaussures de sécurité obligatoires » ;
\item « 3 ateliers visités : torréfaction, broyage-conchage, moulage-emballage » ;
\item « un ouvrier témoigne : "Depuis la formation d'octobre, nous n'avons plus eu d'accident de machine" » ;
\item « visite terminée à 12 h 30 ; dégustation de chocolat au lait » ;
\item « les fèves de cacao viennent surtout du Centre et du Sud-Ouest » ;
\item « la température de torréfaction est contrôlée automatiquement » ;
\item « 2 élèves ont oublié leur autorisation parentale et ont failli rester au portail » ;
\item « M. Kamga : "Nous recrutons des techniciens en maintenance chaque année" » ;
\item « bruit élevé dans l'atelier de broyage ; bouchons d'oreilles distribués ».
\end{itemize}

\textbf{Document 2 : tableau de données chiffrées} (fourni par le service production)

\begin{center}
\begin{tabularx}{\linewidth}{|l|X|X|}\hline
\rowcolor{popblueL} \textbf{Trimestre} & \textbf{Cacao transformé (tonnes)} & \textbf{Chocolat produit (tonnes)} \\ \hline
T1 (janvier--mars) & 1\,200 & 780 \\ \hline
T2 (avril--juin) & 1\,050 & 690 \\ \hline
T3 (juillet--septembre) & 900 & 600 \\ \hline
T4 (octobre--décembre) & 1\,350 & 880 \\ \hline
\end{tabularx}
\end{center}

\textbf{Document 3 : deux photographies} : la chaîne de torréfaction (fèves dans le tambour rotatif) et la chaîne de moulage (tablettes sortant des moules).

\textbf{Document 4 : affiche de sécurité de l'usine} : pictogrammes barrés de rouge (pas de téléphone portable, pas de bijoux), pictogrammes bleus d'obligation (charlotte, gants, chaussures de sécurité, bouchons d'oreilles), avec la consigne écrite : « La sécurité d'abord : zéro accident, c'est possible ».

\courssub{Consignes}

Par groupes de quatre élèves, et à l'aide du dossier documentaire :

\begin{enumerate}
\item \textbf{Trier} les dix notes de terrain en les classant dans trois rubriques : déroulement de la visite, observations sur la production, observations sur la sécurité. Barrer les notes hors sujet pour un rapport adressé au proviseur et justifier ce choix.
\item \textbf{Établir le plan} du rapport : une introduction (contexte, objectif, déroulement de la mission), un corps en trois parties, une conclusion avec suggestions.
\item \textbf{Interpréter le tableau chiffré} : calculer le total annuel de cacao transformé, identifier le trimestre le plus et le moins productif, calculer le taux de transformation moyen (chocolat produit / cacao transformé) et formuler deux constats rédigés.
\item \textbf{Lire l'affiche de sécurité} : identifier son émetteur, son destinataire, sa fonction, et distinguer les pictogrammes d'interdiction et d'obligation. En déduire un constat sur la politique de sécurité de l'usine.
\item \textbf{Relever dans le dossier} cinq termes de lexique spécialisé de l'industrie chocolatière et les définir avec des mots du lexique commun.
\item \textbf{Rédiger le rapport complet} (une page à une page et demie) : structure externe soignée (en-tête, lieu et date, objet, annonces « Monsieur le Proviseur », formule de politesse, signatures), introduction, corps en trois parties rédigé sous forme de texte explicatif (définitions, exemples, connecteurs logiques), conclusion avec deux suggestions.
\item \textbf{Présenter} le rapport devant la classe (5 minutes par groupe), puis évaluer les autres groupes à l'aide de la grille-critères fournie.
\item \textbf{Produire collectivement}, en remédiation, une version corrigée du rapport à partir des meilleurs passages de chaque groupe.
\end{enumerate}

\courssub{Ressources à mobiliser}

\begin{aretenir}[Ce qu'il faut se rappeler]
\begin{itemize}
\item \textbf{Structure externe du rapport} : en-tête (identité du rédacteur, fonction), lieu et date, objet, formule d'appel (« Monsieur le Proviseur »), formule de politesse, signature.
\item \textbf{Structure interne} : introduction (qui ? quoi ? où ? quand ? pourquoi ?), corps organisé en parties numérotées, conclusion (bilan + suggestions).
\item \textbf{Texte explicatif} : il définit, illustre par des exemples, compare, et utilise des connecteurs (\emph{d'abord, ensuite, en effet, c'est-à-dire, par conséquent, enfin}).
\item \textbf{Lexique} : le lexique commun est compris de tous ; le lexique spécialisé appartient à un domaine technique et doit être défini.
\item \textbf{Lecture des textes fonctionnels et iconiques} : identifier émetteur, destinataire, fonction ; dégager l'information chiffrée ou visuelle et la transformer en constat rédigé.
\item \textbf{Calculs utiles} : total = somme des valeurs ; taux de transformation $= \dfrac{\text{chocolat produit}}{\text{cacao transformé}} \times 100$.
\end{itemize}
\end{aretenir}

\courssub{Démarche et corrigé commenté}

\begin{exoresolu}[Rapport de mission : visite d'étude à l'usine CHOCOCAM]
Énoncé : à partir du dossier documentaire, produire le rapport de mission demandé par le proviseur, en suivant les huit consignes.

\tcblower

\textbf{Étape 1 : tri des notes.} On classe les dix notes en trois rubriques :
\begin{itemize}
\item \emph{Déroulement de la visite} : arrivée à 9 h 10 et accueil par M. Kamga ; visite des 3 ateliers ; fin à 12 h 30 avec dégustation.
\item \emph{Production} : origine des fèves (Centre et Sud-Ouest) ; température de torréfaction contrôlée automatiquement ; recrutement de techniciens en maintenance.
\item \emph{Sécurité} : affiche près de la fèveuse ; témoignage de l'ouvrier (zéro accident depuis octobre) ; bruit élevé et bouchons d'oreilles.
\end{itemize}
La note sur les deux élèves sans autorisation parentale est \textbf{barrée} : elle relève de l'organisation interne de la classe, pas du compte rendu de la mission adressé au proviseur. Justifier ce choix montre qu'un rapport sélectionne l'information pertinente.

\textbf{Étape 2 : plan du rapport.}
\begin{enumerate}
\item Introduction : contexte, objectif, date et encadrement de la visite.
\item Partie 1 : déroulement de la visite (accueil, trois ateliers, dégustation).
\item Partie 2 : observations sur la production (filière cacao, données chiffrées, emplois).
\item Partie 3 : observations sur la sécurité (affiche, équipements, témoignage).
\item Conclusion : bilan positif + deux suggestions.
\end{enumerate}

\textbf{Étape 3 : interprétation du tableau.}
\begin{align*}
\text{Total annuel} &= 1200 + 1050 + 900 + 1350 = 4500 \text{ tonnes de cacao}\\
\text{Taux moyen} &= \frac{780+690+600+880}{4500} \times 100 = \frac{2950}{4500} \times 100 \approx 65{,}6\,\%
\end{align*}
Constats rédigés : « Le quatrième trimestre est le plus productif (1\,350 tonnes), ce qui correspond à la grande campagne cacaoyère ; le troisième trimestre est le plus faible (900 tonnes). En moyenne, 100 kg de fèves donnent environ 65,6 kg de chocolat. »

\textbf{Étape 4 : lecture de l'affiche.} Émetteur : la direction de CHOCOCAM ; destinataires : ouvriers et visiteurs ; fonction : prévenir les accidents (texte prescriptif). Les pictogrammes rouges barrés signalent des \emph{interdictions} (téléphone, bijoux) ; les pictogrammes bleus signalent des \emph{obligations} (charlotte, gants, chaussures, bouchons d'oreilles). Constat : l'usine applique une politique de sécurité stricte, confirmée par le témoignage de l'ouvrier.

\textbf{Étape 5 : lexique spécialisé.} \emph{Torréfaction} : grillage des fèves pour développer l'arôme. \emph{Broyage} : transformation des fèves en pâte de cacao. \emph{Conchage} : brassage prolongé qui rend le chocolat onctueux. \emph{Moulage} : coulage du chocolat dans des moules. \emph{Charlotte} : coiffe jetable qui couvre les cheveux.

\textbf{Étape 6 : rédaction du rapport (modèle).}

\begin{quote}
\textbf{Classe de Première technique — Lycée technique de Bonabéri}\\
Douala, le 16 mars 2026\\
\textbf{Objet : rapport de mission sur la visite d'étude à l'usine CHOCOCAM}\\
Monsieur le Proviseur,

\textbf{Introduction.} Dans le cadre du programme de rapprochement école-entreprise, la classe de Première technique, encadrée par M. Etoundi et Mme Ngo Bell, a effectué le vendredi 14 mars 2026 une visite d'étude à l'usine CHOCOCAM de Douala-Bassa. La mission visait à découvrir le fonctionnement d'une industrie agroalimentaire camerounaise. Le présent rapport rend compte du déroulement de la visite, de nos observations sur la production et sur la sécurité.

\textbf{1. Déroulement de la visite.} D'abord, le groupe est arrivé à l'usine à 9 h 10 et a été accueilli par M. Kamga, responsable de production. Ensuite, nous avons visité trois ateliers : la torréfaction, c'est-à-dire le grillage des fèves, le broyage-conchage et le moulage-emballage. Enfin, la visite s'est achevée à 12 h 30 par une dégustation de chocolat au lait.

\textbf{2. Observations sur la production.} Les fèves proviennent principalement du Centre et du Sud-Ouest. La température de torréfaction est contrôlée automatiquement, ce qui garantit une qualité régulière. Selon les données du service production, l'usine a transformé 4\,500 tonnes de cacao en 2025, avec un pic au quatrième trimestre (1\,350 tonnes) lié à la grande campagne cacaoyère ; le taux de transformation moyen est d'environ 65,6\,\%. Par ailleurs, M. Kamga a précisé que l'usine recrute chaque année des techniciens en maintenance : c'est un débouché concret pour notre filière.

\textbf{3. Observations sur la sécurité.} Une affiche placée près de la fèveuse impose le port de la charlotte, des gants et des chaussures de sécurité, et interdit téléphones et bijoux. Dans l'atelier de broyage, où le bruit est élevé, des bouchons d'oreilles sont distribués. Un ouvrier a témoigné : « Depuis la formation d'octobre, nous n'avons plus eu d'accident de machine. » Ces éléments montrent que l'usine applique une politique de prévention efficace.

\textbf{Conclusion.} Cette mission a été très instructive : elle nous a permis de relier nos cours aux réalités industrielles. Nous suggérons, d'une part, de renouveler ce partenariat l'année prochaine et, d'autre part, d'inviter un technicien de CHOCOCAM à animer une conférence au lycée.

Nous vous prions d'agréer, Monsieur le Proviseur, l'expression de notre respectueuse considération.\\
\emph{Les délégués de la classe de Première technique} (signatures)
\end{quote}

\textbf{Étapes 7 et 8 : évaluation et remédiation.} La grille-critères note sur 20 : structure externe (4 pts), introduction et conclusion (4 pts), corps en trois parties avec texte explicatif (6 pts), exploitation des documents (3 pts), correction de la langue (3 pts). La version collective reprend les meilleures formulations de chaque groupe et corrige les erreurs relevées (objet oublié, parties non numérotées, chiffres non interprétés).
\end{exoresolu}

\courssub{Bilan de l'activité}

Cette activité a permis de mettre en évidence que le rapport est un \cle{écrit utilitaire} qui ne se contente pas de raconter : il \textbf{sélectionne}, \textbf{organise} et \textbf{interprète} l'information pour éclairer une décision (ici, le renouvellement du partenariat). Les élèves ont vérifié concrètement que :
\begin{itemize}
\item la \cle{préparation} (tri des notes, plan) conditionne la clarté du texte final ;
\item la \cle{structure externe} donne au document sa valeur officielle, tandis que la \cle{structure interne} guide le lecteur ;
\item un tableau chiffré ou une affiche ne parle pas tout seul : le rapporteur doit transformer les données en \textbf{constats rédigés} ;
\item le \cle{texte explicatif} rend le lexique spécialisé accessible grâce aux définitions et aux connecteurs logiques ;
\item l'évaluation par grille-critères et la remédiation collective développent l'esprit critique et la capacité à \textbf{réviser un texte}, compétence essentielle au Probatoire technique comme dans la vie professionnelle.
\end{itemize}

\newpage

\coursec{Méthodes et exercices résolus}

\coursec{Produire des écrits utilitaires : le rapport}

\courssub{Méthode 1 : Appliquer les notions de l'unité à une situation concrète — préparer un rapport}

\begin{methode}[Préparer la rédaction d'un rapport]
Un rapport ne s'improvise pas : il se prépare. Voici la démarche en étapes.
\begin{enumerate}
\item \cle{Analyser la commande} : qui demande le rapport (le destinataire) ? Pourquoi (l'objectif : informer, rendre compte d'une visite, d'un stage, d'un accident, d'une enquête) ? Sur quoi (le sujet précis) ?
\item \cle{Identifier le type de rapport} : rapport de stage, rapport de visite, rapport d'activité, rapport d'incident ou d'accident, rapport d'enquête. Chaque type a ses rubriques attendues.
\item \cle{Collecter les informations} : notes prises sur le terrain, documents, chiffres, témoignages, dates, lieux. Classer ces informations par thème.
\item \cle{Trier et hiérarchiser} : éliminer ce qui est hors sujet ; ordonner les faits (ordre chronologique ou ordre logique : constat, causes, conséquences, propositions).
\item \cle{Établir un plan} : introduction (contexte, objet, annonce du plan), corps (faits exposés objectivement), conclusion (bilan, recommandations éventuelles).
\item \cle{Choisir le bon lexique} : lexique spécialisé du domaine concerné (atelier, chantier, entreprise) employé avec exactitude, dans un registre neutre et objectif.
\end{enumerate}
\textbf{Réflexe à retenir} : un rapport rapporte des \emph{faits}, pas des sentiments. On écrit à la troisième personne ou à la première personne du pluriel (\og nous avons constaté \fg{}), jamais de ton familier.
\end{methode}

\begin{exoresolu}[Préparer un rapport de visite d'entreprise]
\textbf{Énoncé (type Bac)} : Dans le cadre de ta formation, ta classe a visité l'entreprise SOCAMETAL de Douala le 12 mars. Ton professeur te demande de préparer le rapport de cette visite. Présente : le destinataire, l'objectif, le type de rapport, les informations à collecter et le plan que tu suivras.
\tcblower
\textbf{Solution détaillée}
\begin{enumerate}
\item \textbf{Analyse de la commande.} Le destinataire est le professeur (et, par ricochet, l'administration du lycée). L'objectif est de rendre compte fidèlement d'une visite pédagogique. Le sujet : la visite de l'entreprise SOCAMETAL à Douala.
\item \textbf{Type de rapport.} Il s'agit d'un \emph{rapport de visite} (ou rapport de sortie pédagogique). Il comporte : les références de la visite (date, lieu, participants), le déroulement, les observations, le bilan.
\item \textbf{Informations à collecter.} Date et horaires (12 mars, départ 8 h, retour 14 h) ; nom et fonction du guide ; ateliers visités (soudure, tôlerie, contrôle qualité) ; effectif de l'entreprise ; matériels observés ; consignes de sécurité affichées ; réponses aux questions posées.
\item \textbf{Tri et hiérarchisation.} On écarte les détails sans rapport avec l'objectif (le menu de la collation, une panne du car). On retient les faits professionnels, classés dans l'ordre chronologique de la visite.
\item \textbf{Plan retenu.} Introduction : cadre de la visite (qui, quand, où, pourquoi). Corps : accueil et présentation de l'entreprise ; visite des ateliers ; échanges avec le personnel. Conclusion : bilan pédagogique et remerciements.
\end{enumerate}
\textbf{Conclusion} : la préparation aboutit à un plan clair et à un dossier de faits vérifiés, conditions indispensables d'un rapport objectif et complet.
\end{exoresolu}

\courssub{Méthode 2 : Construire la structure externe et interne du rapport}

\begin{methode}[Respecter la structure d'un rapport]
\begin{enumerate}
\item \cle{Structure externe} (la présentation matérielle) : la page de garde (titre, nom de l'auteur, destinataire, date, lieu), éventuellement un sommaire, les pages numérotées, les annexes (photos, tableaux, documents).
\item \cle{Structure interne} (l'organisation du contenu) : \textbf{introduction} (contexte et circonstances, objet du rapport, annonce du plan) ; \textbf{corps} divisé en parties numérotées et titrées (I, II, III ou 1., 2., 3.), chaque partie présentant des faits précis ; \textbf{conclusion} (bilan, suggestions ou recommandations).
\item \cle{Numéroter} les parties et sous-parties pour faciliter la lecture (1.1, 1.2...).
\item \cle{Signer et dater} le rapport à la fin, avec la mention de la fonction de l'auteur.
\end{enumerate}
\textbf{Formule à retenir} : \og Un rapport bien construit se lit comme une enquête : contexte, faits, bilan. \fg{}
\end{methode}

\begin{exoresolu}[Reconstituer la structure d'un rapport]
\textbf{Énoncé (type Bac)} : Voici, dans le désordre, les éléments d'un rapport d'accident survenu dans un atelier : (a) \og Nous suggérons de renforcer le port des équipements de protection. \fg{} ; (b) \og Rapport d'accident survenu le 5 mai à l'atelier de soudure — rédigé par M. Ngo Bell, chef d'atelier, à l'attention du directeur \fg{} ; (c) \og Le 5 mai à 10 h 15, l'ouvrier Etoundi s'est brûlé la main gauche en manipulant une pièce sans gants. \fg{} ; (d) \og Les témoins confirment que la consigne de sécurité était affichée mais non respectée. \fg{} Replace ces éléments dans l'ordre et justifie en nommant les parties de la structure interne et externe.
\tcblower
\textbf{Solution détaillée}
\begin{enumerate}
\item L'élément (b) vient \textbf{en premier} : c'est la \emph{page de garde / l'en-tête}, élément de la \textbf{structure externe} (titre, auteur, destinataire, date).
\item L'élément (c) constitue le début du \textbf{corps} du rapport : l'exposé des faits, dans l'ordre chronologique (date, heure, lieu, victime, nature de l'accident).
\item L'élément (d) prolonge le \textbf{corps} : c'est l'enquête, c'est-à-dire les causes et circonstances établies grâce aux témoignages. Il se place donc après (c).
\item L'élément (a) forme la \textbf{conclusion} : le bilan débouche sur des recommandations, ce qui est la fonction propre de la conclusion d'un rapport d'incident.
\end{enumerate}
\textbf{Ordre correct} : (b) — (c) — (d) — (a).
\textbf{Conclusion} : tout rapport suit la progression contexte, faits, causes, recommandations ; respecter cette structure interne garantit la clarté et l'utilité du document.
\end{exoresolu}

\courssub{Méthode 3 : Rédiger un texte explicatif (stylistique et rhétorique)}

\begin{methode}[Écrire un passage explicatif dans un rapport]
Le rapport repose sur le \cle{texte explicatif}, qui informe et fait comprendre sans chercher à convaincre.
\begin{enumerate}
\item \cle{Annoncer ce qu'on explique} : \og Nous allons expliquer le fonctionnement de la machine \fg{}, \og Voici comment s'est déroulée l'opération \fg{}.
\item \cle{Employer les marques de l'explication} : connecteurs logiques (\emph{d'abord, ensuite, enfin} ; \emph{parce que, car, en effet, c'est-à-dire, ainsi, par conséquent}), reformulations, définitions des termes techniques.
\item \cle{Utiliser les temps de l'objectivité} : présent de vérité générale pour décrire un fonctionnement (\og la presse découpe la tôle \fg{}), passé composé ou imparfait pour rapporter des faits passés.
\item \cle{Préférer les tournures neutres} : voix passive ou \og on \fg{} (\og la pièce a été contrôlée \fg{}), vocabulaire précis, phrases courtes.
\item \cle{Illustrer si nécessaire} par un exemple, un chiffre, un schéma légendé.
\end{enumerate}
\textbf{Piège} : ne pas glisser vers l'argumentation (jugements, opinions) : l'explication expose, elle ne débat pas.
\end{methode}

\begin{exoresolu}[Transformer un récit subjectif en texte explicatif]
\textbf{Énoncé (type Bac)} : Un élève écrit dans son rapport de stage : \og Franchement, la machine à découper, c'était génial, j'ai adoré la voir marcher, elle coupait la tôle super vite ! \fg{} Réécris ce passage sous la forme d'un texte explicatif adapté à un rapport, en trois ou quatre phrases.
\tcblower
\textbf{Solution détaillée}

On identifie d'abord les défauts du passage : jugements subjectifs (\og génial \fg{}, \og j'ai adoré \fg{}), registre familier (\og super vite \fg{}), absence de précision technique.

Correction étape par étape :
\begin{enumerate}
\item Supprimer les marques de subjectivité et le registre familier.
\item Remplacer par des faits observables et un lexique spécialisé : nom exact de la machine, fonction, résultat mesurable.
\item Employer le présent de vérité générale (fonctionnement) et le passé composé (observation ponctuelle).
\end{enumerate}
\textbf{Texte explicatif proposé} : \og L'atelier est équipé d'une machine de découpe laser. Cette machine découpe les tôles d'acier selon un tracé programmé sur ordinateur : d'abord, l'opérateur fixe la tôle sur la table ; ensuite, le faisceau laser réalise la découpe avec une précision de l'ordre du dixième de millimètre. Lors de notre visite, nous avons observé la découpe d'une tôle de 2 mm en moins d'une minute. \fg{}
\textbf{Conclusion} : le passage réécrit informe objectivement (nature, fonctionnement, performance chiffrée), ce qui correspond exactement à la fonction du texte explicatif dans un rapport.
\end{exoresolu}

\courssub{Méthode 4 : Distinguer lexique commun et lexique spécialisé}

\begin{methode}[Choisir le bon lexique]
\begin{enumerate}
\item \cle{Lexique commun} : mots de la langue courante, compris de tous (\emph{machine, outil, travail, danger}).
\item \cle{Lexique spécialisé} : termes propres à un domaine technique, précis et sans ambiguïté (\emph{tour à métaux, établi, disjoncteur, EPI} pour \og équipement de protection individuelle \fg{}).
\item Dans un \cle{rapport technique}, employer le lexique spécialisé dès qu'il existe un terme exact, mais \cle{définir} à la première occurrence les sigles et les termes que le destinataire pourrait ignorer : \og les EPI (équipements de protection individuelle) \fg{}.
\item Vérifier l'\cle{exactitude} du terme : un mot spécialisé mal employé discrédite tout le rapport.
\end{enumerate}
\textbf{Réflexe} : à la relecture, souligner chaque terme technique et se demander : est-il exact ? est-il défini si nécessaire ?
\end{methode}

\begin{exoresolu}[Lexique commun ou spécialisé ?]
\textbf{Énoncé (type Bac)} : Dans le rapport suivant, relève quatre termes de lexique spécialisé et deux de lexique commun, puis remplace les expressions vagues soulignées par le terme technique exact : \og L'ouvrier a réparé \emph{la chose qui coupe le courant} avec \emph{un outil pour visser}. Il a ensuite réglé le tour et nettoyé l'établi. \fg{}
\tcblower
\textbf{Solution détaillée}
\begin{enumerate}
\item \textbf{Termes de lexique spécialisé} présents : \emph{tour} (machine-outil pour l'usinage des pièces), \emph{établi} (table de travail de l'ouvrier), \emph{disjoncteur} (une fois corrigé), \emph{tournevis} (une fois corrigé).
\item \textbf{Termes de lexique commun} : \emph{ouvrier}, \emph{réparer}, \emph{nettoyer} (deux suffisent).
\item \textbf{Corrections} : \og la chose qui coupe le courant \fg{} devient \emph{le disjoncteur} (appareil qui coupe automatiquement le circuit électrique en cas de défaut) ; \og un outil pour visser \fg{} devient \emph{un tournevis}.
\end{enumerate}
\textbf{Texte corrigé} : \og L'ouvrier a réparé le disjoncteur avec un tournevis. Il a ensuite réglé le tour et nettoyé l'établi. \fg{}
\textbf{Conclusion} : le lexique spécialisé rend le rapport précis, court et professionnel ; les expressions vagues du lexique commun doivent être remplacées par le terme technique exact.
\end{exoresolu}

\courssub{Méthode 5 : Lire les textes fonctionnels et les affiches de sécurité}

\begin{methode}[Lire un texte fonctionnel ou une affiche de sécurité]
\begin{enumerate}
\item \cle{Identifier la nature du document} : notice, consigne, affiche, panneau, formulaire. Qui l'émet ? À qui s'adresse-t-il ?
\item \cle{Dégager la fonction} : informer, interdire, obliger, avertir d'un danger, guider.
\item \cle{Lire le texte iconique} : pour une affiche de sécurité, analyser le \emph{pictogramme} (forme, couleur : rouge = interdiction ou incendie ; jaune = avertissement ; bleu = obligation ; vert = sauvetage et secours), le symbole (flamme, tête de mort, flèche) et le court message écrit.
\item \cle{Repérer les marques linguistiques} : infinitifs de consigne (\og Porter un casque \fg{}), impératifs, phrases nominales (\og Danger de mort \fg{}), vocabulaire impératif ou interdit (\og Il est interdit de... \fg{}, \og Défense de... \fg{}).
\item \cle{Déduire le comportement attendu} : que doit faire concrètement le lecteur ?
\end{enumerate}
\end{methode}

\begin{exoresolu}[Analyser une affiche de sécurité]
\textbf{Énoncé (type Bac)} : À l'entrée d'un atelier, une affiche carrée à fond bleu représente, en blanc, un personnage portant un casque, avec l'inscription \og PORT DU CASQUE OBLIGATOIRE \fg{}. Analyse cette affiche : nature du document, signification de la couleur et du pictogramme, fonction, comportement attendu.
\tcblower
\textbf{Solution détaillée}
\begin{enumerate}
\item \textbf{Nature du document} : c'est un texte fonctionnel iconique, une \emph{affiche d'indication sécuritaire} placée à l'entrée d'un atelier, émise par l'employeur et adressée à toute personne entrant dans l'atelier.
\item \textbf{La couleur bleue} : dans le code des panneaux de sécurité, le fond bleu (forme ronde ou carrée selon les normes) signale une \emph{obligation}. Ce n'est ni une interdiction (rouge) ni un simple avertissement (jaune).
\item \textbf{Le pictogramme} : le personnage coiffé d'un casque désigne la partie du corps à protéger (la tête) et l'équipement exigé (le casque de chantier, élément des EPI).
\item \textbf{Le message écrit} : la phrase nominale \og Port du casque obligatoire \fg{} et l'adjectif \emph{obligatoire} expriment une contrainte absolue, sans possibilité de négociation.
\item \textbf{Comportement attendu} : toute personne qui franchit l'entrée de l'atelier doit porter un casque ; l'affiche vise à prévenir les accidents (chute d'objets, chocs).
\end{enumerate}
\textbf{Conclusion} : l'efficacité de l'affiche repose sur la combinaison du code couleur, du pictogramme universel et d'un message bref et impératif, compris immédiatement même par un lecteur pressé.
\end{exoresolu}

\courssub{Méthode 6 : Rédiger le rapport complet (introduction, corps, conclusion)}

\begin{methode}[Rédiger et justifier une démarche : produire un rapport]
\begin{enumerate}
\item \cle{Rédiger l'introduction} : présenter le contexte (qui a ordonné le rapport, quand, pourquoi), l'objet précis du rapport, puis annoncer le plan (\og Nous présenterons d'abord..., ensuite..., enfin... \fg{}).
\item \cle{Rédiger le corps} : une partie par thème, numérotée et titrée ; faits précis (dates, chiffres, noms), texte explicatif objectif, lexique spécialisé exact, connecteurs logiques.
\item \cle{Rédiger la conclusion} : bilan synthétique (ce que le rapport a établi), puis recommandations ou suggestions concrètes ; formule de politesse administrative si le destinataire est un supérieur.
\item \cle{Justifier la démarche} : être capable d'expliquer ses choix (ordre des parties, temps employés, vocabulaire) : c'est ce qui est demandé lorsqu'on \og rédige et justifie une démarche, une réponse ou une conclusion \fg{}.
\item \cle{Relire} : orthographe, cohérence des temps, exactitude des faits et des chiffres.
\end{enumerate}
\end{methode}

\begin{exoresolu}[Rédiger un rapport de stage (introduction et conclusion)]
\textbf{Énoncé (type Bac)} : Tu as effectué un stage de deux semaines (du 2 au 13 juin) au garage moderne de Bafoussam. Rédige l'introduction et la conclusion du rapport de stage que tu adresses à ton professeur principal, puis justifie en deux phrases ton choix de plan.
\tcblower
\textbf{Solution détaillée}

\textbf{Introduction rédigée} : \og Dans le cadre de notre formation en mécanique automobile, le professeur principal de la classe de Première technique nous a demandé d'effectuer un stage d'observation et de pratique. C'est ainsi que nous avons été accueillis au Garage moderne de Bafoussam, du 2 au 13 juin. Le présent rapport rend compte de ce stage : nous présenterons d'abord l'entreprise d'accueil, ensuite les tâches effectuées, enfin les enseignements tirés de cette expérience. \fg{}

Justification de l'introduction : elle respecte les trois éléments obligatoires — le \emph{contexte} (cadre de la formation, commande du professeur), l'\emph{objet} (rendre compte du stage), l'\emph{annonce du plan} (d'abord, ensuite, enfin).

\textbf{Conclusion rédigée} : \og En bilan, ce stage nous a permis de découvrir le fonctionnement réel d'un garage et de mettre en pratique les notions apprises en classe, notamment la vidange, le réglage du freinage et l'accueil de la clientèle. Nous suggérons que de tels stages soient renouvelés et prolongés, car ils complètent utilement la formation théorique. Nous remercions le patron du garage et son équipe pour leur encadrement, et nous prions Monsieur le professeur principal d'agréer l'expression de notre respectueuse considération. \fg{}

Justification de la conclusion : elle présente le \emph{bilan} (acquis du stage), des \emph{recommandations} (renouvellement) et une \emph{formule de politesse} adaptée au destinataire.

\textbf{Justification du plan} : l'ordre retenu (entreprise, tâches, bilan) est l'ordre logique et chronologique, qui permet au lecteur de suivre le déroulement du stage comme s'il y assistait ; chaque partie répond à une question précise : où ? quoi ? qu'en retirer ?

\textbf{Conclusion de l'exercice} : un rapport complet et bien structuré, rédigé dans un français objectif et correct, remplit pleinement la fonction utilitaire attendue au Baccalauréat technique.
\end{exoresolu}

\courssub{Synthèse et points de vigilance}

\begin{aretenir}[Les 5 erreurs qui coûtent des points au Bac]
\begin{enumerate}
\item \textbf{Confondre rapport et récit subjectif} : écrire \og j'ai trouvé ça super \fg{} au lieu de rapporter des faits observés. Le rapport exige l'objectivité : faits, dates, chiffres, troisième personne ou \og nous \fg{}.
\item \textbf{Négliger la structure externe} : oublier la page de garde (titre, auteur, destinataire, date), la numérotation des parties ou la signature. Un rapport sans en-tête ni signature est un document incomplet, donc pénalisé.
\item \textbf{Déséquilibrer la structure interne} : rédiger une introduction sans annonce de plan, ou une conclusion qui introduit des faits nouveaux au lieu de faire le bilan et de formuler des recommandations.
\item \textbf{Employer un lexique vague ou familier} : \og la chose \fg{}, \og le truc \fg{}, \og super vite \fg{}. Il faut le terme technique exact (\emph{disjoncteur, tour, établi, EPI}) et définir les sigles à la première occurrence.
\item \textbf{Mal lire les consignes de sécurité} : confondre les codes couleurs (bleu = obligation, rouge = interdiction, jaune = avertissement, vert = secours) ou ignorer le pictogramme. À l'examen, analyser toujours ensemble image, couleur et texte de l'affiche.
\end{enumerate}
\end{aretenir}

\begin{savaistu}[Le rapport dans la vie professionnelle camerounaise]
Au Cameroun, le rapport est un document administratif courant : rapport de stage (insertion professionnelle), rapport d'activité (évaluation annuelle), rapport de mission (déplacement), rapport d'incident (sécurité au travail, code du travail). La maîtrise de sa structure et de son langage objectif est une compétence transversale, évaluée au Probatoire et indispensable pour l'insertion dans les entreprises (SONARA, CAMRAIL, ENEO, PME du secteur informel structuré). Les codes couleurs des panneaux de sécurité (NF EN ISO 7010) sont identiques à ceux utilisés dans l'industrie camerounaise.
\end{savaistu}

\newpage

\coursec{Exercices}

\courssub{Je vérifie mes connaissances}

\begin{exercice}[QCM : Structure et vocabulaire du rapport]
Parmi les propositions suivantes, une seule est correcte. Choisis-la et justifie ton choix en une phrase.
\begin{enumerate}
    \item La structure externe d'un rapport comprend obligatoirement : 
    \begin{itemize}
        \item[A)] L'introduction, le développement, la conclusion.
        \item[B)] L'en-tête, l'objet, la référence, la date, la signature.
        \item[C)] Le titre, les remerciements, la bibliographie.
        \item[D)] Le sommaire, les annexes, l'index.
    \end{itemize}
    \item Le lexique spécialisé dans un rapport technique se caractérise par :
    \begin{itemize}
        \item[A)] L'utilisation de mots du langage courant uniquement.
        \item[B)] L'emploi de termes précis, propres à un domaine professionnel, souvent monosémiques.
        \item[C)] La présence exclusive de néologismes.
        \item[D)] L'absence de connecteurs logiques.
    \end{itemize}
    \item Dans un texte explicatif, le connecteur « \emph{par conséquent} » introduit :
    \begin{itemize}
        \item[A)] Une cause.
        \item[B)] Une conséquence.
        \item[C)] Une opposition.
        \item[D)] Un but.
    \end{itemize}
    \item L'introduction d'un rapport ne doit pas contenir :
    \begin{itemize}
        \item[A)] Le contexte et le déclencheur de la mission.
        \item[B)] L'objet précis de la mission.
        \item[C)] Les résultats détaillés et l'analyse des causes.
        \item[D)] Le plan annoncé du rapport.
    \end{itemize}
\end{enumerate}
\end{exercice}

\begin{exercice}[Vrai / Faux justifié : Notions de stylistique et réception]
Indique si chaque affirmation est Vraie ou Fausse, puis justifie ta réponse en te basant sur le cours.
\begin{enumerate}
    \item Le texte explicatif a pour unique fonction de divertir le lecteur.
    \item Une affiche de sécurité utilise un code couleur normé : le rouge indique l'interdiction ou le danger, le bleu l'obligation, le vert le secours ou la sortie.
    \item La conclusion d'un rapport se limite à remercier le destinataire.
    \item La préparation d'un rapport (collecte d'informations, tri, planification) est une étape facultative pour un rédacteur expérimenté.
\end{enumerate}
\end{exercice}

\begin{exercice}[Définitions et distinctions]
Donne une définition précise et concise pour chacun des termes suivants :
\begin{enumerate}
    \item Rapport (définition fonctionnelle).
    \item Structure interne vs structure externe.
    \item Lexique commun vs lexique spécialisé (donne un exemple pour chacun dans le contexte d'un atelier de menuiserie).
    \item Texte explicatif (finalité et procédés principaux).
\end{enumerate}
\end{exercice}

\begin{exercice}[Calculs directs et manipulations langagières]
\begin{enumerate}
    \item Classe les connecteurs suivants dans le tableau ci-dessous selon la relation logique qu'ils expriment : \emph{car, donc, afin que, puisque, par suite, dans le but de, c'est pourquoi, comme, pour que, ainsi}.
    \begin{center}
    \begin{tabularx}{\linewidth}{|X|X|X|X|}\hline
    \textbf{Cause} & \textbf{Conséquence} & \textbf{But} & \textbf{Autre (préciser)} \\ \hline
     &  &  &  \\ \hline
    \end{tabularx}
    \end{center}
    \item Transforme les phrases suivantes en remplaçant le lexique commun par du lexique spécialisé approprié (domaine : BTP / Chantier) :
    \begin{itemize}
        \item[a)] « Les ouvriers ont mis du ciment dans les trous. »
        \item[b)] « Le chef a regardé si le mur était droit. »
        \item[c)] « On a mis de l'eau dans le mélange. »
    \end{itemize}
\end{enumerate}
\end{exercice}

\courssub{Je m'entraîne}

\begin{exercice}[Réorganisation et justification : Rapport d'accident d'atelier]
Voici les rubriques d'un rapport d'accident survenu dans un atelier de soudure, présentées dans le désordre. 
\begin{enumerate}
    \item Remets-les dans l'ordre logique de la structure externe puis interne standard.
    \item Justifie l'enchaînement de trois rubriques successives de ton choix.
\end{enumerate}
\textbf{Rubriques mélangées :}
\begin{itemize}
    \item Témoignages des témoins (M. Ngono, M. Essomba)
    \item Objet : Accident de travail – Brûlure au 2e degré (soudeur M. Abena)
    \item Conclusion et recommandations (port obligatoire EPI, révision extincteurs)
    \item Analyse des causes (défaut de mise à la terre, absence gants isolants)
    \item En-tête : Lycée Technique de New-Bell / Atelier de Soudure
    \item Description chronologique des faits (10h15 : arc électrique ; 10h16 : projection ; 10h20 : prise en charge infirmerie)
    \item Référence : N° 045/AT/ST/2025
    \item Introduction (contexte, mission du délégué sécurité)
    \item Pièces jointes : Certificat médical, Fiche de poste, Photos du poste de travail
    \item Date et signature du rédacteur (Délégué Sécurité)
\end{itemize}
\end{exercice}

\begin{exercice}[Analyse d'un texte explicatif : Procédure de sécurité]
Lis l'extrait suivant (procédure de verrouillage / \emph{Lockout-Tagout}) et réponds aux questions.
\textit{« Avant toute intervention sur une machine, il est impératif de la mettre à l'arrêt complet. \textbf{En effet}, l'énergie résiduelle peut provoquer un redémarrage intempestif. \textbf{Afin de} garantir la sécurité de l'opérateur, on procède au cadenassage des sources d'énergie. \textbf{Par conséquent}, nulle personne ne peut retirer le cadenas sauf l'auteur du verrouillage. \textbf{Ainsi}, le risque d'écrasement ou de sectionnement est éliminé. »}
\begin{enumerate}
    \item Releve trois procédés d'explication utilisés dans ce texte (définition, illustration, comparaison, reformulation, etc.) et cite le passage concerné.
    \item Identifie cinq connecteurs logiques, précise pour chacun la relation exprimée (cause, conséquence, but, conclusion).
    \item Réécris le dernier paragraphe en remplaçant les connecteurs par des synonymes sans changer le sens.
\end{enumerate}
\end{exercice}

\begin{exercice}[Classement lexical et définition : Rapport de chantier]
Voici une liste de 15 mots extraits d'un rapport de chantier de construction (BTP). 
\begin{enumerate}
    \item Classe-les en deux colonnes : \textbf{Lexique commun} / \textbf{Lexique spécialisé (BTP)}.
    \item Choisis cinq termes du lexique spécialisé et propose une définition accessible à un lecteur profane (non initié).
\end{enumerate}
\textbf{Liste :} Béton, couler, ferraillage, grue, vibrer, bancher, dalle, coffrage, armature, ciment, niveau, maçon, équerre, adjuvants, cure.
\end{exercice}

\begin{exercice}[Réception d'image et rédaction intégrée : Affiche de sécurité]
Observe la description de l'affiche ci-dessous (pas d'image fournie, utilise la description).
\textbf{Description :} Fond bleu circulaire. Pictogramme blanc : visage de profil portant un masque de protection respiratoire (type FFP2). Texte en bas : « Port du masque obligatoire – Zone poussiéreuse ». Panneau carré à fond vert à côté : pictogramme croix blanche + « Poste de premiers secours – 50 m ».
\begin{enumerate}
    \item Analyse cette signalisation : type d'affiche (obligation/secours), code couleur, pictogrammes, message textuel.
    \item Rédige le paragraphe « Constat / Preuves » du rapport d'incident (inhalation de poussières par un élève) qui cite cette affiche comme élément de preuve du non-respect de la consigne. Respecte l'objectivité et le style explicatif.
\end{enumerate}
\end{exercice}

\courssub{Je me prépare au Bac}

\begin{exercice}[Type Bac : Rapport de mission suite à une catastrophe naturelle (Inondations Ndogpassi)]
\textbf{Contexte :} Suite aux inondations exceptionnelles du quartier Ndogpassi à Douala (octobre 2025), le proviseur du Lycée Technique de Ndogpassi a mandaté le délégué de la classe de 1re TMA (Technique Mécanique Automobile) pour rendre visite aux familles des élèves sinistrés et évaluer les besoins urgents.
\textbf{Consignes :} Rédige le \textbf{rapport complet} que le délégué adresse au Proviseur.
\begin{enumerate}
    \item Respecte scrupuleusement la \textbf{structure externe} (en-tête, référence, objet, date, destinataire, expéditeur, signature, pièces jointes).
    \item Respecte la \textbf{structure interne} : Introduction (contexte, mission, objet, plan), Corps (au moins 3 parties logiques : 1. Étendue des dégâts / Situation des familles visitées ; 2. Besoins identifiés (logement, matériel scolaire, santé) ; 3. Actions menées / Propositions de solutions), Conclusion (synthèse, recommandations, ouverture).
    \item Adopte le style explicatif (objectivité, connecteurs logiques, lexique spécialisé adapté : \emph{sinistré, relogement, kit scolaire, vecteur paludéen, solidarité}).
    \item Longueur visée : environ 300-400 mots.
\end{enumerate}
\end{exercice}

\begin{exercice}[Type Bac : Rapport de stage – Introduction et Conclusion uniquement]
\textbf{Contexte :} Tu as effectué un stage de 4 semaines à la Cameroon Tea Estates (CTE) de Ndawara (Région de l'Adamaoua), service « Maintenance Industrielle ».
\textbf{Plan du corps du rapport (fourni) :}
\begin{enumerate}
    \item Présentation de l'entreprise (historique, capacité de production, organigramme).
    \item Description du service maintenance (moyens humains, matériel, planning préventif/curatif).
    \item Activités réalisées : Révision complète broyeurs à thé (semaine 2) ; Dépannage convoyeur à bande n°3 (semaine 3) ; Participation inventaire pièces de rechange (semaine 4).
    \item Analyse des compétences acquises (techniques : soudure TIG, lecture plans hydrauliques ; transverses : sécurité, travail en équipe).
    \item Difficultés rencontrées et solutions apportées.
\end{enumerate}
\textbf{Consignes :} Rédige \textbf{uniquement l'Introduction et la Conclusion} de ce rapport de stage.
\begin{enumerate}
    \item \textbf{Introduction :} Contexte (cadre scolaire, dates, lieu), Objet du stage, Problématique/Intérêt (lien formation/entreprise), Annonce du plan.
    \item \textbf{Conclusion :} Bilan global (réponse à la problématique), Compétences validées / à consolider, Apport pour l'insertion professionnelle, Remerciements formels, Perspective (projet professionnel).
\end{enumerate}
Soigne la cohérence avec le plan de corps fourni. Style impersonnel ou « je » modéré selon usage technique.
\end{exercice}

\begin{exercice}[Type Probatoire Technique : Correction de brouillon – Détection d'erreurs]
Voici un brouillon de rapport d'intervention technique (remplacement d'un variateur de fréquence sur une pompe de forage) qui contient \textbf{5 fautes de subjectivité}, \textbf{3 ruptures de cohérence} (logique ou chronologique) et \textbf{2 erreurs de structure externe}.
\textbf{Texte du brouillon :}
\textit{« Lycée Technique de Bafoussam. Moi, Kamga, je vais vous dire ce qui s'est passé. C'était nul au début.
Objet : Changement variateur.
Le 12/03/2025. Pour : Monsieur le Proviseur.
Introduction : La pompe marchait mal. C'est la faute à l'ancien variateur qui est pourri. J'ai dû tout faire tout seul parce que mes camarades sont paresseux. Mon but c'est de montrer que je suis le meilleur.
Corps : 1. D'abord j'ai acheté le nouveau variateur au marché. 2. Ensuite je l'ai monté sans couper le courant, c'était plus vite. 3. Heureusement il n'y a pas eu d'étincelles. 4. Après j'ai réglé les paramètres au feeling. 5. Finalement la pompe a démarré, c'est magnifique.
Conclusion : En conclusion, je suis content. Le proviseur doit m'acheter un cadeau. Merci.
Signé : Kamga. Pièces jointes : Facture du marché, photo floue. »}
\textbf{Consignes :}
\begin{enumerate}
    \item Recopie le tableau ci-dessous et complète-le en relevant \textbf{exactement} les 10 erreurs demandées (cite le passage fautif, nomme le type d'erreur : subjectivité / rupture cohérence / erreur structure externe, propose la correction formelle).
    \begin{center}
    \begin{tabularx}{\linewidth}{|X|X|X|X|}\hline
    \textbf{Passage fautif (citation)} & \textbf{Type d'erreur} & \textbf{Explication brève} & \textbf{Correction proposée (style rapport)} \\ \hline
     &  &  &  \\ \hline
     &  &  &  \\ \hline
     &  &  &  \\ \hline
     &  &  &  \\ \hline
     &  &  &  \\ \hline
     &  &  &  \\ \hline
     &  &  &  \\ \hline
     &  &  &  \\ \hline
     &  &  &  \\ \hline
     &  &  &  \\ \hline
    \end{tabularx}
    \end{center}
    \item Réécris l'\textbf{Introduction} et le \textbf{Corps} (partie 1 à 5) de ce rapport en version finale, corrigée, structurée et professionnelle.
\end{enumerate}
\end{exercice}

\newpage

\coursec{Corrigés des exercices}

\coursec{Leçon 16 : Produire des écrits utilitaires : le rapport}

\begin{corrige}[Exercice 1 — QCM : Structure et vocabulaire du rapport]
\begin{enumerate}
    \item \textbf{Réponse B.} La structure externe d'un rapport comprend les éléments qui encadrent le texte proprement dit : en-tête, objet, référence, date, signature (et pièces jointes éventuelles). L'introduction, le développement et la conclusion relèvent de la structure \emph{interne}.
    \item \textbf{Réponse B.} Le lexique spécialisé est constitué de termes précis, propres à un domaine professionnel, et souvent \cle{monosémiques} (un seul sens), ce qui garantit la clarté technique.
    \item \textbf{Réponse B.} « Par conséquent » est un connecteur de \cle{conséquence} : il introduit le résultat logique de ce qui précède.
    \item \textbf{Réponse C.} Les résultats détaillés et l'analyse des causes appartiennent au \emph{corps} du rapport ; l'introduction ne présente que le contexte, la mission, l'objet et le plan.
\end{enumerate}
\end{corrige}

\begin{corrige}[Exercice 2 — Vrai / Faux justifié]
\begin{enumerate}
    \item \textbf{Faux.} Le texte explicatif a pour fonction d'\emph{informer}, d'\emph{expliquer} et de faire comprendre un phénomène, une procédure ou un fonctionnement ; il relève du registre didactique, non du divertissement.
    \item \textbf{Vrai.} La signalétique de sécurité est normée (ISO 7010) : rouge = interdiction/danger/incendie ; bleu = obligation ; vert = secours, évacuation, issue de secours ; jaune = avertissement.
    \item \textbf{Faux.} La conclusion d'un rapport fait la \emph{synthèse} des constats, formule des \emph{recommandations} et peut ouvrir sur des perspectives ; les remerciements ne sont qu'un élément éventuel et secondaire.
    \item \textbf{Faux.} La préparation (collecte des informations, vérification des faits, tri, planification) est une étape \emph{indispensable} : sans données fiables et un plan clair, aucun rapport, même rédigé par un professionnel expérimenté, ne peut être objectif et cohérent.
\end{enumerate}
\end{corrige}

\begin{corrige}[Exercice 3 — Définitions et distinctions]
\begin{enumerate}
    \item \textbf{Rapport :} écrit fonctionnel et utilitaire, objectif et structuré, par lequel une personne \emph{mandatée} rend compte à un destinataire déterminé d'une mission, d'une visite, d'un incident ou d'une situation, en exposant des faits constatés, leur analyse et des propositions.
    \item \textbf{Structure externe :} ensemble des éléments d'encadrement du document (en-tête, référence, objet, date, destinataire, expéditeur, signature, pièces jointes). \textbf{Structure interne :} organisation du contenu en introduction, corps (parties logiques) et conclusion.
    \item \textbf{Lexique commun :} mots de la langue courante, compris de tous, souvent polysémiques ; ex. en menuiserie : \emph{bois, clou, couper, trou}. \textbf{Lexique spécialisé :} termes précis d'un domaine professionnel ; ex. en menuiserie : \emph{tenon, mortaise, dégauchir, raboter, assemblage à mi-bois}.
    \item \textbf{Texte explicatif :} texte dont la finalité est de faire comprendre un phénomène, un fonctionnement ou une procédure. Ses procédés principaux : définitions, exemples et illustrations, comparaisons, reformulations, énumérations, connecteurs logiques (cause, conséquence, but), lexique spécialisé défini pour le lecteur, style objectif et impersonnel.
\end{enumerate}
\end{corrige}

\begin{corrige}[Exercice 4 — Manipulations langagières : connecteurs et lexique spécialisé]
\begin{enumerate}
    \item Classement des connecteurs :
    \begin{center}
    \begin{tabularx}{\linewidth}{|X|X|X|X|}\hline
    \textbf{Cause} & \textbf{Conséquence} & \textbf{But} & \textbf{Autre} \\ \hline
    car ; puisque ; comme & donc ; par suite ; c'est pourquoi ; ainsi & afin que ; dans le but de ; pour que & --- \\ \hline
    \end{tabularx}
    \end{center}
    \emph{Remarque :} « ainsi » peut aussi exprimer la manière ; ici, dans un raisonnement, il marque la conséquence (« ainsi » = « de cette façon, il en résulte »).
    \item Reformulations en lexique spécialisé (BTP) :
    \begin{itemize}
        \item[a)] « Les maçons ont \textbf{coulé du béton} dans les \textbf{fouilles} (ou : ont procédé au \textbf{remblaiement en béton} des excavations). »
        \item[b)] « Le \textbf{chef de chantier} a \textbf{contrôlé l'aplomb} du mur (à l'aide du fil à plomb / du niveau). »
        \item[c)] « On a procédé au \textbf{gâchage} du mortier (ajout de l'eau de gâchage au mélange). »
    \end{itemize}
\end{enumerate}
\end{corrige}

\begin{corrige}[Exercice 5 — Réorganisation : Rapport d'accident d'atelier]
\begin{enumerate}
    \item \textbf{Ordre correct de la structure externe et interne :}
    \begin{enumerate}
        \item En-tête : Lycée Technique de New-Bell / Atelier de Soudure
        \item Référence : N° 045/AT/ST/2025
        \item Objet : Accident de travail – Brûlure au 2e degré (soudeur M. Abena)
        \item Date (placée en en-tête) ; Destinataire (À Monsieur le Proviseur) ; Expéditeur (De : Délégué Sécurité)
        \item Introduction (contexte, mission du délégué sécurité, objet, plan)
        \item Description chronologique des faits (10h15 ; 10h16 ; 10h20)
        \item Témoignages des témoins (M. Ngono, M. Essomba)
        \item Analyse des causes (défaut de mise à la terre, absence de gants isolants)
        \item Conclusion et recommandations (port obligatoire des EPI, révision des extincteurs)
        \item Signature du rédacteur (Nom + Qualité)
        \item Pièces jointes : certificat médical, fiche de poste, photos
    \end{enumerate}
    \item \textbf{Justification de trois enchaînements logiques :}
    \begin{itemize}
        \item La \emph{description chronologique des faits} précède les \emph{témoignages}, car il faut d'abord établir la version factuelle du rédacteur avant de la confronter aux déclarations des témoins.
        \item Les \emph{témoignages} précèdent l'\emph{analyse des causes}, car ils fournissent les éléments de preuve sur lesquels repose l'analyse.
        \item L'\emph{analyse des causes} précède la \emph{conclusion et les recommandations}, car on ne peut proposer des mesures préventives (port des EPI) qu'après avoir identifié les causes de l'accident.
    \end{itemize}
\end{enumerate}
\end{corrige}

\begin{corrige}[Exercice 6 — Analyse d'un texte explicatif : procédure de verrouillage (consignation)]
\begin{enumerate}
    \item \textbf{Procédés d'explication identifiés :}
    \begin{itemize}
        \item \emph{Explicitation causale / justification :} « \textbf{En effet}, l'énergie résiduelle peut provoquer un redémarrage intempestif » (on explique pourquoi l'arrêt complet est impératif).
        \item \emph{Énumération des étapes d'une procédure :} arrêt complet $\rightarrow$ cadenassage $\rightarrow$ interdiction de retrait (progression méthodique).
        \item \emph{Définition implicite par la finalité :} « on procède au cadenassage des sources d'énergie \textbf{afin de} garantir la sécurité de l'opérateur » (le but définit l'action).
        \item \emph{Illustration du risque :} « le risque d'écrasement ou de sectionnement » donne des exemples concrets de dangers.
    \end{itemize}
    \item \textbf{Connecteurs et relations logiques :}
    \begin{center}
    \begin{tabularx}{\linewidth}{|l|X|}\hline
    \textbf{Connecteur} & \textbf{Relation exprimée} \\ \hline
    En effet & Cause / justification \\ \hline
    Afin de & But \\ \hline
    Par conséquent & Conséquence \\ \hline
    Ainsi & Conséquence / conclusion \\ \hline
    Avant toute intervention & Antériorité (ordre chronologique de la procédure) \\ \hline
    \end{tabularx}
    \end{center}
    \item \textbf{Réécriture du passage final (remplacement de connecteurs) :}\\
    « \textbf{De ce fait}, nulle personne ne peut retirer le cadenas sauf l'auteur du verrouillage. \textbf{C'est pourquoi} le risque d'écrasement ou de sectionnement est éliminé. » (Autres possibles : \emph{donc, par suite, en conséquence, si bien que}.)
\end{enumerate}
\end{corrige}

\begin{corrige}[Exercice 7 — Classement lexical : rapport de chantier (BTP)]
\begin{enumerate}
    \item \textbf{Classement :}
    \begin{center}
    \begin{tabularx}{\linewidth}{|X|X|}\hline
    \textbf{Lexique commun} & \textbf{Lexique spécialisé (BTP)} \\ \hline
    béton, couler, grue, vibrer, ciment, niveau, maçon, équerre & ferraillage, bancher, dalle, coffrage, armature, adjuvants, cure \\ \hline
    \end{tabularx}
    \end{center}
    \emph{Remarque :} certains mots sont frontaliers : « couler » et « vibrer » appartiennent au vocabulaire courant mais prennent un sens technique précis au chantier (couler une dalle, vibrer le béton). On les accepte dans les deux colonnes si la justification est donnée. « Dalle » est un mot courant mais désigne ici un ouvrage technique précis (dalle en béton armé).
    \item \textbf{Définitions pour un profane (style explicatif) :}
    \begin{itemize}
        \item \emph{Ferraillage :} ensemble des barres d'acier placées dans le béton pour le rendre plus résistant à la traction.
        \item \emph{Coffrage :} moule provisoire (en bois ou en métal) dans lequel on verse le béton frais pour lui donner sa forme définitive.
        \item \emph{Bancher :} action de couler du béton dans un coffrage pour réaliser un mur ou une voile.
        \item \emph{Dalle :} plaque horizontale en béton armé formant un plancher ou un sol.
        \item \emph{Cure :} période pendant laquelle on maintient le béton humide après le coulage pour qu'il durcisse correctement et atteigne sa résistance.
        \item \emph{Adjuvants :} produits ajoutés en petite quantité au béton pour modifier ses propriétés (fluidité, prise plus rapide ou plus lente, résistance au gel).
    \end{itemize}
\end{enumerate}
\end{corrige}

\begin{corrige}[Exercice 8 — Réception d'image et rédaction intégrée (signalétique)]
\begin{enumerate}
    \item \textbf{Analyse de la signalisation :}
    \begin{itemize}
        \item \emph{Panneau 1 :} panneau d'\textbf{obligation} : forme ronde, fond bleu, pictogramme blanc (code couleur normé ISO 7010 : bleu = obligation). Le pictogramme représente un visage portant un masque respiratoire FFP2 ; le texte « Port du masque obligatoire – Zone poussiéreuse » précise la consigne et la zone concernée.
        \item \emph{Panneau 2 :} panneau de \textbf{secours / sauvetage} : forme carrée (ou rectangulaire), fond vert, pictogramme blanc (vert = secours, évacuation, issue de secours). Il indique la localisation du poste de premiers secours à 50 m.
        \item \emph{Complémentarité :} le premier impose une conduite préventive (protection respiratoire), le second indique les moyens de secours curatifs en cas d'accident (premiers soins).
    \end{itemize}
    \item \textbf{Paragraphe « Constat / Preuves » (modèle de rédaction objective) :}\\
    « Lors de notre visite de l'atelier, le 14 octobre 2025 à 9 h 30, nous avons constaté que l'élève N., affecté au poste de ponçage, travaillait sans masque de protection respiratoire. Or, un panneau d'obligation à fond bleu, apposé à l'entrée de la zone, prescrit explicitement le port du masque FFP2 avec la mention "Port du masque obligatoire – Zone poussiéreuse". Par conséquent, l'incident (inhalation de poussières) résulte du non-respect d'une consigne de sécurité clairement affichée. Le panneau de secours situé à proximité (poste de premiers secours à 50 m) a permis une prise en charge rapide de la victime. »
\end{enumerate}
\textbf{Points de vigilance :} style objectif et impersonnel (« nous avons constaté », pas de jugement sur l'élève) ; faits datés et localisés ; connecteurs logiques (« or », « par conséquent ») ; citation exacte de l'affiche comme preuve.
\end{corrige}

\begin{corrige}[Exercice 9 — Type Bac : Rapport de mission (inondations Ndogpassi)]
\textbf{Modèle de rapport complet :}

\begin{center}
\textbf{Lycée Technique de Ndogpassi — Douala}\\
Réf. : N° 012/DEL/1TMA/2025\\
Douala, le 20 octobre 2025\\
\textbf{À Monsieur le Proviseur du Lycée Technique de Ndogpassi}\\
De : Le délégué de la classe de Première TMA\\
\textbf{Objet :} Rapport de mission de visite aux familles des élèves sinistrés du quartier Ndogpassi.\\
Pièces jointes : liste des familles visitées, fiches d'évaluation des besoins, photographies.
\end{center}

\textbf{Introduction}\\
À la suite des inondations exceptionnelles qui ont frappé le quartier Ndogpassi au début du mois d'octobre 2025, Monsieur le Proviseur m'a mandaté, en ma qualité de délégué de la classe de Première TMA, pour rendre visite aux familles des élèves sinistrés et évaluer leurs besoins urgents. La présente mission a été menée du 15 au 18 octobre 2025. Ce rapport expose d'abord l'étendue des dégâts et la situation des familles visitées, puis les besoins identifiés, enfin les actions menées et les solutions proposées.

\textbf{1. Étendue des dégâts et situation des familles}\\
Huit familles d'élèves ont été visitées. Dans six cas, les habitations ont été envahies par les eaux sur une hauteur de 0,80 m à 1,50 m ; le mobilier, les denrées alimentaires et les documents administratifs ont été détruits. Deux familles ont dû être provisoirement relogées chez des proches. En outre, les eaux stagnantes favorisent la prolifération du vecteur paludéen, ce qui constitue un risque sanitaire réel.

\textbf{2. Besoins identifiés}\\
Les besoins relevés sont de trois ordres. D'abord, le logement : assèchement, désinfection et réparation des habitations. Ensuite, le matériel scolaire : douze élèves ont perdu leurs manuels, cahiers et tenues ; la fourniture de kits scolaires est urgente pour éviter le décrochage. Enfin, la santé : une sensibilisation antipaludique et la distribution de moustiquaires imprégnées s'imposent.

\textbf{3. Actions menées et propositions}\\
Une collecte de vivres et de vêtements a été organisée au sein de l'établissement ; les dons ont été remis aux familles le 18 octobre. Nous proposons, par conséquent, la mise en place d'un fonds de solidarité de l'établissement, l'octroi de kits scolaires aux élèves sinistrés et une campagne d'assainissement menée avec les services municipaux.

\textbf{Conclusion}\\
En synthèse, les inondations de Ndogpassi ont gravement affecté huit familles d'élèves, dont deux relogées d'urgence. Les besoins les plus urgents concernent le relogement, le matériel scolaire et la prévention sanitaire. Nous recommandons l'adoption rapide des mesures proposées afin de garantir la continuité de la scolarité des élèves sinistrés. La solidarité déjà manifestée par la communauté scolaire laisse espérer un retour rapide à la normale.

Signature du délégué.

\textbf{Barème indicatif (sur 20) :} structure externe complète (4 pts) ; introduction conforme : contexte, mission, plan (3 pts) ; corps en 3 parties logiques et factuelles (6 pts) ; conclusion : synthèse et recommandations (3 pts) ; style explicatif : objectivité, connecteurs, lexique spécialisé (3 pts) ; correction de la langue (1 pt).

\textbf{Points de vigilance :} ne pas confondre rapport et lettre (pas de formule de politesse oratoire) ; rester objectif (faits, chiffres, pas de pathos) ; employer correctement le lexique attendu (\emph{sinistré, relogement, kit scolaire, vecteur paludéen, solidarité}) ; annoncer un plan et le respecter.
\end{corrige}

\begin{corrige}[Exercice 10 — Type Bac : Rapport de stage — Introduction et Conclusion]
\textbf{Introduction (modèle) :}\\
« Dans le cadre de notre formation en classe de Première technique au Lycée Technique de Bafoussam, un stage d'immersion professionnelle de quatre semaines a été effectué du 2 juin au 27 juin 2025 à la Cameroon Tea Estates (CTE) de Ndawara, dans la Région de l'Adamaoua, au sein du service Maintenance Industrielle. Ce stage avait pour objet de confronter les enseignements théoriques reçus en atelier aux réalités de la maintenance d'une unité industrielle de transformation du thé. Il s'agissait de déterminer dans quelle mesure les compétences acquises en formation permettent de participer efficacement aux opérations de maintenance préventive et curative d'une entreprise agro-industrielle. Le présent rapport présente, en cinq parties : 1) l'entreprise d'accueil et son environnement, 2) le service maintenance et son organisation, 3) les activités réalisées et les techniques mises en œuvre, 4) les compétences acquises et validées, 5) les difficultés rencontrées et les solutions apportées. »

\textbf{Conclusion (modèle) :}\\
« Au terme de ce stage, il apparaît que la formation reçue au lycée constitue une base solide, qu'il convient toutefois de compléter par l'expérience des procédures industrielles réelles : la révision des broyeurs à thé et le dépannage du convoyeur à bande n°3 ont montré l'écart entre l'exercice scolaire et la contrainte de production. Sur le plan technique, les compétences en soudure TIG et en lecture de plans hydrauliques ont été validées ; la rigueur du planning de maintenance préventive reste à consolider. Sur le plan humain, le respect des consignes de sécurité et le travail en équipe ont été pleinement intégrés. Ce stage confirme notre projet professionnel : l'obtention du Baccalauréat technique, puis la spécialisation en maintenance industrielle. Nous remercions la Direction de la CTE de Ndawara, le chef du service Maintenance ainsi que l'ensemble de l'équipe pour leur accueil et leur encadrement, et nos enseignants pour le suivi assuré tout au long de la période. »

\textbf{Barème indicatif (sur 10) :} introduction : contexte complet (dates, lieu, cadre scolaire) (2 pts), objet et problématique (1 pt), annonce du plan en 5 parties cohérente avec le corps (1 pt) ; conclusion : bilan répondant à la problématique (2 pts), compétences et insertion professionnelle (1 pt), remerciements et perspective (1 pt) ; cohérence avec le plan de corps, style et langue (2 pts).

\textbf{Points de vigilance :} l'annonce du plan doit reprendre \emph{exactement} les cinq parties du corps fourni ; la conclusion doit répondre à la problématique posée en introduction ; le « je » modéré ou le « nous » impersonnel est admis, mais pas de familiarité ; citer les éléments du corps (broyeurs, convoyeur n°3, TIG) pour assurer la cohérence.
\end{corrige}

\begin{corrige}[Exercice 11 — Type Probatoire : Correction de brouillon (révision professionnelle)]
\begin{enumerate}
    \item \textbf{Tableau d'analyse des erreurs (brouillon de Kamga) :}
    \begin{center}
    \begin{tabularx}{\linewidth}{|X|X|X|X|}\hline
    \textbf{Passage fautif} & \textbf{Type} & \textbf{Explication} & \textbf{Correction} \\ \hline
    « Moi, Kamga, je vais vous dire ce qui s'est passé » & Subjectivité / Structure & Ton familier, « je » emphatique, apostrophe au lecteur, absence d'en-tête & « Le soussigné rend compte de l'intervention effectuée » (en-tête formel requis) \\ \hline
    « C'était nul au début » & Subjectivité & Jugement de valeur familier, sans précision technique & « Le fonctionnement de la pompe était défaillant » \\ \hline
    « l'ancien variateur qui est pourri » & Subjectivité & Registre familier, accusation non étayée techniquement & « le variateur défectueux (hors service) » \\ \hline
    « mes camarades sont paresseux » / « je suis le meilleur » & Subjectivité & Attaques personnelles et vantardise, hors de propos professionnel & Supprimer ; le rapport décrit des faits techniques, pas des personnes \\ \hline
    « réglé les paramètres au feeling » / « c'est magnifique » & Subjectivité & Approximation et enthousiasme non professionnels & « les paramètres ont été réglés selon la notice du constructeur ; les essais sont concluants » \\ \hline
    « monté sans couper le courant, c'était plus vite » & Rupture de cohérence / Sécurité & Violation de la procédure de consignation, inadmissible dans un rapport technique & « après consignation (coupure et verrouillage de l'alimentation), le variateur a été monté » \\ \hline
    « Heureusement il n'y a pas eu d'étincelles » & Rupture de cohérence & Contredit la logique de sécurité : le risque ne devait pas exister grâce à la consignation & Supprimer ; la consignation supprime le risque électrique à la source \\ \hline
    Chronologie : achat au marché avant tout diagnostic & Rupture de cohérence & On ne remplace pas une pièce sans diagnostic préalable (méthode scientifique) & « 1. Diagnostic : variateur hors service. 2. Approvisionnement. 3. Montage. 4. Paramétrage. 5. Essais » \\ \hline
    En-tête incomplet : pas de référence & Erreur de structure externe & La référence est obligatoire pour la traçabilité & Ajouter : « Réf. : N° 015/IT/2025 » \\ \hline
    « Pour : Monsieur le Proviseur » mal placé ; signature sans qualité & Erreur de structure externe & Destinataire et expéditeur doivent être identifiés en tête ; la signature précise nom et fonction & « À Monsieur le Proviseur / De : Kamga X., élève de 1re Électrotechnique » ; signature : nom + qualité \\ \hline
    \end{tabularx}
    \end{center}
    \item \textbf{Version corrigée (rapport professionnel) :}

    \textbf{Lycée Technique de Bafoussam — Atelier Électrotechnique}\\
    Réf. : N° 015/IT/2025\\
    Bafoussam, le 12 mars 2025\\
    \textbf{À Monsieur le Proviseur du Lycée Technique de Bafoussam}\\
    De : Kamga X., élève de 1re Électrotechnique\\
    \textbf{Objet :} Rapport d'intervention — Remplacement du variateur de fréquence de la pompe de forage.\\
    Pièces jointes : facture d'achat, notice constructeur, procès-verbal d'essais.

    \textbf{Introduction}\\
    À la suite de la défaillance de la pompe de forage de l'établissement, constatée le 10 mars 2025, une intervention de remplacement du variateur de fréquence a été effectuée le 12 mars 2025. Le présent rapport décrit le diagnostic établi, les opérations réalisées et les résultats des essais.

    \textbf{Corps du rapport}\\
    \textbf{1. Diagnostic :} le contrôle de l'installation a montré que le variateur de fréquence était hors service ; le moteur et la pompe étaient en bon état.\\
    \textbf{2. Approvisionnement :} un variateur neuf, de caractéristiques identiques, a été acquis auprès d'un fournisseur agréé (facture jointe).\\
    \textbf{3. Montage :} après consignation de l'installation (coupure et verrouillage de l'alimentation électrique), l'ancien variateur a été déposé et le nouveau monté conformément au schéma de câblage.\\
    \textbf{4. Paramétrage :} les paramètres ont été réglés selon la notice du constructeur (fréquence, rampes d'accélération et de décélération).\\
    \textbf{5. Essais :} la pompe a démarré normalement ; le débit et la pression mesurés sont conformes aux valeurs attendues.

    \textbf{Conclusion}\\
    L'intervention a permis de rétablir le fonctionnement nominal de la pompe de forage. Le respect de la procédure de consignation a garanti la sécurité de l'opérateur. Le matériel remplacé est traçable via la facture jointe.

    Signature : Kamga X., élève de 1re Électrotechnique.
\end{enumerate}
\textbf{Barème indicatif (sur 20) :} 10 erreurs correctement relevées, typées et corrigées (10 pts) ; introduction réécrite conforme (3 pts) ; corps réécrit en 5 étapes logiques, objectif et sécurisé (5 pts) ; langue et présentation (2 pts).

\textbf{Points de vigilance :} distinguer nettement les trois catégories d'erreurs (subjectivité, rupture de cohérence, erreur de structure) ; la correction doit toujours proposer une formulation professionnelle, pas seulement supprimer ; la sécurité (consignation avant intervention) est un impératif technique et rédactionnel ; vérifier la présence de tous les éléments de structure externe (en-tête, référence, objet, date, destinataire, expéditeur, signature, pièces jointes).
\end{corrige}

\newpage

\coursec{Fiche bilan}

\begin{popfigure}
\begin{tikzpicture}[mindmap, grow cyclic, every node/.style={concept, text=white, font=\small}, concept color=popblue, level 1/.append style={level distance=3.6cm, sibling angle=60}, level 2/.append style={level distance=2.2cm, sibling angle=45}]
\node[concept, font=\small\bfseries] {Le rapport}
  child[concept color=poporange] { node {Définition et types}
    child { node[concept color=poporange!80] {administratif, technique, de stage} }
    child { node[concept color=poporange!80] {objectivité, exactitude} }
  }
  child[concept color=popgreen] { node {Préparation}
    child { node[concept color=popgreen!80] {collecte des données} }
    child { node[concept color=popgreen!80] {plan, fiches} }
  }
  child[concept color=popblue] { node {Structure externe}
    child { node[concept color=popblue!80] {page de garde, sommaire} }
    child { node[concept color=popblue!80] {annexes, signature} }
  }
  child[concept color=poppurple] { node {Structure interne}
    child { node[concept color=poppurple!80] {introduction} }
    child { node[concept color=poppurple!80] {corps développé} }
    child { node[concept color=poppurple!80] {conclusion} }
  }
  child[concept color=poppink] { node {Texte explicatif}
    child { node[concept color=poppink!80] {définition, exemple, comparaison} }
    child { node[concept color=poppink!80] {connecteurs logiques} }
  }
  child[concept color=popgold] { node {Lexique}
    child { node[concept color=popgold!80] {lexique commun} }
    child { node[concept color=popgold!80] {lexique spécialisé} }
  };
\end{tikzpicture}
\legende{Carte mentale de la leçon : les six savoirs essentiels autour du rapport.}
\end{popfigure}

\begin{bilan}
\textbf{Définitions clés}
\begin{itemize}
\item Le \cle{rapport} : écrit utilitaire qui rend compte de façon \textbf{objective}, exacte et ordonnée d'une mission, d'une enquête, d'un stage ou d'une activité, à destination d'un destinataire précis (chef de service, jury, administration).
\item Le \cle{texte explicatif} : texte qui fait comprendre un phénomène, un fonctionnement ou une notion ; il procède par définition, exemple, comparaison, cause et conséquence.
\item Le \cle{lexique commun} : ensemble des mots compris et utilisés par tous les locuteurs (ex. \emph{machine, outil, réparer}).
\item Le \cle{lexique spécialisé} : ensemble des termes propres à un métier, une science ou une technique (ex. \emph{culasse, tension, cahier des charges}).
\item Le \cle{texte fonctionnel} : texte de la vie pratique (notice, affiche, consigne, formulaire) qui vise une action ou une information immédiate.
\end{itemize}

\textbf{À connaître par c\oe ur}
\begin{center}
\begin{tabularx}{\linewidth}{|l|X|}
\hline
\rowcolor{popblueL} \textbf{Élément} & \textbf{Contenu obligatoire} \\
\hline
Structure externe & page de garde (titre, auteur, date, destinataire), sommaire, pagination, annexes, signature \\
\hline
Introduction & contexte, objet du rapport, méthode, annonce du plan \\
\hline
Corps & faits constatés, données chiffrées, analyse, résultats ; paragraphes liés par des connecteurs \\
\hline
Conclusion & bilan, réponse à la mission, suggestions ou recommandations \\
\hline
Qualités du style & clarté, précision, objectivité (pas de « je » lyrique), concision \\
\hline
\end{tabularx}
\end{center}

\textbf{Méthodes express}
\begin{itemize}
\item \textbf{Préparer} : définir l'objet $\rightarrow$ collecter et trier les données $\rightarrow$ établir un plan $\rightarrow$ rédiger $\rightarrow$ relire.
\item \textbf{Expliquer} : définir le terme, donner un exemple, comparer, tirer la conséquence ; employer \emph{c'est-à-dire, en effet, par conséquent, ainsi}.
\item \textbf{Lire une affiche de sécurité} : identifier le pictogramme, la couleur (rouge = interdiction, jaune = danger, vert = issue/secours), la consigne et le public visé.
\item \textbf{Choisir le lexique} : adapter le vocabulaire au destinataire ; définir tout terme spécialisé à sa première occurrence.
\end{itemize}
\end{bilan}

\begin{aretenir}[Checklist avant le Bac]
\begin{itemize}
\item[$\square$] Je sais définir le rapport et citer au moins trois types de rapports.
\item[$\square$] Je connais les étapes de la préparation d'un rapport.
\item[$\square$] Je maîtrise la structure externe (page de garde, sommaire, annexes, signature).
\item[$\square$] Je sais rédiger une introduction : contexte, objet, méthode, plan.
\item[$\square$] Je sais construire le corps : faits, données, analyse, connecteurs logiques.
\item[$\square$] Je sais rédiger une conclusion avec bilan et recommandations.
\item[$\square$] Je reconnais les procédés du texte explicatif (définition, exemple, comparaison).
\item[$\square$] Je distingue lexique commun et lexique spécialisé et je les emploie à bon escient.
\item[$\square$] Je sais lire et interpréter une affiche d'indications sécuritaires.
\item[$\square$] Je rédige de façon objective, claire et concise, sans fautes d'orthographe.
\item[$\square$] Je relis mon rapport pour vérifier la cohérence et la présentation.
\end{itemize}
\end{aretenir}

\findecours

\end{document}
